% Le travail et les échanges — Philosophie 1re Tech — Cours signé OnBuch+
% Programme officiel MINESEC. Compiler avec : tectonic N13-travail-echanges.tex
\newcommand{\DOCMATIERE}{Philosophie}
\newcommand{\DOCNIVEAU}{1re Tech}
\newcommand{\DOCMODULE}{Philosophie – classe de Première technique}
\newcommand{\DOCLECON}{N13}
\newcommand{\DOCTITRE}{Le travail et les échanges}
% OnBuch+ — préambule « cours signés » ENRICHI (compilé avec tectonic / XeLaTeX)
% Système visuel « Pop » d'OnBuch+ : fond crème (jamais blanc), bordures encre
% épaisses, ombres « dures » décalées, titres Archivo Black en capitales.
% Version MATHÉMATIQUES : graphes (pgfplots/tikz), boîtes pédagogiques
% (définition, propriété, méthode, attention, le savais-tu, exercice résolu,
% exercice, TP, bilan), page de garde et signature OnBuch+ ancrée en bas.
%
% Macros attendues dans main.tex AVANT \input{preamble} :
%   \DOCMATIERE  \DOCNIVEAU  \DOCMODULE  \DOCLECON (numéro)  \DOCTITRE
\documentclass[11pt]{article}

\usepackage{fontspec}
\usepackage[french]{babel}
\usepackage[a4paper,margin=2cm,top=3.4cm,bottom=2.8cm,headheight=1.4cm,footskip=1.3cm]{geometry}
\usepackage{amsmath,amssymb,amsfonts}
\usepackage{mathtools} % \xmapsto, \coloneqq... souvent utilisés par les modèles
\usepackage{cancel} % simplifications barrées (bilans de forces)
% \abs{z} (module, valeur absolue) et \norm{u} (norme), utilisés par les modèles
\providecommand{\abs}{}\let\abs\relax\DeclarePairedDelimiter{\abs}{\lvert}{\rvert}
\providecommand{\norm}{}\let\norm\relax\DeclarePairedDelimiter{\norm}{\lVert}{\rVert}
\usepackage[table]{xcolor} % [table] : fournit aussi \rowcolors (alternance de lignes)
\usepackage{pagecolor}
\usepackage{tikz}
\usetikzlibrary{arrows.meta,positioning,calc,shapes.geometric,shapes.misc,shapes.arrows,decorations.pathmorphing,decorations.pathreplacing,decorations.markings,patterns,shadows,mindmap,trees,fit,backgrounds,matrix,intersections,angles,quotes}
\usepackage{pgfplots}
\usepackage[version=4]{mhchem} % équations nucléaires \ce{^{235}_{92}U}
\usepackage[european,siunitx]{circuitikz} % schémas électriques
\pgfplotsset{compat=1.18}
\usepgfplotslibrary{fillbetween,groupplots} % aires entre courbes (calcul intégral)
\usepackage{enumitem}
\usepackage{fancyhdr}
% [most] charge toutes les bibliothèques tcolorbox (listings, minted, raster,
% documentation...) et gonfle inutilement la mémoire TeX sur un document long
% (~300 boîtes) : on ne charge que ce qui est réellement utilisé.
\usepackage[skins,breakable]{tcolorbox}
\usepackage{graphicx}
\usepackage{colortbl}
\usepackage{booktabs}
\usepackage{tabularx}
\usepackage{multirow}
\usepackage{diagbox} % case d'en-tête diagonale des tableaux à double entrée
% Colonnes extensibles centrée / gauche / droite (C, L, R), souvent utilisées par les modèles
\newcolumntype{C}{>{\centering\arraybackslash}X}
\newcolumntype{L}{>{\raggedright\arraybackslash}X}
\newcolumntype{R}{>{\raggedleft\arraybackslash}X}
\usepackage{array}
\usepackage{multicol}
\usepackage{siunitx}
% parse-numbers=false : accepte aussi \SI{2,5 \times 10^{-3}}{...}
\sisetup{locale=FR,output-decimal-marker={,},group-minimum-digits=4,parse-numbers=false}
\DeclareSIUnit\ohm{\ensuremath{\symup{\Omega}}} % U+2126 absent de la police
\DeclareSIUnit\division{div} % réglages d'oscilloscope (ms/div, V/div)
\DeclareSIUnit\tr{tr} % tours (vitesse de rotation, ex. \SI{1500}{\tr\per\minute})
\DeclareSIUnit\molar{M} % concentration molaire (ex. \SI{3}{\molar}, \SI{1}{\micro\molar})
\DeclareSIUnit\an{an} % année (le modèle confond parfois avec \year, un compteur TeX) — ex. \SI{5}{\kilo\meter\per\an}
\DeclareSIUnit\FCFA{FCFA} % franc CFA (ex. \SI{500}{\FCFA\per\kilo\gram})
\DeclareSIUnit\degreeBrix{\textdegree Brix} % teneur en sucre d'un sirop/jus (réfractométrie)
\DeclareSIUnit\month{mois} % le modèle utilise parfois \month, un compteur TeX, comme unité de durée
\DeclareSIUnit\microliter{\micro\liter} % le modèle écrit parfois \microliter en un seul token
\DeclareSIUnit\million{millions} % dénombrement (ex. \SI{15}{\million\per\mL} spermatozoïdes)
\usepackage{qrcode}
% \degree : commande fréquemment utilisée par les modèles pour un angle ou une
% température en dehors de \SI{}{} (non fournie par les paquets chargés).
\providecommand{\degree}{^{\circ}}
% \qed (fin de démonstration), utilisé par les modèles sans amsthm
\providecommand{\qed}{\hfill$\square$}
% \barre : double barre des valeurs interdites dans un tableau de variations (tkz-tab)
\providecommand{\barre}{\ensuremath{\|}}
% environnement proof (amsthm non chargé) utilisé par les modèles
\ifdefined\proof\else\newenvironment{proof}[1][Démonstration]{\par\noindent\textit{#1.}\ }{\qed\par}\fi
\usepackage{lastpage}
\usepackage{hyperref}
\hypersetup{hidelinks}

% ============================================================
% Palette « Pop » OnBuch+ — mêmes hex que la classe P (plus_ui.dart)
% ============================================================
\definecolor{poporange}{HTML}{F59321}
\definecolor{popdark}{HTML}{E07A0C}
\definecolor{popcream}{HTML}{FFF6E8}
\definecolor{popink}{HTML}{1C1714}
\definecolor{poppurple}{HTML}{7A5AE0}
\definecolor{popgreen}{HTML}{2E9E6A}
\definecolor{poppink}{HTML}{E0457B}
\definecolor{popblue}{HTML}{2F7DE1}
\definecolor{popgold}{HTML}{C98A12}
\definecolor{popmuted}{HTML}{8A7F76}
\definecolor{popline}{HTML}{EADFD2}
\definecolor{popgray}{HTML}{8A7F76}
\colorlet{popgrey}{popgray}
% Alias tolérants (le modèle nomme parfois les couleurs différemment)
\colorlet{popwhite}{white}
\colorlet{popblack}{popink}
\colorlet{popred}{poppink}
\colorlet{popyellow}{popgold}
% Teintes claires (fonds de tableaux/encadrés) — le modèle nomme parfois
% les couleurs avec un suffixe L (ex. popblueL) sans les avoir définies.
\colorlet{poporangeL}{poporange!20}
\colorlet{popdarkL}{popdark!20}
\colorlet{poppurpleL}{poppurple!20}
\colorlet{popgreenL}{popgreen!20}
\colorlet{poppinkL}{poppink!20}
\colorlet{popblueL}{popblue!20}
\colorlet{popgoldL}{popgold!20}
\colorlet{popredL}{popred!20}
\colorlet{popgrayL}{popgray!20}
% Teintes claires pour fonds de tableaux / zones
\colorlet{popblueL}{popblue!10!white}
\colorlet{poporangeL}{poporange!14!white}
\colorlet{popgreenL}{popgreen!12!white}
\colorlet{poppurpleL}{poppurple!12!white}
\colorlet{poppinkL}{poppink!12!white}
\colorlet{popgoldL}{popgold!14!white}

\pagecolor{popcream}

% ============================================================
% Typographie
% ============================================================
\defaultfontfeatures{Ligatures=TeX}
\setmainfont{PlusJakartaSans-Medium.ttf}[Path=../../../fonts/,BoldFont=PlusJakartaSans-Bold.ttf]
% Maths en sans-serif (Fira Math) avec chiffres et lettres droites de Plus
% Jakarta, pour que les formules se fondent dans le texte.
\usepackage[math-style=ISO,bold-style=ISO]{unicode-math}
\setmathfont{FiraMath-Regular.otf}[Path=../../../fonts/]
\setmathfont{PlusJakartaSans-Medium.ttf}[Path=../../../fonts/,range={up/num,up/{Latin,latin},\mathup}]
\usepackage{icomma} % virgule décimale sans espace en mode math
% Caractères mathématiques Unicode absents des polices de texte (Plus Jakarta,
% Archivo Black) : rendus par leur équivalent mathématique, en texte comme en
% formule — sinon ils s'affichent en carré vide (ex. « Arithmétique dans ℤ »).
\usepackage{newunicodechar}
\newunicodechar{ℝ}{\ensuremath{\mathbb{R}}}
\newunicodechar{ℤ}{\ensuremath{\mathbb{Z}}}
\newunicodechar{ℂ}{\ensuremath{\mathbb{C}}}
\newunicodechar{ℕ}{\ensuremath{\mathbb{N}}}
\newunicodechar{ℚ}{\ensuremath{\mathbb{Q}}}
\newunicodechar{𝔻}{\ensuremath{\mathbb{D}}}
\newunicodechar{α}{\ensuremath{\alpha}}
\newunicodechar{β}{\ensuremath{\beta}}
\newunicodechar{γ}{\ensuremath{\gamma}}
\newunicodechar{δ}{\ensuremath{\delta}}
\newunicodechar{ε}{\ensuremath{\varepsilon}}
\newunicodechar{θ}{\ensuremath{\theta}}
\newunicodechar{λ}{\ensuremath{\lambda}}
\newunicodechar{μ}{\ensuremath{\mu}}
\newunicodechar{σ}{\ensuremath{\sigma}}
\newunicodechar{τ}{\ensuremath{\tau}}
\newunicodechar{φ}{\ensuremath{\varphi}}
\newunicodechar{ψ}{\ensuremath{\psi}}
\newunicodechar{ω}{\ensuremath{\omega}}
\newunicodechar{Σ}{\ensuremath{\Sigma}}
% majuscules produites par \MakeUppercase dans les titres (α-aminés -> Α)
\newunicodechar{Α}{\ensuremath{\alpha}}
\newunicodechar{Β}{\ensuremath{\beta}}
\newunicodechar{Γ}{\ensuremath{\Gamma}}
\newunicodechar{Δ}{\ensuremath{\Delta}}
\newunicodechar{Ε}{\ensuremath{\varepsilon}}
\newunicodechar{Θ}{\ensuremath{\Theta}}
\newunicodechar{Λ}{\ensuremath{\Lambda}}
\newunicodechar{Μ}{\ensuremath{\mu}}
\newunicodechar{Τ}{\ensuremath{\tau}}
\newunicodechar{Φ}{\ensuremath{\Phi}}
\newunicodechar{Ψ}{\ensuremath{\Psi}}
\newunicodechar{Ω}{\ensuremath{\Omega}}
\newunicodechar{↦}{\ensuremath{\mapsto}}
\newunicodechar{∈}{\ensuremath{\in}}
\newunicodechar{∉}{\ensuremath{\notin}}
\newunicodechar{∧}{\ensuremath{\wedge}}
\newunicodechar{⇔}{\ensuremath{\Leftrightarrow}}
\newunicodechar{⟺}{\ensuremath{\Longleftrightarrow}}
\newunicodechar{⇒}{\ensuremath{\Rightarrow}}
\newunicodechar{⟹}{\ensuremath{\Longrightarrow}}
\newunicodechar{∘}{\ensuremath{\circ}}
\newunicodechar{∪}{\ensuremath{\cup}}
\newunicodechar{∩}{\ensuremath{\cap}}
\newunicodechar{≡}{\ensuremath{\equiv}}
\newunicodechar{⊂}{\ensuremath{\subset}}
\newunicodechar{⊥}{\ensuremath{\perp}}
\newunicodechar{∃}{\ensuremath{\exists}}
\newunicodechar{∀}{\ensuremath{\forall}}
\newunicodechar{∛}{\ensuremath{\sqrt[3]{\,}}}
\newunicodechar{∜}{\ensuremath{\sqrt[4]{\,}}}
\newunicodechar{⃗}{\ensuremath{{}^{\to}}}
\newunicodechar{ⁿ}{\ifmmode^{n}\else\textsuperscript{\ensuremath{n}}\fi}
\newunicodechar{ˣ}{\ifmmode^{x}\else\textsuperscript{\ensuremath{x}}\fi}
\newunicodechar{ᵇ}{\ifmmode^{b}\else\textsuperscript{\ensuremath{b}}\fi}
\newunicodechar{ʳ}{\ifmmode^{r}\else\textsuperscript{\ensuremath{r}}\fi}
\newunicodechar{ᵘ}{\ifmmode^{u}\else\textsuperscript{\ensuremath{u}}\fi}
\newunicodechar{ʸ}{\ifmmode^{y}\else\textsuperscript{\ensuremath{y}}\fi}
\newunicodechar{ᵃ}{\ifmmode^{a}\else\textsuperscript{\ensuremath{a}}\fi}
\newunicodechar{ᵉ}{\ifmmode^{e}\else\textsuperscript{\ensuremath{e}}\fi}
\newunicodechar{ᵏ}{\ifmmode^{k}\else\textsuperscript{\ensuremath{k}}\fi}
\newunicodechar{ᵖ}{\ifmmode^{p}\else\textsuperscript{\ensuremath{p}}\fi}
\newunicodechar{ᴺ}{\ifmmode^{N}\else\textsuperscript{\ensuremath{N}}\fi}
\newunicodechar{ᶜ}{\ifmmode^{c}\else\textsuperscript{\ensuremath{c}}\fi}
\newunicodechar{ʰ}{\ifmmode^{h}\else\textsuperscript{\ensuremath{h}}\fi}
\newunicodechar{ᵠ}{\ifmmode^{q}\else\textsuperscript{\ensuremath{q}}\fi}
\newunicodechar{ₙ}{\ifmmode_{n}\else\textsubscript{\ensuremath{n}}\fi}
\newunicodechar{ₐ}{\ifmmode_{a}\else\textsubscript{\ensuremath{a}}\fi}
\newunicodechar{ᵢ}{\ifmmode_{i}\else\textsubscript{\ensuremath{i}}\fi}
\newunicodechar{ᵣ}{\ifmmode_{r}\else\textsubscript{\ensuremath{r}}\fi}
\newunicodechar{ₖ}{\ifmmode_{k}\else\textsubscript{\ensuremath{k}}\fi}
\newunicodechar{ᵦ}{\ifmmode_{\beta}\else\textsubscript{\ensuremath{\beta}}\fi}
\newunicodechar{ₓ}{\ifmmode_{x}\else\textsubscript{\ensuremath{x}}\fi}
\newfontfamily\popdisplay{ArchivoBlack-Regular.ttf}[Path=../../../fonts/]
\newfontfamily\poplogo{SpaceGrotesk-Bold.ttf}[Path=../../../fonts/]
\newfontfamily\popsemi{PlusJakartaSans-SemiBold.ttf}[Path=../../../fonts/]
\newfontfamily\popxbold{PlusJakartaSans-ExtraBold.ttf}[Path=../../../fonts/]
\newfontfamily\popxboldls{PlusJakartaSans-ExtraBold.ttf}[Path=../../../fonts/,LetterSpace=3.0]

\setlist[enumerate]{leftmargin=1.7em,itemsep=5pt,topsep=4pt}
\setlist[itemize]{leftmargin=1.7em,itemsep=4pt,topsep=4pt,label={\color{poporange}\textbullet}}
\setlength{\parindent}{0pt}
\setlength{\parskip}{5pt}
\renewcommand{\arraystretch}{1.25}

% pgfplots : style Pop par défaut
\pgfplotsset{
  every axis/.append style={axis line style={popink,line width=1pt},tick style={popink},
    grid style={popline,line width=0.5pt},label style={font=\small},tick label style={font=\footnotesize}},
  popcurve/.style={poporange,line width=1.8pt,no markers,smooth},
}
% Même style disponible dans un \draw TikZ simple (hors pgfplots)
\tikzset{popcurve/.style={poporange,line width=1.8pt}}
\tikzset{popgrid/.style={popline,very thin}}
\colorlet{popcurve}{poporange} % le nom du style est parfois utilisé comme couleur

% ============================================================
% Pill / squiggle
% ============================================================
\newcommand{\poppill}[3]{%
  \tikz[baseline=-0.55ex]\node[draw=popink,line width=1.2pt,rounded corners=2.6pt,%
    fill=#2,inner xsep=6.5pt,inner ysep=3pt]{%
    \popxboldls\fontsize{7}{7}\selectfont\color{#3}\MakeUppercase{#1}};%
}
\newcommand{\popsquiggle}[1]{%
  \tikz[baseline=-0.4ex]{
    \draw[#1,line width=2.6pt,line cap=round]
      plot [smooth] coordinates {(0,0.09) (0.34,0.01) (0.68,0.21) (1.02,0.01) (1.36,0.10)};
  }%
}
% Mot-clé surligné (marqueur orange sous le texte)
\newcommand{\cle}[1]{{\popxbold\color{popdark}#1}}

% ============================================================
% En-tête / pied
% ============================================================
\pagestyle{fancy}
\fancyhf{}
\renewcommand{\headrulewidth}{0pt}
\renewcommand{\footrulewidth}{0pt}
\lhead{\raisebox{-0.15ex}{\poplogo\fontsize{13}{13}\selectfont\color{popink}OnBuch\color{poporange}+}}
\rhead{\poppill{\DOCMATIERE\ \textbullet\ \DOCNIVEAU\ \textbullet\ Leçon \DOCLECON}{white}{popink}}
\lfoot{\popxbold\fontsize{7.6}{7.6}\selectfont\color{popmuted}\MakeUppercase{Cours signé OnBuch+}\ \textbullet\ {\popsemi\DOCTITRE}}
\rfoot{\tikz[baseline=-0.6ex]\node[rounded corners=8.5pt,fill=popink,minimum height=17pt,minimum width=17pt,inner xsep=5pt,inner sep=0]{\popxbold\fontsize{8}{8}\selectfont\color{popcream}\thepage\,/\,\pageref*{LastPage}};}

% ============================================================
% Titres
% ============================================================
\newcommand{\chaptitre}[2]{%
  \begin{tcolorbox}[enhanced,colback=poporange,colframe=popink,boxrule=2.6pt,arc=11pt,
      drop shadow=popink,left=15pt,right=15pt,top=12pt,bottom=11pt,before skip=3pt,after skip=18pt]
    \poppill{#2}{popink}{popcream}\par\vspace{9pt}
    {\raggedright\hyphenpenalty=10000\exhyphenpenalty=10000\popdisplay\fontsize{22}{24}\selectfont\color{popcream}\MakeUppercase{#1}\par}
  \end{tcolorbox}
}
% Section numérotée : pastille orange avec le numéro + titre Archivo
\newcounter{popsec}
\newcommand{\coursec}[1]{%
  \stepcounter{popsec}\setcounter{popsub}{0}%
  \par\vspace{16pt}\noindent
  \begin{minipage}{\linewidth}
  \tikz[baseline=-0.7ex]\node[draw=popink,line width=1.6pt,fill=poporange,rounded corners=4pt,minimum size=22pt,inner sep=0]{\popdisplay\fontsize{12}{12}\selectfont\color{popcream}\thepopsec};%
  \hspace{8pt}{\popdisplay\fontsize{15}{17}\selectfont\color{popink}\MakeUppercase{#1}}\par
  \vspace{4pt}\noindent{\color{poporange}\rule{36pt}{3.6pt}}
  \end{minipage}\par\nopagebreak\vspace{8pt}
}
\newcounter{popsub}
\newcommand{\courssub}[1]{%
  \stepcounter{popsub}%
  \par\vspace{9pt}\noindent{\popxbold\fontsize{11.5}{13}\selectfont\color{popink}{\color{poporange}\thepopsec.\thepopsub}\hspace{6pt}#1}\par\nopagebreak\vspace{4pt}
}

% ============================================================
% Boîtes pédagogiques — même recette : fond blanc, bordure épaisse colorée,
% ombre dure encre, badge plein. Titre optionnel [#1].
% ============================================================
\tcbset{popbase/.style={breakable,enhanced,colback=white,boxrule=2.3pt,arc=9pt,
  shadow={3pt}{-3pt}{0pt}{popink},left=14pt,right=14pt,top=11pt,bottom=11pt,before skip=11pt,after skip=15pt,
  fontupper=\popsemi\fontsize{10.6}{14.5}\selectfont\color{popink},
  fontlower=\popsemi\fontsize{10.6}{14.5}\selectfont\color{popink}}}
% Coche dessinée (le glyphe ✓ n'existe pas dans les polices du document)
\newcommand{\popcheck}[1]{\tikz[baseline=-0.5ex]\draw[#1,line width=1.6pt,line cap=round,line join=round](-0.13,0.0)--(-0.04,-0.09)--(0.13,0.1);}
\providecommand{\coche}{\popcheck{popgreen}}
\providecommand{\resultat}[1]{\par\smallskip\noindent\fcolorbox{popgreen}{popgreenL}{\parbox{\dimexpr\linewidth-2\fboxsep-2\fboxrule\relax}{\textbf{Résultat.} #1}}\par\smallskip}
\newcommand{\popboxhead}[3]{\poppill{#1}{#2}{white}\ifx\relax#3\relax\else\hspace{6pt}{\popxbold\fontsize{10.6}{12}\selectfont\color{popink}#3}\fi\par\vspace{7pt}}

\newtcolorbox{aretenir}[1][]{popbase,colframe=popgreen,before upper={\popboxhead{À retenir}{popgreen}{#1}}}
\newtcolorbox{exemplebox}[1][]{popbase,colframe=poporange,before upper={\popboxhead{Exemple}{poporange}{#1}}}
\newtcolorbox{definition}[1][]{popbase,colframe=popblue,before upper={\popboxhead{Définition}{popblue}{#1}}}
\newtcolorbox{propriete}[1][]{popbase,colframe=poppurple,before upper={\popboxhead{Propriété}{poppurple}{#1}}}
\newtcolorbox{methode}[1][]{popbase,colframe=popgold,before upper={\popboxhead{Méthode}{popgold}{#1}}}
\newtcolorbox{attention}[1][]{popbase,colframe=poppink,before upper={\popboxhead{Attention}{poppink}{#1}}}
\newtcolorbox{experience}[1][]{popbase,colframe=popblue,colback=popblueL,before upper={\popboxhead{Expérience}{popblue}{#1}}}
% « Le savais-tu ? » : carte violette pleine (contexte, culture, Cameroun)
\newtcolorbox{savaistu}[1][]{popbase,colback=poppurple,colframe=popink,
  fontupper=\popsemi\fontsize{10.4}{14.2}\selectfont\color{white},
  before upper={\poppill{Le savais-tu\,?}{white}{poppurple}\ifx\relax#1\relax\else\hspace{6pt}{\popxbold\color{white}#1}\fi\par\vspace{7pt}}}
% Exercice résolu : énoncé en haut, \tcblower puis la solution sur fond crème
\newcounter{popexo}\newcounter{popexr}
\newtcolorbox{exoresolu}[1][]{popbase,colframe=poporange,colbacklower=popcream,
  segmentation style={popink,line width=1.2pt,dashed},
  before upper={\stepcounter{popexr}\popboxhead{Exercice résolu \thepopexr}{poporange}{#1}},
  before lower={\poppill{Solution}{popink}{popcream}\par\vspace{6pt}}}
% Exercice (non résolu en place) — la correction est dans la partie Corrigés
\newtcolorbox{exercice}[1][]{popbase,colframe=popink,
  before upper={\stepcounter{popexo}\popboxhead{Exercice \thepopexo}{popink}{#1}}}
% Corrigé
\newtcolorbox{corrige}[1][]{popbase,colframe=popgreen,colback=popgreenL,
  before upper={\popboxhead{Corrigé}{popgreen}{#1}}}
% Objectifs / prérequis / bilan
\newtcolorbox{objectifs}{popbase,colframe=popink,colback=poporangeL,
  before upper={\popboxhead{Objectifs de la leçon}{popink}{}}}
\newtcolorbox{prerequis}{popbase,colframe=popink,colback=popblueL,
  before upper={\popboxhead{Prérequis}{popblue}{}}}
\newtcolorbox{bilan}{popbase,colframe=popink,colback=white,boxrule=2.8pt,
  before upper={\popboxhead{Fiche bilan}{poporange}{L'essentiel à connaître pour l'examen}}}
% Figure encadrée avec légende
\newtcolorbox{popfigure}[1][]{enhanced,breakable=false,colback=white,colframe=popline,boxrule=1.4pt,arc=8pt,
  left=10pt,right=10pt,top=10pt,bottom=8pt,before skip=10pt,after skip=12pt,halign=center,
  lower separated=false,halign lower=center,
  fontlower=\popsemi\fontsize{9}{11.5}\selectfont\color{popmuted},
  #1}
\newcommand{\legende}[1]{\par\vspace{6pt}{\popsemi\fontsize{9}{11.5}\selectfont\color{popmuted}#1}}

% ============================================================
% Signature OnBuch+
% ============================================================
\newcommand{\poplogoinline}[1]{{\poplogo\fontsize{#1}{#1}\selectfont\color{popink}OnBuch\color{poporange}+}}
\newcommand{\signatureonbuch}{%
  \begin{tcolorbox}[enhanced,colback=popink,colframe=popink,boxrule=2.2pt,arc=10pt,
      drop shadow=poporange,left=14pt,right=14pt,top=10pt,bottom=10pt,before skip=0pt,after skip=0pt]
    \tikz[baseline=-0.7ex]\node[circle,fill=popgreen,draw=popcream,line width=1.3pt,minimum size=22pt,inner sep=0]{\popcheck{white}};%
    \hspace{9pt}\begin{minipage}[c]{0.62\linewidth}
      {\popxbold\fontsize{12}{13}\selectfont\color{popcream}Signé\ \textbullet\ {\poplogo OnBuch\color{poporange}+}}\par\vspace{2pt}
      {\raggedright\popsemi\fontsize{8}{10.5}\selectfont\color{popline}\MakeUppercase{Équipe pédagogique OnBuch+}\\\MakeUppercase{Conforme au programme officiel MINESEC}\par}
    \end{minipage}\hfill
    {\poplogo\fontsize{22}{22}\selectfont\color{popcream}OnBuch\color{poporange}+}
  \end{tcolorbox}%
}

% ============================================================
% Page de garde
% #1 titre · #2 matière · #3 niveau · #4 module · #5 n° leçon · #6 durée ·
% #7 description / objectifs (texte ou itemize)
% ============================================================
\newcommand{\pagegarde}[7]{%
  \thispagestyle{empty}
  \newgeometry{a4paper,margin=1.7cm,top=1.6cm,bottom=1.4cm}%
  \begin{tikzpicture}[remember picture,overlay]
    \fill[poporange!18] ([xshift=-3.2cm,yshift=-3.6cm]current page.north east) circle (4.6cm);
    \fill[poppurple!14] ([xshift=2.4cm,yshift=7.6cm]current page.south west) circle (3.8cm);
  \end{tikzpicture}%
  {\poplogo\fontsize{30}{30}\selectfont\color{popink}OnBuch\color{poporange}+}\hfill
  \raisebox{0.6ex}{\poppill{Cours signé \textbullet\ Premium}{popink}{popcream}}\par
  \vspace{1pt}\popsquiggle{poppurple}\par
  \vspace{8pt}
  \begin{tcolorbox}[enhanced,colback=poporange,colframe=popink,boxrule=2.8pt,arc=13pt,
      drop shadow=popink,left=18pt,right=18pt,top=12pt,bottom=13pt,after skip=10pt]
    \poppill{#2 \textbullet\ #3}{popink}{popcream}\hspace{5pt}\poppill{#4}{white}{popink}\par\vspace{8pt}
    {\popdisplay\fontsize{14}{14}\selectfont\color{popink}LEÇON #5}\par\vspace{4pt}
    {\raggedright\hyphenpenalty=10000\exhyphenpenalty=10000\popdisplay\fontsize{25}{27}\selectfont\color{popcream}\MakeUppercase{#1}\par}
    \vspace{8pt}
    {\popxbold\fontsize{9.5}{9.5}\selectfont\color{popink}\MakeUppercase{Durée conseillée : #6}}
  \end{tcolorbox}
  \begin{tcolorbox}[enhanced,colback=white,colframe=popink,boxrule=2.2pt,arc=10pt,
      drop shadow=popink,left=16pt,right=16pt,top=10pt,bottom=10pt,after skip=8pt]
    \poppill{Dans cette leçon}{poppurple}{white}\par\vspace{5pt}
    \popsemi\fontsize{9.3}{12.6}\selectfont\color{popink}#7
  \end{tcolorbox}
  \noindent\begin{minipage}[t]{0.487\linewidth}
    \begin{tcolorbox}[enhanced,colback=popgreen,colframe=popink,boxrule=2.2pt,arc=10pt,
        drop shadow=popink,left=11pt,right=11pt,top=10pt,bottom=10pt,height=3.1cm,valign=center]
      \tcbox[colback=white,colframe=popink,boxrule=1.3pt,arc=5pt,boxsep=0pt,nobeforeafter,
        left=3pt,right=3pt,top=3pt,bottom=3pt,valign=center]{\qrcode[height=1.9cm]{https://play.google.com/store/apps/details?id=com.luvvix.onbuch&hl=fr}}%
      \hspace{8pt}\begin{minipage}[c]{0.5\linewidth}
        {\popdisplay\fontsize{11}{12.5}\selectfont\color{white}\MakeUppercase{Télécharge OnBuch}}\par\vspace{4pt}
        {\raggedright\popsemi\fontsize{8.4}{11}\selectfont\color{white}Gratuit \textbullet\ Google Play\\Tuteur IA, annales, concours, résultats\par}
      \end{minipage}
    \end{tcolorbox}
  \end{minipage}\hfill
  \begin{minipage}[t]{0.487\linewidth}
    \begin{tcolorbox}[enhanced,colback=poppurple,colframe=popink,boxrule=2.2pt,arc=10pt,
        drop shadow=popink,left=11pt,right=11pt,top=10pt,bottom=10pt,height=3.1cm,valign=center]
      \tikz[baseline=-0.7ex]\node[circle,fill=white,draw=popink,line width=1.3pt,minimum size=1.6cm,inner sep=0]{\popdisplay\fontsize{24}{24}\selectfont\color{poppurple}+};%
      \hspace{8pt}\begin{minipage}[c]{0.6\linewidth}
        {\popdisplay\fontsize{11}{12.5}\selectfont\color{white}\MakeUppercase{Passe à OnBuch+}}\par\vspace{4pt}
        {\raggedright\popsemi\fontsize{8.4}{11}\selectfont\color{white}Tous les cours signés, Tuteur IA illimité, fiches PDF — onglet OnBuch+ de l'app\par}
      \end{minipage}
    \end{tcolorbox}
  \end{minipage}
  \vspace{8pt}
  \popcontenu
  \vfill
  \signatureonbuch
  \restoregeometry
  \newpage
}

% Rangée « Ce cours contient » (page de garde)
\newcommand{\poptile}[3]{% #1 couleur #2 gros texte #3 légende
  \begin{tcolorbox}[enhanced,colback=#1,colframe=popink,boxrule=2pt,arc=9pt,shadow={2.5pt}{-2.5pt}{0pt}{popink},
      left=6pt,right=6pt,top=9pt,bottom=9pt,halign=center,valign=center,height=1.75cm,width=\linewidth,nobeforeafter]
    {\popdisplay\fontsize{12.5}{13}\selectfont\color{white}\MakeUppercase{#2}}\par\vspace{3pt}
    {\centering\popsemi\fontsize{7.8}{9.5}\selectfont\color{white}#3\par}
  \end{tcolorbox}}
\newcommand{\popcontenu}{%
  \par\noindent{\popxboldls\fontsize{7.5}{7.5}\selectfont\color{popmuted}CE COURS SIGNÉ CONTIENT}\par\vspace{7pt}
  \noindent\begin{minipage}[t]{0.235\linewidth}\poptile{popblue}{Cours}{complet, illustré, pas à pas}\end{minipage}\hfill
  \begin{minipage}[t]{0.235\linewidth}\poptile{poporange}{Méthodes}{et exercices résolus}\end{minipage}\hfill
  \begin{minipage}[t]{0.235\linewidth}\poptile{poppink}{Exercices}{type examen + corrigés}\end{minipage}\hfill
  \begin{minipage}[t]{0.235\linewidth}\poptile{popgreen}{Fiche bilan}{l'essentiel à retenir}\end{minipage}\par}

% Sommaire
\newtcolorbox{sommaire}{enhanced,colback=white,colframe=popink,boxrule=2.2pt,arc=10pt,
  drop shadow=popink,left=18pt,right=18pt,top=13pt,bottom=13pt,before skip=6pt,after skip=20pt,
  before upper={\poppill{Sommaire}{poppurple}{white}\par\vspace{9pt}\popsemi\fontsize{10.6}{14.5}\selectfont\color{popink}}}

% Signature de fin de cours — poussée tout en bas de la dernière page
\newcommand{\findecours}{%
  \par\vfill\nopagebreak
  \begin{center}\popsquiggle{poporange}\end{center}\vspace{4pt}
  \signatureonbuch
}
\AtBeginDocument{\def\llbracket{\mathopen{[\mkern-3.5mu[}}\def\rrbracket{\mathclose{]\mkern-3.5mu]}}} % intervalles d'entiers (glyphe ⟦ absent de la police)
% Glyphes absents de Fira Math : redéfinis à partir de glyphes disponibles
\AtBeginDocument{%
  \def\setminus{\mathbin{\backslash}}\let\smallsetminus\setminus
  \def\vdots{\mathord{\vbox{\baselineskip3.2pt\lineskiplimit0pt\kern4pt\hbox{.}\hbox{.}\hbox{.}}}}%
  \def\ddots{\mathinner{\mkern1mu\raise6.5pt\hbox{.}\mkern2mu\raise3.6pt\hbox{.}\mkern2mu\raise0.7pt\hbox{.}\mkern1mu}}%
  \def\triangle{\mathord{\scalebox{0.9}{$\symup{\Delta}$}}}%
  \def\nabla{\mathord{\raisebox{\depth}{\scalebox{1}[-1]{$\symup{\Delta}$}}}}%
  \def\checkmark{\popcheck{popdark}}%
  \def\star{\mathbin{\ast}}\def\bigstar{\mathord{\ast}}%
  \def\diamond{\mathbin{\scalebox{0.6}{$\Diamond$}}}%
  \def\bigcup{\mathop{\mathchoice{\raisebox{-0.3em}{\scalebox{1.7}{$\cup$}}}{\scalebox{1.25}{$\cup$}}{\cup}{\cup}}\displaylimits}%
  \def\bigcap{\mathop{\mathchoice{\raisebox{-0.3em}{\scalebox{1.7}{$\cap$}}}{\scalebox{1.25}{$\cap$}}{\cap}{\cap}}\displaylimits}%
  \def\lesseqqgtr{\mathrel{\lessgtr}}\def\bigoplus{\mathop{\oplus}\displaylimits}%
}
\providecommand{\ssi}{si et seulement si} % abréviation fréquente
\AtBeginDocument{\providecommand{\wideparen}{\overparen}} % arc de cercle

\begin{document}

\pagegarde{Le travail et les échanges}{Philosophie}{1re Tech}{Philosophie – classe de Première technique}{N13}{2 séances (fiche de progression harmonisée)}{Cette leçon interroge deux réalités fondamentales de l'existence humaine : le travail, par lequel l'homme transforme la nature et se réalise, et les échanges, qui fondent la vie en société. À partir de situations concrètes camerounaises et de textes philosophiques, l'élève apprend à définir, analyser et problématiser ces notions pour construire une réflexion argumentée.
\vspace{4pt}
\begin{multicols}{2}\small
\begin{itemize}
\item À la fin de cette leçon, je sais définir le travail et le distinguer de l'activité animale et du simple loisir
\item À la fin de cette leçon, je sais expliquer les différentes valeurs du travail (valeur économique, sociale, morale et existentielle)
\item À la fin de cette leçon, je sais analyser le travail comme malédiction, comme libération ou comme aliénation à l'aide d'auteurs étudiés
\item À la fin de cette leçon, je sais définir les échanges et identifier leurs formes (économiques, linguistiques, symboliques)
\item À la fin de cette leçon, je sais montrer en quoi les échanges sont le fondement de la société
\item À la fin de cette leçon, je sais appliquer les notions de travail et d'échange à des situations concrètes de la vie courante au Cameroun
\end{itemize}\end{multicols}}

\begin{sommaire}
\begin{multicols}{2}
\begin{enumerate}
\item Qu'est-ce que le travail ?
\item Les valeurs du travail
\item Travail : malédiction, libération ou aliénation ?
\item Qu'est-ce que l'échange ?
\item Les échanges, fondement de la société
\item Méthode : appliquer les notions et rédiger une argumentation
\item Activité d'intégration
\item Méthodes et exercices résolus
\item Exercices
\item Corrigés des exercices
\item Fiche bilan
\end{enumerate}
\end{multicols}
\end{sommaire}

\begin{objectifs}
\begin{itemize}
\item À la fin de cette leçon, je sais définir le travail et le distinguer de l'activité animale et du simple loisir
\item À la fin de cette leçon, je sais expliquer les différentes valeurs du travail (valeur économique, sociale, morale et existentielle)
\item À la fin de cette leçon, je sais analyser le travail comme malédiction, comme libération ou comme aliénation à l'aide d'auteurs étudiés
\item À la fin de cette leçon, je sais définir les échanges et identifier leurs formes (économiques, linguistiques, symboliques)
\item À la fin de cette leçon, je sais montrer en quoi les échanges sont le fondement de la société
\item À la fin de cette leçon, je sais appliquer les notions de travail et d'échange à des situations concrètes de la vie courante au Cameroun
\item À la fin de cette leçon, je sais rédiger et justifier une démarche argumentative (introduction problématisée, thèse, antithèse, synthèse)
\item À la fin de cette leçon, je sais analyser un texte philosophique et en dégager la thèse et les arguments
\end{itemize}
\end{objectifs}

\begin{prerequis}
\begin{itemize}
\item La notion de culture et de nature (leçons précédentes de Première)
\item La distinction entre l'homme et l'animal (Seconde)
\item La notion de liberté et de conscience (Seconde)
\item La méthode de l'explication de texte et de la dissertation (acquis de Seconde)
\end{itemize}
\end{prerequis}

\newpage

\coursec{Qu'est-ce que le travail ?}

Chaque matin à Douala, des milliers de jeunes vendeurs à la sauvette, mototaximen et couturiers gagnent leur vie par le travail, tandis qu'au marché Mokolo de Yaoundé, vendeuses et clients échangent marchandises, paroles et services. Un élève de Première technique se demande : pourquoi travailler si l'on peut vivre autrement ? Le travail est-il seulement une contrainte pour survivre, ou est-il ce qui fait de nous des hommes ? Et pourquoi aucun homme ne peut-il vivre sans échanger avec les autres ? Ces questions, au cœur de la vie quotidienne camerounaise, sont aussi de grandes questions philosophiques.

\courssub{De l'étymologie à la définition conceptuelle}

\begin{savaistu}[L'origine douloureuse du mot]
Le mot français « travail » ne vient pas d'un terme noble. Il dérive du latin \emph{tripalium} (de \emph{tripalis} : « à trois pieux »), qui désignait un instrument de torture constitué de trois pieux sur lequel on attachait les esclaves pour les battre. Le verbe \emph{tripaliare} signifiait « torturer sur le tripalium », puis par extension « fatiguer, tourmenter ». En ancien français, \emph{travailler} a gardé le sens de « souffrir, peiner, tourmenter » (on disait « travailler de la fièvre » pour « avoir de la fièvre »). Ce n'est qu'au XIIe siècle que le sens glisse vers « peine prise pour faire quelque chose » puis « activité pénible » et enfin « occupation professionnelle ». L'étymologie nous livre donc une première vérité : le travail est d'abord perçu comme une \textbf{peine}, une \textbf{contrainte} subie, une souffrance physique.
\end{savaistu}

\begin{methode}[Définir une notion philosophique : la démarche en trois temps]
Pour construire une définition rigoureuse d'une notion comme le travail, on suit trois étapes :
\begin{enumerate}
    \item \textbf{L'analyse étymologique :} retrouver le sens originel du mot pour en saisir la charge historique (ici : la torture, la peine).
    \item \textbf{L'analyse du sens commun :} observer comment le mot s'emploie dans le langage ordinaire (ex. : « aller au travail », « c'est du travail », « travail manuel / intellectuel ») pour en dégager les invariants (effort, finalité, production).
    \item \textbf{La conceptualisation philosophique :} dépasser le sens commun pour proposer une définition structurée, opératoire, qui distingue l'essence de la notion de ses accidents (rémunération, pénibilité, statut social).
\end{enumerate}
\end{methode}

Au terme de cette démarche, nous pouvons poser la \textbf{définition conceptuelle} :

\begin{definition}[Le travail]
Le travail est une \textbf{activité humaine consciente, volontaire et finalisée}, par laquelle l'homme transforme la nature à l'aide d'outils et de techniques, pour produire des biens ou des services utiles, satisfaire ses besoins et se réaliser comme sujet.
\end{definition}

Cette définition repose sur quatre piliers que nous allons examiner : la \textbf{conscience} (projet), la \textbf{transformation de la nature} (matière), la \textbf{médiation technique} (outil), la \textbf{finalité} (besoin, utilité, réalisation de soi).

\courssub{Travail humain vs activité animale : le projet et l'outil}

Une erreur fréquente consiste à dire que « les animaux travaillent aussi » (la fourmi, l'abeille, le castor). Marx, dans \emph{Le Capital} (Livre I, Ch. 7), propose une distinction décisive à partir de l'exemple de l'architecte et de l'abeille.

\begin{exemplebox}[L'architecte et l'abeille (Marx)]
« L'abeille met à honte plus d'un architecte par la construction de ses cellules en cire. Mais ce qui distingue dès le début le plus mauvais architecte de la meilleure des abeilles, c'est qu'il a construit la cellule dans sa tête avant de la construire en cire. »
\end{exemplebox}

\begin{popfigure}
\begin{tikzpicture}[
    node distance=1.2cm and 2.5cm,
    box/.style={draw=popdark, thick, rounded corners, fill=popcream, minimum width=3.5cm, minimum height=1.2cm, align=center, font=\small},
    header/.style={draw=popblue, thick, rounded corners, fill=popblueL, minimum width=3.5cm, minimum height=1cm, align=center, font=\bfseries\small},
    arrow/.style={-Stealth, thick, popdark}
]
% Title
\node[header] (title) {Comparaison : Activité Animale vs Travail Humain};
% Column headers
\node[header, below=1.5cm of title, xshift=-3.5cm] (animal) {Activité Animale (Instinct)};
\node[header, below=1.5cm of title, xshift=3.5cm] (human) {Travail Humain (Projet)};
% Rows
% 1. Origine
\node[box, below=0.2cm of animal] (a1) {Déclenchée par un \textbf{stimulus extérieur} (saison, faim, danger)};
\node[box, below=0.2cm of human] (h1) {Née d'une \textbf{représentation intérieure} (idée, plan, désir)};
% 2. Outil
\node[box, below=0.1cm of a1] (a2) {Utilise son \textbf{corps} comme outil (mandibules, pattes, bec) \emph{ou} objets trouvés non modifiés};
\node[box, below=0.1cm of h1] (h2) {Fabrique et utilise des \textbf{outils médiateurs} (marteau, machine, logiciel) qui prolongent la main};
% 3. Conscience / Transmission
\node[box, below=0.1cm of a2] (a3) {Transmission \textbf{génétique} (innée) ; peu de variation individuelle};
\node[box, below=0.1cm of h2] (h3) {Transmission \textbf{culturelle} (apprise, enseignée) ; perfectible, cumulative};
% 4. Finalité
\node[box, below=0.1cm of a3] (a4) {Finalité \textbf{immédiate} : survie de l'espèce (nourriture, abri, reproduction)};
\node[box, below=0.1cm of h3] (h4) {Finalité \textbf{médiate} : besoins présents + projets futurs, création de sens, beauté, savoir};
% 5. Résultat
\node[box, below=0.1cm of a4] (a5) {Résultat \textbf{stéréotypé}, identique pour tous les individus de l'espèce};
\node[box, below=0.1cm of h4] (h5) {Résultat \textbf{singulier}, porteur de la marque de son auteur (style, innovation)};
% Braces
\draw[popdark, thick, decorate, decoration={brace, amplitude=5pt}] (a1.north west) -- (a5.south west) node[midway, left=8pt, rotate=90, font=\bfseries\small] {Hétérodéterminé};
\draw[popdark, thick, decorate, decoration={brace, amplitude=5pt}] (h1.north east) -- (h5.south east) node[midway, right=8pt, rotate=-90, font=\bfseries\small] {Autodéterminé (Liberté)};
\end{tikzpicture}
\legende{Tableau comparatif : l'activité animale est close sur l'instinct et la répétition ; le travail humain s'ouvre sur le projet, la technique et l'histoire.}
\end{popfigure}

\begin{aretenir}[La rupture anthropologique]
L'animal est \emph{dans} la nature ; l'homme, par le travail, se \emph{détache} de la nature pour s'y opposer et la transformer. Le travail est l'acte fondateur de l'humanité : « Le travail a créé l'homme lui-même » (Engels, \emph{Le rôle du travail dans la transformation du singe en homme}). C'est parce qu'il travaille que l'homme développe son intelligence, son langage et sa société.
\end{aretenir}

\courssub{Travail, jeu et loisir : la contrainte et l'utilité}

Tout n'est pas travail dans l'activité humaine. Il faut distinguer le travail du \textbf{jeu} et du \textbf{loisir}.

\begin{propriete}[Critères de distinction]
\begin{itemize}
    \item \textbf{Le travail :} Contrainte (nécessité vitale ou sociale), finalité \textbf{extérieure} à l'activité (le produit, le salaire, l'utilité), effort soutenu, hétéronomie partielle (horaires, normes, commande).
    \item \textbf{Le jeu :} Activité \textbf{gratuite} (fin en soi), libre (on peut s'arrêter), régie par des règles acceptées volontairement, source de plaisir immédiat (Caillois, \emph{Les Jeux et les Hommes}).
    \item \textbf{Le loisir :} Temps \textbf{libéré} des contraintes professionnelles, consacré au repos, à la culture, à la vie associative. Ce n'est pas de l'oisiveté : c'est le temps de la \emph{scholè} (école, culture) chez les Grecs.
\end{itemize}
\end{propriete}

\begin{attention}[Piège : confondre occupation et travail]
\emph{Erreur fréquente :} « Je travaille sur mon jeu vidéo » ou « Le jardinage du dimanche, c'est du travail. »
\emph{Correction :} Tout travail est une activité, mais toute activité n'est pas un travail.
\begin{itemize}
    \item Le jardinage du retraité passionné est un \textbf{loisir} (liberté, fin en soi, pas de contrainte de rendement).
    \item Le jardinage du maraîcher de Mfoundi est du \textbf{travail} (contrainte de rendement, de calendrier, de vente pour vivre).
    \item Le « work » dans les jeux vidéo (farming, grinding) imite la structure du travail (effort, récompense), mais reste du \textbf{jeu} (sortie possible, finalité ludique).
\end{itemize}
La frontière est poreuse (l'artiste, l'artisan d'art), mais le critère décisif reste : \textbf{l'activité est-elle subie ou choisie ? Le produit est-il visé pour lui-même ou pour un usage/échange ?}
\end{attention}

\courssub{Les quatre caractéristiques structurelles du travail humain}

Au-delà de la définition, le travail possède une structure invariante que l'on retrouve chez le cultivateur de l'Ouest comme chez l'informaticien de Yaoundé.

\begin{popfigure}
\begin{tikzpicture}[
    center/.style={circle, draw=popblue, thick, fill=popblueL, minimum size=2.2cm, font=\bfseries, text width=2cm, align=center},
    arm/.style={rectangle, draw=popdark, thick, rounded corners, fill=popcream, minimum width=3.2cm, minimum height=1.4cm, font=\small, align=center},
    line/.style={-Stealth, thick, popblue}
]
\node[center, align=center] (center){LE\\TRAVAIL\\HUMAIN};
% Four arms
\node[arm, above left=1.5cm and 1.5cm of center, align=center] (f1){\textbf{FINALITÉ}\\\emph{« Pour quoi ? »}\\\textbullet\ Besoins vitaux\\\textbullet\ Désirs culturels\\\textbullet\ Projet de vie};
\node[arm, above right=1.5cm and 1.5cm of center, align=center] (f2){\textbf{EFFORT / PEINE}\\\emph{« À quel prix ? »}\\\textbullet\ Dépense physique\\\textbullet\ Charge mentale\\\textbullet\ Temps contraint};
\node[arm, below right=1.5cm and 1.5cm of center, align=center] (f3){\textbf{TECHNIQUE / OUTIL}\\\emph{« Comment ? »}\\\textbullet\ Médiation instrumentale\\\textbullet\ Savoir-faire transmis\\\textbullet\ Innovation};
\node[arm, below left=1.5cm and 1.5cm of center, align=center] (f4){\textbf{SOCIALISATION}\\\emph{« Avec / Pour qui ? »}\\\textbullet\ Division du travail\\\textbullet\ Échange / Don\\\textbullet\ Reconn. sociale};
\draw[line] (center.north west) -- (f1.south east);
\draw[line] (center.north east) -- (f2.south west);
\draw[line] (center.south east) -- (f3.north west);
\draw[line] (center.south west) -- (f4.north east);
% Anchors for labels on lines (optional, kept simple)
\end{tikzpicture}
\legende{Les quatre caractéristiques structurelles du travail : elles sont indissociables. Supprimer l'une (ex. finalité), c'est changer la nature de l'activité.}
\end{popfigure}

\begin{exemplebox}[Quatre visages du travail au Cameroun]
\begin{itemize}
    \item \textbf{Le cultivateur de l'Ouest (Bamileke) :} Il transforme la pente en terrasses (technique), porte des charges lourdes (effort), vise la récolte de maïs/igname pour nourrir sa famille et vendre (finalité), participe aux tontines et aux marchés (socialisation).
    \item \textbf{Le pêcheur de Limbé :} Il connaît les courants et les saisons (savoir-faire), utilise pirogue et filets (outils), brave la mer la nuit (effort/risque), vend sa prise aux « bayam-sellam » (échange social), assure la protéine de la ville.
    \item \textbf{L'artisan forgeron de Foumban :} Il domine le feu et le métal (technique haute), crée des objets d'art et d'usage (finalité esthétique + utilitaire), transmet son savoir par apprentissage (socialisation), son travail fonde son statut et son identité.
    \item \textbf{L'informaticien de Yaoundé (Silicon Mountain) :} Son outil est le code/clavier (technique immatérielle), son effort est cognitif/attentionnel (charge mentale), sa finalité est la création de service/valeur (startup, e-gov), il travaille en équipe agile, pour des clients globaux (socialisation déterritorialisée).
\end{itemize}
Dans les quatre cas : \textbf{projet conscient + transformation médiatisée + effort + lien social}.
\end{exemplebox}

\courssub{Exercice résolu : Appliquer la définition à une situation concrète}

\begin{exoresolu}[Analyser une situation : « Le tailleur de pierre de la cathédrale »]
\textbf{Énoncé :} Un tailleur de pierre travaille sur le chantier de la cathédrale Notre-Dame-des-Victoires de Yaoundé. Il taille un bloc de granit selon les plans de l'architecte. Expliquez en quoi cette activité illustre parfaitement la définition philosophique du travail.

\textbf{Solution détaillée :}
\begin{enumerate}
    \item \textbf{Conscience et Projet (Finalité) :} Le tailleur ne frappe pas la pierre au hasard. Il a \emph{intériorisé} le plan de l'architecte : il « a construit la pierre dans sa tête avant de la construire en granit » (Marx). Son geste est guidé par une \textbf{représentation anticipée} du résultat final (la voûte, le pilier). C'est la marque de l'humanité : l'idée précède l'acte.
    \item \textbf{Transformation de la nature (Matière) :} Il prend un bloc brut (nature donnée) et lui impose une forme géométrique précise (nature transformée). Il \emph{humanise} la matière : la pierre devient « pierre de taille », porteuse de sens architectural.
    \item \textbf{Médiation technique (Outil) :} Il n'utilise pas ses ongles. Il manie le ciseau, la massette, le niveau, l'équerre. Ces outils sont le \textbf{dépôt de savoirs accumulés} par des générations d'artisans. Ils prolongent sa force et sa précision. La technique est le « bras séculaire » de l'intelligence.
    \item \textbf{Effort et Contrainte (Peine) :} La poussière, le bruit, la fatigue physique, le respect des cotes millimétriques, les horaires du chantier : c'est la dimension \emph{tripalium}, la peine réelle, le coût physique payé pour la réalisation.
    \item \textbf{Socialisation (Échange/Reconnaissance) :} Il ne travaille pas seul ni pour lui seul. Son ouvrage s'insère dans un tout collectif (la cathédrale). Il est payé (échange salaire/travail), reconnu par ses pairs (compagnonnage), et son œuvre dépasse sa vie (héritage culturel).
\end{enumerate}
\textbf{Conclusion :} Cette activité n'est pas une simple occupation (comme casser des cailloux pour passer le temps), ni un jeu (il ne peut pas « recommencer » ou « arrêter » sans conséquence), ni une activité animale (pas d'instinct de tailleur de pierre). C'est du \textbf{travail} au sens fort : l'homme se réalise en objectivant son humanité dans la matière.
\end{exoresolu}

\courssub{Synthèse : le travail, acte fondateur de l'humain}

\begin{aretenir}[Ce qu'il faut retenir de la section 1]
\begin{itemize}
    \item \textbf{Étymologie :} \emph{Tripalium} $\rightarrow$ travail = peine, contrainte, souffrance originelle.
    \item \textbf{Définition :} Activité \textbf{consciente, finalisée, technique, socialisée} par laquelle l'homme transforme la nature pour produire de l'utile et se produire lui-même.
    \item \textbf{Spécificité humaine :} \textbf{Projet} (idée antérieure) + \textbf{Outil médiatisé} + \textbf{Transmission culturelle} (vs instinct animal).
    \item \textbf{Distinctions :} Travail $\neq$ Jeu (gratuité/liberté) ; Travail $\neq$ Loisir (temps libéré) ; Travail $\neq$ Simple occupation (contrainte/fin extérieure).
    \item \textbf{Structure invariante :} Finalité -- Effort -- Technique -- Socialisation (les 4 piliers).
\end{itemize}
\end{aretenir}

\begin{exercice}[Entraînement : Appliquer les notions]
\begin{enumerate}
    \item Un élève passe 3 heures à construire une maison dans Minecraft (mode créatif, ressources infinies). Est-ce du travail ? Justifiez en mobilisant les critères (finalité, contrainte, outil, socialisation).
    \item Une femme prépare le ndolè pour sa famille le dimanche midi. Travail ou loisir ? La réponse change-t-elle si elle est cuisinière dans un restaurant ? Argumentez.
    \item Reprenons l'exemple du cultivateur de l'Ouest. Identifiez dans son activité les quatre caractéristiques structurelles (finalité, effort, technique, socialisation) en donnant un exemple précis pour chacune.
\end{enumerate}
\end{exercice}

\begin{corrige}[Exercice n°1 -- Éléments de réponse]
\begin{enumerate}
    \item \textbf{Minecraft (créatif) :} C'est du \textbf{jeu}, pas du travail. Finalité interne (plaisir de construire), pas de contrainte de survie/salaire, outil virtuel sans résistance matérielle réelle, socialisation optionnelle. (En mode « survie » avec grind, la frontière s'estompe : effort simulé, finalité de progression).
    \item \textbf{Ndolè familial :} \textbf{Loisir / Don / Activité domestique non marchande}. Choix du moment, du menu, plaisir du don, pas de contrainte salariale. \textbf{Cuisinière pro :} \textbf{Travail}. Contrainte d'horaires, de normes (hygiène, carte), finalité extérieure (salaire, satisfaction client), hiérarchie, échange marchand.
    \item \textbf{Cultivateur :} Finalité = récolte pour manger/vendre (besoin + projet d'avenir). Effort = pioche, port de charges, soleil, calendrier imposé par les saisons. Technique = terrasses, variétés améliorées, calendrier cultural, machette/houe/tracteur. Socialisation = main-d'œuvre familiale/embauchée, vente au marché (échange monétaire), tontine pour intrants, transmission aux fils.
\end{enumerate}
\end{corrige}

\coursec{Les valeurs du travail}

Dans la section précédente, nous avons défini le travail comme une activité humaine, consciente et finalisée, par laquelle l'homme transforme la nature pour satisfaire ses besoins. Mais une question demeure : pourquoi les hommes attachent-ils tant d'importance au travail ? Pourquoi celui qui perd son emploi se sent-il souvent diminué, au-delà de la simple perte de revenu ? Pourquoi, lorsqu'on fait connaissance, demande-t-on presque toujours : \emph{« Que fais-tu dans la vie ? »} Ces questions montrent que le travail n'est pas seulement une activité : il possède une \cle{valeur}. C'est ce que nous allons maintenant étudier.

\begin{definition}[La valeur]
La \cle{valeur} d'une chose est l'importance qu'on lui accorde, ce qui la rend digne d'estime ou de désir. Parler de la \cle{valeur du travail}, c'est donc rechercher ce qui rend le travail précieux pour l'individu et pour la société, au-delà de la fatigue qu'il coûte.
\end{definition}

On distingue classiquement quatre valeurs du travail : la valeur \cle{économique}, la valeur \cle{sociale}, la valeur \cle{morale} et la valeur \cle{existentielle}. Nous les examinerons l'une après l'autre, avant de voir comment le chômage, expérience négative, révèle par contraste la valeur du travail.

\courssub{La valeur économique : le travail, source de richesse et de revenu}

\textbf{L'intuition d'abord.} Pourquoi la plupart des gens travaillent-ils ? La première réponse qui vient à l'esprit est simple : pour gagner leur vie. Le motocycliste taxi de Douala, la vendeuse du marché Mokolo à Yaoundé, l'enseignant, l'infirmier : tous attendent de leur activité un \cle{revenu} qui leur permet de se nourrir, de se loger, d'éduquer leurs enfants.

\textbf{L'énoncé rigoureux.} La valeur économique du travail tient à ce qu'il est la source des \cle{richesses}. Les économistes classiques l'ont montré : pour Adam Smith (1723-1790), dans \emph{Recherches sur la nature et les causes de la richesse des nations} (1776), c'est le travail, et non l'or ou la terre, qui est la source première de la richesse d'un pays. Pour l'individu, cette richesse prend la forme du \cle{salaire}, c'est-à-dire la rémunération que le travailleur reçoit en échange de sa force de travail.

\begin{exemplebox}[Le salaire comme contrepartie du travail]
Une couturière de Bafoussam transforme du tissu en pagnes confectionnés. Elle vend le fruit de son travail et en tire un revenu. Un employé d'une entreprise de Buea reçoit un salaire mensuel fixe. Dans les deux cas, le travail produit une richesse qui revient, au moins en partie, au travailleur. Sans travail, pas de production ; sans production, pas de revenu.
\end{exemplebox}

\begin{attention}[Ne pas réduire le travail à son salaire]
Piège fréquent dans les dissertations : identifier la valeur du travail à sa seule valeur économique. Or une mère qui éduque ses enfants, un bénévole qui nettoie son quartier, un élève qui révise ses leçons travaillent réellement, sans recevoir aucun salaire. Le travail dépasse donc l'emploi rémunéré : le salaire est une valeur du travail, non \emph{la} valeur du travail.
\end{attention}

\courssub{La valeur sociale : le travail donne un statut et une place dans la société}

\textbf{L'intuition d'abord.} Observez une conversation entre adultes qui se rencontrent pour la première fois. Après le nom, la question qui revient presque toujours est : \emph{« Que fais-tu dans la vie ? »} Cette question apparemment banale est très révélatrice : on identifie une personne par son travail. Être médecin, enseignant, mécanicien ou cultivateur, ce n'est pas seulement exercer une activité, c'est occuper une \cle{place} dans la société.

\textbf{L'énoncé rigoureux.} La valeur sociale du travail tient à ce qu'il confère à l'individu un \cle{statut social}, c'est-à-dire une position reconnue dans le groupe, avec des droits, des devoirs et une certaine considération. Le sociologue Émile Durkheim (1858-1917) a montré dans \emph{De la division du travail social} (1893) que la division du travail crée la \cle{solidarité} entre les hommes : chacun dépend du travail des autres, et chacun est utile aux autres par son propre travail. Le travail est ainsi un lien social : il nous fait exister aux yeux des autres.

\begin{exemplebox}[Le travail comme carte d'identité sociale]
Dans nos villes et villages camerounais, on présente quelqu'un en disant : « C'est le maître d'école », « C'est le chef mécanicien », « C'est la bayam-sellam du marché ». Le métier devient presque un second nom. Inversement, celui qui n'a pas d'activité reconnue peine à se présenter : il se sent « invisible » socialement.
\end{exemplebox}

\courssub{La valeur morale : le travail forme le caractère}

\textbf{L'intuition d'abord.} Les parents et les enseignants répètent souvent : « Le travail forme l'homme. » Un élève qui s'entraîne régulièrement à résoudre des exercices apprend la patience ; un apprenti menuisier qui rate plusieurs fois sa pièce avant de réussir apprend la persévérance. Le travail ne produit pas seulement des objets : il produit des \emph{qualités}.

\textbf{L'énoncé rigoureux.} La valeur morale du travail tient à ce qu'il éduque le caractère. En exigeant un effort régulier, le respect de règles et la capacité de différer le plaisir, le travail enseigne la \cle{discipline}, la \cle{patience}, la \cle{persévérance} et le sens de l'effort. C'est pourquoi les moralistes traditionnels opposent le travail, vertu, à la paresse et à l'oisiveté, vices. L'adage populaire « l'oisiveté est la mère de tous les vices » exprime cette conviction : celui qui ne travaille pas risque de se perdre, tandis que le travail structure la vie et donne des habitudes solides.

\begin{exemplebox}[L'apprentissage comme formation morale]
Un élève de Première technique qui prépare le Probatoire doit travailler chaque jour, renoncer à certains loisirs, recommencer après un échec à un devoir. Même s'il échoue, il aura acquis des qualités morales — régularité, courage, organisation — qui lui serviront toute sa vie, dans n'importe quel métier.
\end{exemplebox}

\courssub{La valeur existentielle : par le travail, l'homme se réalise et humanise le monde}

C'est la valeur la plus profonde, et la plus difficile. Elle demande un effort de réflexion, mais elle est décisive pour la dissertation.

\textbf{L'intuition d'abord.} Un potier contemple le vase qu'il vient de fabriquer : il y reconnaît \emph{quelque chose de lui-même}. Un jeune informaticien de Yaoundé qui crée sa première application éprouve une fierté qui dépasse l'argent gagné : il a l'impression d'exister davantage, de laisser une trace dans le monde. Le travail serait donc le moyen par lequel l'homme \emph{se réalise}.

\textbf{L'énoncé rigoureux : la leçon de Hegel.} Le philosophe allemand Georg Wilhelm Friedrich Hegel (1770-1831), dans \emph{La Phénoménologie de l'Esprit} (1807), a donné du travail une analyse célèbre. Dans la dialectique du maître et de l'esclave, le maître jouit des choses sans les transformer, tandis que l'esclave est contraint de travailler. Or, paradoxalement, c'est l'esclave qui progresse : en travaillant, il transforme les choses, et en transformant les choses, il se transforme lui-même.

\begin{propriete}[Hegel : le travail, « désir retenu »]
Selon Hegel, le travail est un \emph{désir retenu} : au lieu de détruire immédiatement la chose pour la consommer (comme le fait le maître), le travailleur retient son désir, attend, et transforme la chose selon une idée. Par ce détour, il imprime sa marque dans le monde et prend conscience de ses propres pouvoirs : le travail est \emph{formateur de la conscience de soi}. En humanisant la nature, l'homme s'humanise lui-même. (Propriété admise, sans démonstration exigible.)
\end{propriete}

Autrement dit : l'objet fabriqué est comme un miroir. En le regardant, le travailleur se dit : « C'est moi qui ai fait cela ; je suis capable de transformer le monde. » Le travail a ainsi une valeur \cle{existentielle} : il permet à l'homme de se réaliser, c'est-à-dire de devenir pleinement ce qu'il est, et d'humaniser le monde en y inscrivant sa culture, ses techniques et ses valeurs. Les routes, les ponts, les champs, les villes du Cameroun sont autant de traces de ce travail humanisant.

\begin{popfigure}
\begin{center}
\begin{tikzpicture}
% Pyramide des valeurs du travail
\fill[popgreenL]  (-3,0) -- (3,0) -- (2,1.6) -- (-2,1.6) -- cycle;
\fill[poporangeL] (-2,1.6) -- (2,1.6) -- (1.1,3.2) -- (-1.1,3.2) -- cycle;
\fill[popblueL]   (-1.1,3.2) -- (1.1,3.2) -- (0,4.8) -- cycle;
\draw[popdark,thick] (-3,0) -- (3,0) -- (0,4.8) -- cycle;
\draw[popdark] (-2,1.6) -- (2,1.6);
\draw[popdark] (-1.1,3.2) -- (1.1,3.2);
\node at (0,0.8) {\textbf{Survie} : revenu, salaire};
\node at (0,2.4) {\textbf{Statut} : place dans la société};
\node at (0,3.9) {\textbf{Épanouissement}};
\node at (0,3.55) {\small se réaliser (Hegel)};
\draw[-{Stealth},popmuted,thick] (3.4,0.3) -- (3.4,4.5) node[midway,right,align=left] {\small valeurs\\ \small croissantes};
\end{tikzpicture}
\end{center}
\legende{Pyramide des valeurs du travail : le travail assure d'abord la survie (valeur économique), puis donne un statut (valeur sociale), et permet enfin l'épanouissement de soi (valeurs morale et existentielle).}
\end{popfigure}

\courssub{Le chômage : l'expérience négative qui révèle la valeur du travail}

\textbf{La démarche.} Comment savoir ce que vaut le travail ? Une méthode consiste à observer ce qui se passe lorsqu'il \emph{manque}. Le \cle{chômage} — la situation de celui qui voudrait travailler et ne trouve pas d'emploi — fonctionne comme une expérience négative : il révèle, par contraste, toutes les valeurs du travail.

Observons. Le chômeur perd d'abord un \emph{revenu} : c'est la valeur économique qui fait défaut. Mais il perd aussi un \emph{statut} : il ne sait plus quoi répondre à la question « Que fais-tu dans la vie ? » ; il se sent exclu des échanges sociaux, honteux parfois de sa situation. Il perd encore le rythme moral que donne le travail : les journées sans structure favorisent le découragement. Enfin, privé de la possibilité de transformer le monde, il éprouve un sentiment d'inutilité, une \emph{perte de dignité}. Le chômage est donc une \cle{exclusion sociale} autant qu'une privation d'argent : il prouve que le travail vaut bien plus que le salaire.

\begin{exemplebox}[La situation camerounaise : jeunes diplômés et secteur informel]
Au Cameroun, de nombreux jeunes, même diplômés de l'université, peinent à trouver un emploi salarié correspondant à leur formation. Beaucoup se tournent alors vers le \cle{secteur informel} : les \emph{benskineurs} (conducteurs de motos-taxis), les \emph{call-boxeurs} (vendeurs de crédit téléphonique), les vendeurs ambulants, les petits artisans. Ces activités, souvent exercées sans contrat ni salaire fixe, montrent deux choses : d'une part, la souffrance du chômage et du sous-emploi ; d'autre part, la capacité du travail, même précaire, à rendre sa dignité à celui qui l'exerce — le benskineur qui gagne sa vie honnêtement est plus estimé que le jeune qui reste inactif.
\end{exemplebox}

\begin{exemplebox}[Un exemple chiffré]
Au Cameroun, le secteur informel emploie, selon les enquêtes de l'INS, la très grande majorité de la population active (environ 9 actifs sur 10 y travaillent, souvent sans salaire fixe ni protection sociale). Cette réalité confirme qu'il faut distinguer le \emph{travail} de l'\emph{emploi salarié} : des millions de Camerounais travaillent dur chaque jour sans jamais recevoir de bulletin de paie. La valeur du travail dépasse donc largement le cadre du salariat.
\end{exemplebox}

\begin{popfigure}
\begin{center}
\begin{tikzpicture}
\begin{axis}[
  width=0.85\linewidth, height=6.2cm,
  ybar, bar width=16pt,
  ymin=0, ymax=45,
  ylabel={Pourcentage de réponses (\%)},
  symbolic x coords={Revenu,Statut,Epanouissement,UtiliteSociale},
  xtick=data,
  xticklabels={Revenu,Statut,Épanouissement,Utilité sociale},
  x tick label style={font=\small,align=center},
  nodes near coords, nodes near coords style={font=\small},
  axis lines=left,
  enlarge x limits=0.18,
]
\addplot[fill=popblueL,draw=popblue] coordinates {(Revenu,38) (Statut,22) (Epanouissement,25) (UtiliteSociale,15)};
\end{axis}
\end{tikzpicture}
\end{center}
\legende{Enquête fictive auprès d'élèves de Première technique : « Qu'est-ce qui fait surtout la valeur du travail ? » Le revenu arrive en tête (38 \%), mais les valeurs non économiques réunies (statut, épanouissement, utilité sociale) totalisent 62 \% des réponses — ce qui confirme que le travail ne se réduit pas au salaire. (Données fictives, à but pédagogique.)}
\end{popfigure}

\begin{savaistu}[Le mot « travail » désignait la souffrance]
Le mot \emph{travail} vient du latin populaire \emph{trepalium}, qui désignait un instrument de torture. Au Moyen Âge, travailler était le lot pénible des serfs, et l'Antiquité grecque méprisait le travail manuel, réservé aux esclaves, car elle estimait que l'homme libre devait se consacrer à la pensée et à la vie politique. Ce n'est qu'avec la modernité — la Réforme protestante, puis la révolution industrielle — que le travail devient une valeur positive, source de dignité, et finalement un \emph{droit} : la Déclaration universelle des droits de l'homme de 1948 proclame que « toute personne a droit au travail ». L'histoire du mot illustre le chemin que nous avons parcouru : d'une malédiction, le travail est devenu une valeur.
\end{savaistu}

\begin{exoresolu}[Appliquer les notions à une situation concrète]
\textbf{Énoncé.} Marlyse, jeune diplômée d'un BTS en comptabilité à Douala, ne trouve pas d'emploi salarié depuis deux ans. Elle décide d'ouvrir un petit stand de vente de beignets-haricots devant un lycée. Son activité marche bien, mais elle dit à sa sœur : « Je fais ce métier seulement pour l'argent, ce n'est pas un vrai travail. » En mobilisant les quatre valeurs du travail, montrez que l'analyse de Marlyse est incomplète.
\tcblower
\textbf{Solution détaillée.}
\begin{enumerate}
\item \textbf{Valeur économique.} Marlyse a raison sur un point : son activité lui procure un revenu qui assure sa survie et son autonomie financière. C'est la première valeur du travail.
\item \textbf{Valeur sociale.} Mais elle a tort de dire que ce n'est pas un « vrai travail ». Son stand lui donne une place reconnue dans son quartier : les élèves et les enseignants la connaissent, elle est utile chaque matin à des centaines de personnes. Elle possède un statut social.
\item \textbf{Valeur morale.} Se lever chaque jour avant l'aube, gérer ses achats, tenir ses comptes (elle y utilise d'ailleurs sa formation !) exige discipline, régularité et persévérance : son activité forme son caractère.
\item \textbf{Valeur existentielle.} En créant son entreprise de toutes pièces, Marlyse transforme une situation de chômage subi en projet maîtrisé : au sens de Hegel, elle imprime sa marque dans le monde et prend conscience de ses capacités. Elle se réalise.
\end{enumerate}
\textbf{Conclusion.} L'erreur de Marlyse est de réduire le travail à l'emploi salarié et au salaire. Or son activité, même informelle, possède les quatre valeurs du travail : c'est bien un vrai travail, qui la fait vivre, la situe socialement, la forme moralement et la réalise comme personne.
\end{exoresolu}

\begin{aretenir}[L'essentiel de la section]
\begin{itemize}
\item La \cle{valeur} d'une chose est l'importance qu'on lui accorde ; le travail possède quatre valeurs principales.
\item Valeur \cle{économique} : le travail est source de richesse et de revenu (le salaire) — mais il ne s'y réduit pas.
\item Valeur \cle{sociale} : le travail donne un statut, une place reconnue dans la société (« Que fais-tu dans la vie ? »).
\item Valeur \cle{morale} : le travail forme le caractère (discipline, patience, persévérance).
\item Valeur \cle{existentielle} : selon Hegel, le travail, « désir retenu », est formateur de la conscience de soi ; en transformant le monde, l'homme se transforme et se réalise.
\item Le \cle{chômage} révèle par contraste ces valeurs : il entraîne exclusion sociale et perte de dignité, comme le montre la situation des jeunes diplômés camerounais, dont beaucoup retrouvent dans le secteur informel (benskineurs, call-boxeurs) un travail qui dépasse le seul salaire.
\end{itemize}
\end{aretenir}

\coursec{Travail : malédiction, libération ou aliénation ?}

Nous venons de voir que le travail possède une \cle{valeur} multiple : valeur économique (il produit des richesses), valeur sociale (il intègre l'individu dans la société) et valeur morale (il forme le caractère). Mais ces valeurs positives sont-elles universelles ? L'expérience quotidienne semble parfois les contredire : le carrier qui casse des pierres sous le soleil de Yaoundé, l'ouvrier qui répète le même geste huit heures par jour, le mototaximan exposé aux accidents — tous travaillent, mais éprouvent-ils la dignité et l'épanouissement que promet la valeur du travail ? Cette contradiction nous oblige à problématiser : le travail est-il une \cle{malédiction} qui pèse sur l'homme, une \cle{libération} qui le rend maître de son destin, ou une \cle{aliénation} qui le dépossède de lui-même ? Nous examinerons successivement ces trois thèses, puis nous verrons qu'elles peuvent se concilier si l'on distingue les \emph{conditions} dans lesquelles le travail s'exerce.

\courssub{Première thèse : le travail comme malédiction}

\textbf{L'intuition de départ.} Demandez à un élève de première s'il préférerait passer sa journée à jouer ou à travailler aux champs sous la chaleur : la réponse est presque toujours immédiate. Le travail fatigue, le travail fait mal, le travail s'impose. Cette expérience commune fonde la première thèse : le travail serait une \emph{punition}, un fardeau que l'homme subit et dont il rêve de se débarrasser.

\textbf{Le fondement religieux : la Genèse.} Cette conception trouve son expression la plus célèbre dans la \cle{Bible}. Après la désobéissance d'Adam et Ève, Dieu déclare à l'homme : « Tu gagneras ton pain à la sueur de ton front » (Genèse, III, 19). Le travail apparaît ici comme la \emph{conséquence du péché originel} : dans le jardin d'Éden, avant la faute, l'homme ne travaillait pas ; la terre donnait ses fruits spontanément. Le travail est donc pensé comme une \cle{malédiction}, c'est-à-dire une condamnation qui accompagne la condition humaine déchue. L'étymologie latine semble confirmer cette vision : le mot \emph{travail} viendrait de \emph{trepalium}, un instrument de torture. Travailler, ce serait littéralement être torturé.

\textbf{Les arguments de cette thèse.} On peut la justifier par trois observations concrètes :
\begin{itemize}
\item La \cle{pénibilité} physique : le travail épuise le corps. Le mineur, le docker du port de Douala, la lavandière, usent leur santé dans l'effort répété.
\item La \cle{fatigue} et la souffrance : le travail consomme le temps de vie. Les heures passées au travail sont des heures soustraites au repos, à la famille, au plaisir.
\item La \cle{contrainte} : on ne travaille généralement pas par choix mais par nécessité. « Il faut bien vivre », dit le langage courant, ce qui signifie que le travail est subi comme une obligation extérieure.
\end{itemize}

\begin{exemplebox}[La journée du carrier]
Dans les carrières de pierre en périphérie des grandes villes camerounaises, des hommes cassent des blocs de granit à la masse, dix heures par jour, pour un revenu très faible. Dos voûté, mains abîmées, poussière dans les poumons : leur travail ressemble trait pour trait à la malédiction biblique — ils gagnent littéralement leur pain « à la sueur de leur front ». Difficile de parler ici d'épanouissement.
\end{exemplebox}

\textbf{Limite de la thèse.} Si le travail n'était qu'une malédiction, comment expliquer que le chômeur, précisément privé de travail, se sente malheureux et humilié ? Le paradoxe du chômage suggère que le travail n'est pas seulement une punition : il serait aussi un besoin. C'est ce que défend la deuxième thèse.

\courssub{Deuxième thèse : le travail comme libération}

\textbf{L'intuition de départ.} Observez un jeune qui fabrique de ses mains un meuble, répare un moteur ou récolte le champ qu'il a semé : sur son visage, on lit la fierté. Le travail ne fait pas que fatiguer ; il accomplit. C'est l'expérience que théorisent les philosophes pour qui le travail est le propre de l'homme.

\textbf{Hegel : le travail éduque et affranchit.} Pour \cle{Hegel} (1770-1831), dans la célèbre dialectique du maître et de l'esclave (\emph{Phénoménologie de l'Esprit}, 1807), le travail est le chemin de la liberté. Le maître jouit sans travailler, mais il reste dépendant de l'esclave ; l'esclave, lui, en travaillant, \emph{transforme la nature} et, ce faisant, se transforme lui-même. En façonnant les objets, il découvre sa propre puissance : il se reconnaît dans son œuvre. Le travail est ainsi une \cle{éducation} : il discipline les désirs, développe les capacités, et prépare l'esclave à devenir un homme libre et autonome. La leçon est claire : c'est par le travail que l'homme s'arrache à la pure animalité et devient pleinement humain.

\textbf{Le travail affranchit de la nature.} L'animal subit la nature : il dépend des saisons, des proies, du hasard. L'homme, par le travail, \emph{domine} la nature : il cultive au lieu de cueillir, construit des maisons au lieu de chercher des abris, fabrique des outils qui démultiplient sa force. Le travail est donc le moyen par lequel l'humanité conquiert son \cle{autonomie} : plus elle travaille, moins elle dépend des caprices du milieu naturel.

\textbf{Marx (en partie) : le travail, essence de l'homme.} On cite souvent Marx comme le critique du travail ; c'est une erreur que nous corrigerons plus bas. En réalité, \cle{Marx} (1818-1883) voit dans le travail l'\emph{essence même de l'homme} : par le travail, l'homme « objectivise » ses idées — il projette dans la matière ce qu'il a d'abord conçu dans sa tête, comme l'architecte qui imagine la maison avant de la bâtir. Le travail libre, créateur, est pour Marx la plus haute activité humaine, celle qui développe toutes les capacités de l'individu.

\begin{exemplebox}[Le menuisier de quartier]
Le menuisier ébéniste de Bafoussam qui conçoit, découpe, assemble et polit une armoire vit une expérience de libération : il choisit son projet, maîtrise chaque étape, voit naître l'objet sous ses mains et le vend lui-même. Son travail développe son habileté, sa réputation, son autonomie financière. Il se reconnaît dans son œuvre : « C'est moi qui l'ai faite. »
\end{exemplebox}

\textbf{Limite de la thèse.} Cette belle image du travail libérateur correspond-elle à la réalité de l'ouvrier moderne ? Celui qui, sur une chaîne de montage, visse toute la journée la même vis sans jamais voir le produit fini, se reconnaît-il dans son œuvre ? C'est ici qu'intervient la troisième thèse, la plus critique.

\courssub{Troisième thèse : le travail comme aliénation}

\begin{definition}[L'aliénation]
L'\cle{aliénation} (du latin \emph{alienus}, « autre, étranger ») est le processus par lequel l'homme devient \emph{étranger} à lui-même, à ce qu'il produit et aux autres hommes. Ce qu'il crée se retourne contre lui : au lieu de s'accomplir dans son œuvre, il s'y perd. On dit alors que le travailleur est « aliéné » : dépossédé de sa propre activité.
\end{definition}

\textbf{L'intuition de départ.} Imaginez un ouvrier d'une usine de chocolat qui fabrique toute la journée des tablettes qu'il ne pourra jamais s'offrir, pour un patron qu'il ne connaît pas, selon des gestes qu'on lui impose. À qui appartient le chocolat ? Pas à lui. À qui appartient son geste ? Pas à lui non plus : c'est la machine qui dicte la cadence. Que reste-t-il de lui dans son travail ? Presque rien. C'est cette expérience que Marx analyse dans les \emph{Manuscrits de 1844}.

\begin{propriete}[Le paradoxe marxiste de la production capitaliste]
Marx montre que dans le capitalisme, \emph{plus l'ouvrier produit de richesses, plus il s'appauvrit}. \textbf{Justification :} le produit du travail n'appartient pas à l'ouvrier mais au propriétaire de l'usine (le capitaliste), qui ne verse à l'ouvrier qu'un salaire fixe, indépendant de la valeur créée. Donc, plus l'ouvrier produit, plus il enrichit autrui, plus le monde des objets qui lui fait face devient puissant — et plus lui-même devient pauvre et dépendant, proportionnellement. La richesse qu'il produit se dresse contre lui comme une puissance étrangère.
\end{propriete}

\textbf{Les quatre aliénations.} Marx distingue quatre dimensions de cette dépossession :
\begin{enumerate}
\item \textbf{Aliénation au \cle{produit}} : l'objet fabriqué n'appartient pas à celui qui l'a fait. L'ouvrier produit des biens dont il ne jouit pas.
\item \textbf{Aliénation à l'\cle{acte de production}} : le travail lui-même n'est pas choisi ; le geste, le rythme, l'horaire sont imposés de l'extérieur. Le travail n'est pas épanouissement mais « travail forcé » : l'ouvrier ne se sent libre que hors du travail.
\item \textbf{Aliénation à l'\cle{essence humaine}} : l'homme se distingue de l'animal par l'activité libre et créatrice. Réduit à répéter un geste mécanique, l'ouvrier est ramené au niveau d'une simple bête de somme : il est déshumanisé.
\item \textbf{Aliénation à \cle{autrui}} : au lieu d'unir les hommes dans la coopération, le travail salarié les oppose : concurrence entre ouvriers pour l'emploi, opposition entre l'ouvrier et le patron qui s'enrichit de sa peine.
\end{enumerate}

\begin{popfigure}
\begin{center}
\begin{tikzpicture}[node distance=0.4cm,
boite/.style={draw=popink, fill=popblueL, rounded corners, text width=4.6cm, align=center, font=\small, inner sep=5pt},
centre/.style={draw=popdark, fill=poporangeL, rounded corners, text width=3.6cm, align=center, font=\small\bfseries, inner sep=6pt}]
\node[centre] (C) {L'ouvrier aliéné (Marx, 1844)};
\node[boite, above left=0.7cm and 1.2cm of C, align=center] (P){\textbf{1. Au produit}\\ Le produit ne lui appartient pas};
\node[boite, below left=0.7cm and 1.2cm of C, align=center] (A){\textbf{2. À l'acte de production}\\ Geste et rythme imposés};
\node[boite, above right=0.7cm and 1.2cm of C, align=center] (E){\textbf{3. À l'essence humaine}\\ Déshumanisation, geste mécanique};
\node[boite, below right=0.7cm and 1.2cm of C, align=center] (H){\textbf{4. À autrui}\\ Concurrence et opposition patron/ouvrier};
\draw[-{Stealth}, thick, popink] (C) -- (P);
\draw[-{Stealth}, thick, popink] (C) -- (A);
\draw[-{Stealth}, thick, popink] (C) -- (E);
\draw[-{Stealth}, thick, popink] (C) -- (H);
\end{tikzpicture}
\end{center}
\legende{Les quatre aliénations du travailleur selon Marx : le travail salarié capitaliste dépossède l'ouvrier de son produit, de son geste, de son humanité et de son rapport aux autres.}
\end{popfigure}

\begin{exemplebox}[L'ouvrier de l'usine de cacao]
Dans les usines de transformation du cacao de la région du Centre, un ouvrier peut passer ses journées à surveiller des machines qui produisent du chocolat de marque exportée en Europe. Son salaire mensuel ne lui permet pas d'acheter régulièrement le produit fini qu'il a lui-même contribué à fabriquer. Voilà une illustration parfaite de l'\cle{aliénation au produit} : ce que ses mains créent devient un objet étranger, voire inaccessible.
\end{exemplebox}

\begin{attention}[Ne pas caricaturer Marx]
Erreur fréquente dans les copies : présenter Marx comme un \emph{ennemi du travail}. C'est exactement l'inverse ! Marx \textbf{valorise} le travail, qu'il considère comme l'essence créatrice de l'homme. Ce qu'il critique, ce sont les \emph{conditions capitalistes} du travail — le salariat, la propriété privée des moyens de production, la cadence imposée — qui pervertissent cette belle activité et la transforment en aliénation. Écrivez donc : « Marx critique non pas le travail, mais le travail \emph{aliéné} ».
\end{attention}

\courssub{La division du travail et le taylorisme : l'homme devenu machine ?}

\textbf{Le phénomène.} Pour produire plus et moins cher, l'industrie moderne a poussé à l'extrême la \cle{division du travail} : au lieu qu'un ouvrier fabrique un objet entier, la fabrication est découpée en tâches minuscules, chacune répétée à l'infini. L'ingénieur américain \cle{Taylor} (1856-1915) a systématisé cette logique : le \cle{taylorisme} consiste à chronométrer les gestes, à éliminer tout mouvement inutile, et à séparer totalement ceux qui \emph{conçoivent} (les ingénieurs) de ceux qui \emph{exécutent} (les ouvriers). La \cle{chaîne de montage}, popularisée par Ford, en est l'aboutissement : l'objet défile devant l'ouvrier, qui répète le même geste toute la journée, au rythme de la machine.

\textbf{Le problème philosophique.} Le travail à la chaîne \emph{rend-il l'homme machine} ? Trois arguments le suggèrent :
\begin{itemize}
\item L'ouvrier n'a pas besoin de penser : son geste est programmé comme celui d'un automate. L'intelligence est entièrement du côté de l'ingénieur.
\item C'est la machine qui impose le rythme : l'homme s'adapte à la mécanique, et non l'inverse. Le vivant se règle sur l'inerte.
\item L'ouvrier devient interchangeable : n'importe qui peut remplacer n'importe qui, puisque le geste ne demande aucune qualification. L'individu perd sa singularité.
\end{itemize}

\begin{savaistu}[Charlie Chaplin et \emph{Les Temps modernes}]
Dans le film \emph{Les Temps modernes} (1936), Charlie Chaplin joue un ouvrier chargé de serrer des boulons sur une chaîne de montage dont la cadence ne cesse d'accélérer. Avalé par la machine, il en ressort en continuant à visser tout ce qui ressemble à un boulon — y compris les boutons des passantes — avant de devenir fou. Par le comique, Chaplin dénonce une réalité tragique : le travail à la chaîne, en mécanisant le geste, peut détruire l'équilibre mental de l'ouvrier. Le film reste, près d'un siècle plus tard, la critique la plus célèbre du taylorisme.
\end{savaistu}

\textbf{Nuance.} La division du travail a aussi des effets positifs : elle augmente la production, abaisse les prix, rend accessibles des biens autrefois réservés aux riches. Le problème n'est donc pas la technique elle-même, mais l'organisation qui réduit l'homme au rang de simple rouage.

\courssub{Synthèse : le sens du travail dépend de ses conditions}

Les trois thèses semblent s'opposer, mais chacune dit une part de vérité. La synthèse s'impose d'elle-même : le travail n'est ni bon ni mauvais \emph{en soi} — tout dépend des \cle{conditions} dans lesquelles il s'exerce. La distinction décisive est celle du \textbf{travail choisi} et du \textbf{travail subi}.

\begin{itemize}
\item Le \cle{travail choisi} — libre, créateur, dont le travailleur maîtrise le sens et le produit — est libérateur : il développe les capacités, donne une place dans la société, procure la fierté de l'œuvre accomplie. C'est le travail dont rêvaient Hegel et Marx.
\item Le \cle{travail subi} — imposé par la nécessité, morcelé, dépossédé de son produit et de son sens — est aliénant, voire vécu comme une malédiction : c'est celui que dénoncent la Genèse et Marx.
\end{itemize}

La question philosophique « le travail est-il une malédiction ou une libération ? » se transforme ainsi en question pratique et politique : \emph{comment organiser le travail pour qu'il libère au lieu d'aliéner ?} C'est l'enjeu des revendications pour de meilleures conditions de travail, des salaires justes, de la formation professionnelle.

\begin{popfigure}
\begin{center}
\begin{tikzpicture}[cell/.style={draw=popline, fill=#1, text width=3.9cm, minimum height=0.85cm, align=left, font=\scriptsize, inner sep=4pt},
head/.style={draw=popdark, fill=#1, text width=3.9cm, align=center, font=\small\bfseries, inner sep=5pt}]
\node[head=popblueL] (h1) at (0,0) {Malédiction};
\node[head=popgreenL] (h2) at (4.6,0) {Libération};
\node[head=poporangeL] (h3) at (9.2,0) {Aliénation};
\node[cell=popblueL!40, below=0cm of h1] {La Bible (Genèse) : « tu gagneras ton pain à la sueur de ton front »};
\node[cell=popgreenL!40, below=0cm of h2] {Hegel, Marx : le travail humanise et affranchit de la nature};
\node[cell=poporangeL!40, below=0cm of h3] {Marx : le salariat capitaliste dépossède l'ouvrier};
\node[cell=popblueL!40, below=0.85cm of h1] {Pénibilité, fatigue, contrainte, nécessité};
\node[cell=popgreenL!40, below=0.85cm of h2] {Développe les capacités, autonomie, fierté de l'œuvre};
\node[cell=poporangeL!40, below=0.85cm of h3] {Perte du produit, du geste, de soi et du lien aux autres};
\node[cell=popblueL!40, below=1.7cm of h1] {Ex. : le carrier, le travail forcé};
\node[cell=popgreenL!40, below=1.7cm of h2] {Ex. : l'artisan, le menuisier de Bafoussam};
\node[cell=poporangeL!40, below=1.7cm of h3] {Ex. : la chaîne de montage, l'usine de cacao};
\end{tikzpicture}
\end{center}
\legende{Tableau synoptique des trois thèses sur le travail : auteurs, arguments et exemples. La synthèse : le sens du travail dépend de ses conditions (choisi ou subi).}
\end{popfigure}

\begin{methode}[Confronter les thèses avant de conclure]
Pour traiter une question philosophique comme « Le travail est-il une libération ? », respectez la démarche suivante :
\begin{enumerate}
\item \textbf{Définir} les termes du sujet (ici : travail, libération).
\item \textbf{Exposer une première thèse} avec son auteur, ses arguments et un exemple (le travail comme malédiction).
\item \textbf{Exposer la thèse opposée} avec la même rigueur (le travail comme libération ; puis la critique marxiste de l'aliénation).
\item \textbf{Confronter} : montrer ce que chaque thèse a de vrai et où elle bute.
\item \textbf{Proposer une synthèse personnelle justifiée} (ici : tout dépend des conditions du travail — choisi ou subi) et conclure en ouvrant éventuellement sur une nouvelle question.
\end{enumerate}
Une copie qui ne défend qu'une seule thèse sans examiner la thèse adverse reste un exposé, non une dissertation philosophique.
\end{methode}

\begin{exoresolu}[Débat argumenté : le mototaxi, libération ou aliénation ?]
\textbf{Énoncé.} Au Cameroun, des centaines de milliers de jeunes exercent le métier de conducteur de mototaxi (« bendskin » à Douala, « moto » à Yaoundé). En mobilisant les trois thèses étudiées, discutez : ce travail est-il une libération ou une aliénation ? Prenez position et justifiez.
\tcblower
\textbf{Solution détaillée.} Appliquons la méthode de la confrontation des thèses.

\emph{Thèse de la libération.} Le mototaxi offre au jeune sans diplôme une \textbf{autonomie} remarquable : il choisit ses horaires, n'a pas de patron direct, gagne un revenu quotidien qui le soustrait au chômage et à la mendicité. Il développe des capacités (conduite, mécanique, sens du commerce) et occupe une place reconnue dans la vie de la cité : il transporte les élèves, les malades, les marchandises. Au sens de Hegel, il se réalise par son activité et conquiert son indépendance matérielle.

\emph{Thèse de l'aliénation.} Mais ce travail est aussi marqué par la \textbf{précarité} : pas de salaire fixe, pas d'assurance maladie, pas de retraite. Le conducteur est exposé aux accidents mortels, aux intempéries, au racket, à l'insécurité. Souvent, la moto ne lui appartient même pas : il doit reverser chaque jour une « recette » au propriétaire — situation proche de l'aliénation au produit, puisque le fruit de sa peine enrichit d'abord autrui. Sa fatigue est grande, sa marge de choix réelle est faible : il « choisit » ce métier faute de mieux.

\emph{Synthèse personnelle.} Le mototaxi illustre parfaitement notre conclusion générale : c'est la \textbf{condition} du travail qui fait son sens. Tant que le conducteur subit la précarité et le danger sans protection, son travail reste proche du travail subi ; mais lorsqu'il est propriétaire de sa moto, organisé en coopérative, protégé par une assurance, ce même travail devient un instrument d'autonomie et de dignité. La libération ne tient donc pas au métier lui-même, mais aux conditions sociales qui l'encadrent.
\end{exoresolu}

\begin{aretenir}[L'essentiel de la section]
\begin{itemize}
\item Trois thèses s'affrontent : le travail comme \cle{malédiction} (Genèse : punition, pénibilité, fatigue), comme \cle{libération} (Hegel : il éduque et affranchit de la nature ; Marx : il est l'essence créatrice de l'homme), comme \cle{aliénation} (Marx : le salariat dépossède l'ouvrier de son produit, de son acte, de son essence et d'autrui).
\item Le \cle{taylorisme} et la chaîne de montage poussent l'aliénation à l'extrême en réduisant l'ouvrier à un geste mécanique : l'homme risque de devenir machine.
\item Synthèse : le travail n'est ni bon ni mauvais en soi ; son sens dépend de ses \textbf{conditions} — travail \emph{choisi} (libérateur) ou travail \emph{subi} (aliénant).
\item Erreur à éviter : Marx ne critique pas le travail, mais les conditions capitalistes du travail.
\end{itemize}
\end{aretenir}

\begin{exercice}[Pour s'entraîner]
1. Définissez l'aliénation et citez les quatre formes qu'en distingue Marx.\par
2. Expliquez en cinq lignes le paradoxe : « plus l'ouvrier produit, plus il s'appauvrit ».\par
3. Sujet de dissertation : « Le travail libère-t-il l'homme ? » Rédigez l'introduction (amorce, définitions, problématique, annonce du plan) en suivant la méthode de confrontation des thèses.
\end{exercice}

\begin{corrige}[Corrigé des exercices]
\textbf{1.} L'aliénation est le processus par lequel l'homme devient étranger à lui-même, à son produit et aux autres. Marx distingue : l'aliénation au produit (l'objet fabriqué n'appartient pas à l'ouvrier), l'aliénation à l'acte de production (le geste et le rythme sont imposés), l'aliénation à l'essence humaine (l'activité créatrice est remplacée par un geste mécanique), l'aliénation à autrui (concurrence et opposition entre les hommes).\par
\textbf{2.} Dans le capitalisme, le produit appartient au propriétaire de l'usine, non à l'ouvrier qui ne reçoit qu'un salaire fixe. Donc chaque richesse supplémentaire produite enrichit le capitaliste et renforce le monde des objets face à l'ouvrier, qui, lui, reste proportionnellement pauvre et dépendant. Sa propre production se retourne contre lui comme une puissance étrangère.\par
\textbf{3.} \emph{Éléments d'introduction attendus :} amorce sur l'expérience commune du travail (fatigue mais aussi fierté) ; définition du travail (activité humaine transformant la nature) et de la libération (affranchissement d'une contrainte) ; problématique : le travail affranchit-il l'homme ou l'asservit-il davantage ? ; annonce du plan en trois moments : le travail comme malédiction, le travail comme libération, la synthèse fondée sur la distinction entre travail choisi et travail subi.
\end{corrige}

\coursec{Qu'est-ce que l'échange ?}

Après avoir examiné le travail sous ses trois visages possibles — malédiction, libération ou aliénation —, une question s'impose : que devient le produit du travail ? L'homme ne travaille jamais seul et ne consomme jamais tout ce qu'il produit. Le cultivateur de l'Ouest échange son maïs contre du poisson venu de la côte ; l'élève échange ses paroles avec ses camarades ; la famille offre des cadeaux lors des funérailles. Autrement dit, le travail plonge l'homme dans un réseau d'\cle{échanges}. Mais qu'est-ce qu'échanger, au juste ? Et pourquoi l'échange est-il bien plus qu'une simple affaire de marché ?

\courssub{Définition générale de l'échange}

Commençons par l'intuition. Tu sors de chez toi le matin : tu salues ton voisin, tu achètes des beignets à la mamie du carrefour, tu discutes avec un ami, tu donnes un coup de main à ton frère. Sans le savoir, tu as déjà participé à quatre formes d'échanges différentes. L'échange n'est donc pas d'abord une opération de calcul : c'est le fait même de vivre avec les autres.

\begin{definition}[L'échange]
L'\cle{échange} est un \textbf{transfert réciproque} de biens, de services, de paroles ou de symboles entre des personnes ou des groupes. Il suppose au moins deux partenaires, un objet (matériel ou immatériel) qui circule, et une relation qui se crée ou se renforce entre eux.
\end{definition}

Cette définition retient trois éléments essentiels. D'abord, la \textbf{réciprocité} : l'échange n'est pas un vol (transfert forcé) ni une simple perte ; il va dans les deux sens, même si le retour est différé. Ensuite, la \textbf{diversité de l'objet échangé} : on échange des marchandises, mais aussi des paroles, des salutations, des femmes (dans les alliances matrimoniales traditionnelles), des honneurs. Enfin, la \textbf{relation sociale} : échanger, c'est toujours nouer ou entretenir un lien.

\begin{attention}[Piège fréquent]
Ne réduis jamais l'échange au \textbf{commerce}. Acheter du pain au marché est un échange, certes, mais offrir un cadeau à un ami, saluer un voisin, participer à une conversation ou assister aux funérailles d'un parent sont \emph{déjà} des échanges. Dans une dissertation, l'élève qui ne parle que d'argent et de marché montre qu'il n'a pas compris la richesse de la notion.
\end{attention}

\courssub{La triple obligation de Marcel Mauss : donner, recevoir, rendre}

L'analyse la plus célèbre de l'échange est celle du sociologue et ethnologue français \textbf{Marcel Mauss} (1872-1950), dans son \emph{Essai sur le don} (1925). En étudiant les sociétés traditionnelles (Pacifique, Amérique du Nord), Mauss découvre que le cadeau, en apparence gratuit, obéit en réalité à une loi rigoureuse qu'il appelle la \cle{triple obligation} :

\begin{enumerate}
\item \textbf{Donner} : il faut faire le premier pas, offrir quelque chose ; celui qui ne donne jamais se coupe de la communauté.
\item \textbf{Recevoir} : refuser un don, c'est refuser la relation elle-même ; c'est souvent une insulte ou une déclaration d'hostilité.
\item \textbf{Rendre} : celui qui a reçu doit rendre, un jour, d'une manière ou d'une autre ; sinon il perd son honneur et devient débiteur à vie.
\end{enumerate}

\begin{propriete}[Thèse de Mauss]
Selon Mauss, \textbf{le don n'est jamais gratuit} : il crée une obligation de rendre qui tisse le lien social. Le cadeau circule comme un lien invisible entre les personnes et les groupes ; c'est cette circulation qui fait tenir la société ensemble.
\end{propriete}

\begin{popfigure}
\begin{center}
\begin{tikzpicture}[scale=1]
\node[draw=popblue, fill=popblueL, rounded corners, thick, minimum width=2.6cm, minimum height=1cm, font=\bfseries] (donner) at (90:2.6) {DONNER};
\node[draw=poporange, fill=poporangeL, rounded corners, thick, minimum width=2.6cm, minimum height=1cm, font=\bfseries] (recevoir) at (-30:2.6) {RECEVOIR};
\node[draw=popgreen, fill=popgreenL, rounded corners, thick, minimum width=2.6cm, minimum height=1cm, font=\bfseries] (rendre) at (210:2.6) {RENDRE};
\draw[-{Stealth[length=3mm]}, very thick, popblue, bend left=25] (donner) to (recevoir);
\draw[-{Stealth[length=3mm]}, very thick, poporange, bend left=25] (recevoir) to (rendre);
\draw[-{Stealth[length=3mm]}, very thick, popgreen, bend left=25] (rendre) to (donner);
\node[font=\itshape, text=popmuted] at (0,0) {Le lien social};
\end{tikzpicture}
\end{center}
\legende{Le cycle du don selon Marcel Mauss : donner oblige à recevoir, recevoir oblige à rendre, rendre relance le don. Ce cycle sans fin tisse le lien social.}
\end{popfigure}

\begin{exemplebox}[La tontine camerounaise]
Dans la \textbf{tontine}, pratiquée dans tout le Cameroun, les membres d'un groupe cotisent chaque semaine ou chaque mois une somme fixe, et chacun reçoit à tour de rôle la totalité de la cagnotte. L'échange y est double : \emph{économique} (épargne rotative qui permet d'acheter une machine, de payer les frais scolaires) et \emph{symbolique} (solidarité, confiance, appartenance au groupe). Celui qui a reçu sa part doit continuer à cotiser : c'est exactement l'obligation de \textbf{rendre} décrite par Mauss.
\end{exemplebox}

\begin{savaistu}[Le potlatch : l'échange comme rivalité]
Dans certaines sociétés du Pacifique Nord-Ouest étudiées par Mauss, les chefs organisaient des fêtes appelées \textbf{potlatch} au cours desquelles ils distribuaient — et parfois \emph{détruisaient} — d'énormes richesses (couvertures, cuivres, canoës). But : écraser le rival par la générosité, car celui qui ne peut pas rendre davantage perd la face. L'échange devient alors une \textbf{guerre pacifique} par les cadeaux : donner plus que l'autre, c'est affirmer sa supériorité. On retrouve une logique voisine dans certaines cérémonies fastueuses où l'on mesure le prestige à la dépense accomplie.
\end{savaistu}

\courssub{Les formes d'échanges}

On distingue classiquement trois grandes formes d'échanges, qui coexistent dans toutes les sociétés humaines.

\textbf{1. Les échanges économiques.} Ce sont les échanges de biens et de services : le \textbf{troc} (échange direct d'un bien contre un autre : du manioc contre des arachides), puis l'échange monétaire (achat et vente sur le \textbf{marché}). Ici, on cherche une \textbf{équivalence} : combien de kilos de maïs pour un poulet ? Le calcul domine.

\textbf{2. Les échanges linguistiques.} Parler, c'est déjà échanger. Toute conversation est un va-et-vient de paroles : je te salue, tu me réponds ; je pose une question, tu me donnes une information. Le langage est même le \textbf{premier échange} de l'humanité : avant de commercer, les hommes se sont parlé. Entrer en relation avec autrui commence par la parole — c'est pourquoi celui à qui l'on refuse de répondre se sent exclu.

\textbf{3. Les échanges symboliques.} Ce sont les dons et contre-dons, les cadeaux, la \textbf{dot} matrimoniale, les présents offerts lors des funérailles ou des naissances. L'objet échangé vaut moins par sa valeur marchande que par ce qu'il \emph{signifie} : respect, alliance, reconnaissance, amitié.

\begin{popfigure}
\begin{center}
\begin{tikzpicture}[
mindmap, grow cyclic,
concept color=popblue!60,
root concept/.append style={font=\bfseries},
level 1 concept/.append style={level distance=4.2cm, sibling angle=120, font=\bfseries},
level 2 concept/.append style={level distance=2.9cm, sibling angle=50, font=\small},
every node/.style={concept}
]
\node{L'ÉCHANGE}
child[concept color=poporange!60]{ node{Économique}
child{ node{Troc} }
child{ node{Monnaie, marché} }
}
child[concept color=popgreen!60]{ node{Linguistique}
child{ node{Parole, salutation} }
child{ node{Conversation} }
}
child[concept color=poppurple!60]{ node{Symbolique}
child{ node{Don, contre-don} }
child{ node{Dot, cadeaux} }
};
\end{tikzpicture}
\end{center}
\legende{Carte mentale des trois formes d'échanges : économique (biens et services), linguistique (paroles) et symbolique (dons et signes).}
\end{popfigure}

\courssub{Du troc à la monnaie : l'échange marchand}

Dans les sociétés anciennes, on pratiquait le \cle{troc} : échange direct d'un bien contre un autre, sans argent. Le pêcheur de Limbe échange son poisson contre le plantain du cultivateur de Buea. Mais le troc a une limite évidente : il faut que chacun désire \emph{exactement} ce que l'autre offre, au même moment (on parle de « double coïncidence des besoins »). Que faire si le pêcheur veut du plantain mais que le cultivateur ne veut pas de poisson ?

La \cle{monnaie} résout ce problème : elle est un bien intermédiaire accepté par tous, qui permet d'acheter maintenant et de vendre plus tard, à qui l'on veut. Avec la monnaie naît l'\textbf{échange marchand} moderne, fondé sur trois caractères :

\begin{itemize}
\item le \textbf{calcul} : chaque bien a un prix, on compare et on mesure ;
\item l'\textbf{équivalence} : on donne une valeur jugée égale à ce qu'on reçoit ;
\item l'\textbf{anonymat} relatif : une fois le paiement fait, la relation peut s'arrêter là — le vendeur et l'acheteur ne se doivent plus rien.
\end{itemize}

\begin{exemplebox}[Le marché hebdomadaire]
Le marché hebdomadaire d'un village camerounais illustre bien l'échange marchand : la commerçante y vend ses légumes au prix discuté, le client paie en francs CFA, et chacun repart « quitte ». Pourtant, même là, l'échange n'est jamais purement économique : on salue, on marchande en plaisantant, on fait crédit à la cliente fidèle. Le marché reste un lieu de \emph{sociabilité}.
\end{exemplebox}

\courssub{Échange marchand et don : deux logiques}

Il faut donc distinguer soigneusement deux logiques de l'échange :

\begin{center}
\begin{tabularx}{\linewidth}{|l|X|X|}\hline
\rowcolor{popblueL} \textbf{Critère} & \textbf{Échange marchand} & \textbf{Don (échange symbolique)} \\ \hline
Objet & Biens, services, argent & Cadeaux, paroles, honneurs \\ \hline
Règle & Équivalence calculée, prix & Réciprocité sans calcul précis \\ \hline
Temps & Immédiat, quitte à quitte & Différé (on rendra un jour) \\ \hline
Lien créé & Faible, relation terminée après paiement & Fort, dette symbolique durable \\ \hline
Exemple & Achat au marché & Dot, cadeau de funérailles, tontine \\ \hline
\end{tabularx}
\end{center}

La \textbf{dot matrimoniale} camerounaise en est la parfaite illustration : les biens remis par la famille du fiancé à celle de la fiancée ne « paient » pas la femme — ce serait une erreur grave de lecture — mais scellent une \textbf{alliance} entre deux familles. De même, le « cadeau » offert lors des funérailles ne s'achète pas : il honore le défunt et soutient la famille, en attendant que celle-ci « rende » lors d'une cérémonie future.

\begin{exoresolu}[Analyse d'une situation concrète]
\textbf{Énoncé.} À Bafoussam, lors du mariage de Sandrine, la famille du fiancé remet à celle de la mariée des enveloppes d'argent, des sacs de riz, des pagnes et des boissons. Un observateur étranger déclare : « Cette famille a acheté une femme. » Cette interprétation est-elle correcte ? Justifie ta réponse en distinguant échange marchand et échange symbolique.
\tcblower
\textbf{Solution détaillée.} L'interprétation de l'observateur est \textbf{incorrecte}, car elle confond deux logiques d'échange. Dans un \emph{échange marchand}, on donne une valeur équivalente pour acquérir un bien, et la relation s'achève avec le paiement : on achète un objet, pas une personne. Or la dot obéit à la logique du \emph{don} décrite par Mauss : elle relève de la triple obligation — la famille du fiancé \textbf{donne}, celle de la mariée \textbf{reçoit} (refuser serait rompre l'alliance), et elle \textbf{rendra} lors des cérémonies futures ou par l'accueil du gendre. La dot ne fixe pas un prix à la femme ; elle \textbf{scelle une alliance durable} entre deux familles et crée un réseau de dettes symboliques réciproques. Conclusion : réduire la dot à un achat, c'est appliquer la catégorie marchande là où règne la catégorie symbolique — une erreur d'analyse que le philosophe doit savoir éviter.
\end{exoresolu}

\begin{aretenir}[L'essentiel]
L'échange est un transfert réciproque de biens, de paroles ou de symboles. Selon Mauss, tout don obéit à la triple obligation : \textbf{donner, recevoir, rendre}. On distingue trois formes d'échanges : \textbf{économiques} (troc, monnaie, marché), \textbf{linguistiques} (la parole comme premier échange) et \textbf{symboliques} (don, dot, cadeaux). L'échange marchand repose sur le calcul et l'équivalence ; le don repose sur la réciprocité différée et tisse le lien social.
\end{aretenir}

\coursec{Les échanges, fondement de la société}

Dans la section précédente, nous avons défini l'échange comme un transfert réciproque de biens, de services ou de paroles entre des partenaires qui se reconnaissent mutuellement. Mais cette définition ne dit pas encore \emph{pourquoi} l'échange occupe une place si centrale dans la vie humaine. Pourquoi ne pouvons-nous pas vivre chacun de notre côté, en autarcie ? La réponse que nous allons construire est forte : l'échange n'est pas un simple complément de la vie sociale, il en est le \cle{fondement}. C'est parce que les hommes échangent qu'il existe une société. Pour le démontrer, nous partirons d'une observation simple — nul ne ne peut produire seul tout ce dont il a besoin — puis nous montrerons comment l'échange engendre l'interdépendance, le lien social et la paix, avant d'en examiner les limites.

\courssub{L'homme, être de besoins : l'origine de la cité selon Platon}

Commençons par une expérience de pensée à la portée de chacun. Imaginez un élève de Première technique qui déciderait, demain matin, de ne plus rien recevoir de personne : il devrait produire lui-même sa nourriture, tisser ses vêtements, construire sa maison, fabriquer ses outils, soigner ses maladies, et même s'enseigner à lui-même ce qu'il ignore. Il est évident qu'il n'y survivrait pas, ou alors dans une misère extrême. Cette impossibilité n'est pas un accident : elle révèle une vérité sur la nature humaine.

\begin{definition}[L'homme, être de besoins]
L'homme est un \cle{être de besoins} : il a besoin de nourriture, de vêtements, d'abri, de soins, d'éducation, de sécurité. Or aucun individu isolé ne peut produire seul, par son seul travail, la totalité de ce qui satisfait ces besoins. L'insuffisance de l'individu est donc le point de départ de toute vie collective.
\end{definition}

C'est précisément l'intuition que Platon développe dans \emph{La République} (livre II). Socrate y demande : d'où naît une cité ? Sa réponse est célèbre : la cité naît de nos \emph{besoins}, car « chacun de nous ne se suffit pas à lui-même, mais a besoin de beaucoup de choses ». Le paysan produit du blé mais ne fabrique pas ses outils ; le forgeron fabrique des outils mais ne cultive pas son pain ; le tisserand habille les autres mais doit être nourri. Chacun donne donc ce qu'il produit en trop et reçoit ce qui lui manque : la cité est, à l'origine, un \cle{marché de services réciproques}.

\begin{propriete}[Thèse de Platon : la cité naît de l'échange]
Dans \emph{La République}, Platon montre que la cité (la société organisée) naît de l'\cle{échange des services}, parce que chacun, ne se suffisant pas à lui-même, a besoin des autres.
\par\textbf{Justification.} Le raisonnement de Platon se décompose ainsi :
\begin{enumerate}
\item Tout homme a des besoins multiples (nourriture, logement, vêtement\dots).
\item Aucun homme ne peut, à lui seul, produire tout ce qui satisfait ces besoins (manque de temps, de force, de savoir-faire).
\item De plus, chacun a une aptitude naturelle particulière : on travaille mieux et plus quand on fait une seule chose.
\item Il est donc avantageux et nécessaire que chacun se consacre à une tâche et échange le produit de son travail contre celui des autres.
\item Or échanger suppose de vivre ensemble, en un même lieu, avec des règles communes.
\item Donc la vie en société (la cité) est la conséquence nécessaire de l'échange : sans échange, pas de cité.
\end{enumerate}
\end{propriete}

\begin{exemplebox}[Le quartier comme petite cité]
Observez un quartier de Yaoundé ou de Douala : le maquisard prépare le poisson braisé mais achète le poisson au pêcheur, le bois au vendeur de charbon, les épices au commerçant ; le taximan transporte les clients mais fait réparer sa voiture par le mécanicien, qui lui-même envoie ses enfants chez l'enseignant. Personne ne ne se suffit à soi-même : le quartier entier est un tissu d'échanges, exactement comme la cité de Platon.
\end{exemplebox}

\courssub{La division sociale du travail : chacun dépend de tous}

De l'analyse de Platon découle une notion essentielle : celle de \cle{division sociale du travail}.

\begin{definition}[La division sociale du travail]
La \cle{division sociale du travail} est la répartition des tâches productives entre les membres d'une société (cultivateur, commerçant, enseignant, médecin, transporteur\dots), répartition qui les rend \cle{mutuellement dépendants} les uns des autres : chacun se spécialise dans une fonction et compte sur les autres pour tout le reste.
\end{definition}

L'intuition est simple : en se spécialisant, chacun devient plus habile et plus productif dans son domaine — mais aussi plus dépendant. Le médecin sait soigner, mais il ne sait ni cultiver son manioc ni réparer son moto-taxi ; le cultivateur nourrit la ville, mais il a besoin du médecin quand il tombe malade. La spécialisation crée donc un réseau d'\cle{interdépendance} : je dépends du travail des autres, et les autres dépendent du mien.

\begin{popfigure}
\begin{center}
\begin{tikzpicture}[scale=0.95, every node/.style={font=\small}]
% Noeuds
\node[draw, fill=popgreenL, rounded corners, minimum width=2.2cm, minimum height=0.9cm] (paysan) at (0,3) {Paysan};
\node[draw, fill=poporangeL, rounded corners, minimum width=2.2cm, minimum height=0.9cm] (commercant) at (4.2,3) {Commerçant};
\node[draw, fill=popblueL, rounded corners, minimum width=2.2cm, minimum height=0.9cm] (medecin) at (0,0) {Médecin};
\node[draw, fill=poppurpleL, rounded corners, minimum width=2.2cm, minimum height=0.9cm] (enseignant) at (4.2,0) {Enseignant};
\node[draw, fill=popgoldL, rounded corners, minimum width=2.4cm, minimum height=0.9cm] (transporteur) at (2.1,1.5) {Transporteur};
% Flèches d'échanges
\draw[-{Stealth}, thick, popdark] (paysan) -- (commercant);
\draw[-{Stealth}, thick, popdark] (commercant) -- (enseignant);
\draw[-{Stealth}, thick, popdark] (enseignant) -- (medecin);
\draw[-{Stealth}, thick, popdark] (medecin) -- (paysan);
\draw[-{Stealth}, thick, popdark] (paysan) -- (transporteur);
\draw[-{Stealth}, thick, popdark] (transporteur) -- (commercant);
\draw[-{Stealth}, thick, popdark] (transporteur) -- (medecin);
\draw[-{Stealth}, thick, popdark] (transporteur) -- (enseignant);
\end{tikzpicture}
\end{center}
\legende{Le réseau d'interdépendance créé par la division sociale du travail. Chaque métier (nœud) échange avec tous les autres (flèches) : le paysan vend ses vivres au commerçant, le transporteur achemine les marchandises et les malades, le médecin soigne l'enseignant qui instruit les enfants de tous. Supprimez un seul nœud, et tout le réseau souffre : c'est la preuve concrète que chacun dépend de tous.}
\end{popfigure}

Ce réseau a une conséquence philosophique décisive : l'échange n'est pas seulement une opération économique, il est un \cle{lien}. Quand j'achète le pain du boulanger, je ne fais pas que recevoir du pain : j'entretiens une relation avec lui, je contribue à sa subsistance comme il contribue à la mienne. La société tout entière apparaît alors comme un vaste système de services réciproques.

\courssub{L'échange crée le lien social}

Allons plus loin : pourquoi l'échange crée-t-il du \emph{lien social}, et pas seulement de la dépendance matérielle ? La réponse tient à la structure même de l'échange. Pour échanger avec quelqu'un, je dois :
\begin{itemize}
\item le reconnaître comme un \cle{partenaire}, c'est-à-dire comme quelqu'un avec qui je traite d'égal à égal, et non comme une chose que je peux simplement prendre ou détruire ;
\item le reconnaître comme un \cle{semblable}, c'est-à-dire comme un être qui, comme moi, a des besoins, des droits et une parole ;
\item accepter une \cle{règle commune} : je donne parce que je recevrai, je tiens ma parole parce que l'autre tient la sienne.
\end{itemize}

Échanger, c'est donc renoncer à la violence pure. Celui que je pourrais voler, je le paie ; celui que je pourrais combattre, je le salue et je négocie avec lui. L'échange transforme l'étranger en partenaire : il institue entre les hommes une relation de \cle{reconnaissance réciproque}, qui est la forme la plus élémentaire du lien social.

\begin{exemplebox}[Les funérailles villageoises au Cameroun]
Lors des funérailles dans les villages camerounais, les parents, amis et voisins apportent des contributions : argent, vivres, boissons, aide à l'organisation. Ces dons ne sont jamais à sens unique : celui qui a contribué chez vous recevra votre contribution demain. Cet échange matériel \emph{exprime} la solidarité (je partage ta peine) et la \emph{renforce} (nous sommes liés par des services réciproques). Le lien familial et villageois se tisse littéralement par l'échange.
\end{exemplebox}

\courssub{Lévi-Strauss : l'échange à l'origine de la société}

L'anthropologue Claude \cle{Lévi-Strauss}, dans \emph{Les Structures élémentaires de la parenté} (1949), a donné à cette idée une portée radicale. Il s'est demandé : pourquoi toutes les sociétés humaines, sans exception, interdisent-elles l'inceste, c'est-à-dire le mariage entre proches parents ? Sa réponse : la \cle{prohibition de l'inceste} n'est pas d'abord une règle biologique, c'est une règle \emph{sociale} qui oblige les groupes à échanger.

Le raisonnement est le suivant. Si un groupe garde « ses » femmes pour lui-même, il reste fermé, isolé, et risque la guerre avec ses voisins. En interdisant le mariage à l'intérieur du groupe, la prohibition de l'inceste oblige chaque groupe à chercher des épouses \emph{chez les autres} — et à donner les siennes en retour. Les femmes circulent entre les groupes comme des biens précieux, et cette circulation crée des alliances : deux familles unies par un mariage deviennent des alliées, non des ennemies. La prohibition de l'inceste est ainsi, selon Lévi-Strauss, le \cle{premier échange}, celui qui fonde toutes les autres formes de vie sociale : « l'échange est l'origine même de la société ».

\begin{aretenir}[L'essentiel sur Lévi-Strauss]
Pour Lévi-Strauss, la prohibition de l'inceste est la première règle sociale : en obligeant les groupes à échanger des épouses, elle transforme des familles isolées en une société d'alliés. L'échange n'est donc pas une activité parmi d'autres : il est \emph{constitutif} de la société humaine.
\end{aretenir}

\begin{popfigure}
\begin{center}
\begin{tikzpicture}[scale=0.9, every node/.style={font=\small, align=center}]
\node[draw, fill=popblueL, rounded corners, minimum width=2.2cm, minimum height=1cm, align=center] (besoins) at (0,0){Besoins\\(nul ne se\\suffit à soi)};
\node[draw, fill=popgreenL, rounded corners, minimum width=2.2cm, minimum height=1cm, align=center] (echange) at (3.2,0){Échange\\de biens,\\de services};
\node[draw, fill=poporangeL, rounded corners, minimum width=2.2cm, minimum height=1cm, align=center] (interdep) at (6.4,0){Interdépendance\\(division sociale\\du travail)};
\node[draw, fill=poppurpleL, rounded corners, minimum width=2.2cm, minimum height=1cm, align=center] (lien) at (9.6,0){Lien social\\(reconnaissance\\réciproque)};
\node[draw, fill=popgoldL, rounded corners, minimum width=2.2cm, minimum height=1cm, align=center] (societe) at (12.8,0){Société\\(vie commune\\organisée)};
\draw[-{Stealth}, very thick, popdark] (besoins) -- (echange);
\draw[-{Stealth}, very thick, popdark] (echange) -- (interdep);
\draw[-{Stealth}, very thick, popdark] (interdep) -- (lien);
\draw[-{Stealth}, very thick, popdark] (lien) -- (societe);
\end{tikzpicture}
\end{center}
\legende{De l'échange à la société : la chaîne du fondement. Parce que l'homme a des besoins qu'il ne peut satisfaire seul, il échange ; l'échange spécialise les tâches et rend chacun dépendant de tous ; cette dépendance, vécue comme reconnaissance mutuelle, tisse le lien social ; et l'ensemble de ces liens constitue la société.}
\end{popfigure}

\courssub{Les échanges et la paix : « le doux commerce »}

Si l'échange fonde la société à l'intérieur d'un groupe, il peut aussi pacifier les relations \emph{entre} les groupes et les peuples. C'est la thèse du « \cle{doux commerce} », défendue notamment par Montesquieu dans \emph{De l'Esprit des lois} (1748) : « l'effet naturel du commerce est de porter à la paix ». Pourquoi ? Parce que deux peuples qui commercent ensemble ont un \emph{intérêt commun} à préserver leur relation : faire la guerre à son client ou à son fournisseur, c'est détruire sa propre richesse. L'intérêt marchand remplace la passion de conquérir.

\begin{savaistu}[Montesquieu et le commerce qui adoucit les mœurs]
Montesquieu pensait que le commerce « adoucit les mœurs » : les peuples qui commercent ensemble apprennent à se connaître, à négocier, à tenir leurs engagements, et ont moins tendance à se faire la guerre. À l'inverse, observe-t-il, là où le commerce est absent, on se livre plus facilement au brigandage et à la conquête. Mieux vaut vendre à son voisin que le piller !
\end{savaistu}

On peut en donner une illustration camerounaise : les grands marchés (marché Mokolo à Yaoundé, marché Mboppi à Douala, marché de Foumban) réunissent chaque jour des commerçants et des clients de toutes les régions et de toutes les communautés. Autour de l'étal, les différences d'origine s'effacent derrière l'intérêt partagé de la transaction. De même, la coopération intercommunautaire — mariages entre villages, associations de développement, tontines, travaux collectifs — crée entre les groupes des liens durables qui préviennent les conflits.

\courssub{Limites : les échanges peuvent-ils tout régler ?}

Il serait cependant naïf de conclure que l'échange est toujours harmonieux et bienfaisant. Une analyse rigoureuse doit en marquer les limites.

\begin{itemize}
\item \textbf{Les échanges peuvent être inégaux.} Quand un partenaire est beaucoup plus puissant que l'autre, l'échange cesse d'être libre : le petit producteur de cacao, par exemple, vend sa récolte à un prix qu'il ne choisit pas, fixé par des acheteurs plus forts que lui. L'échange formellement libre peut cacher un rapport de force.
\item \textbf{L'échange peut devenir exploitation.} Si l'un profite du besoin de l'autre pour lui imposer des conditions injustes (salaire de misère, travail des enfants), l'échange ne fonde plus la reconnaissance mutuelle mais la domination.
\item \textbf{Tout ne doit pas être marchandisé.} Certaines réalités humaines — la dignité, la justice, l'amour, le corps — ne sont pas des marchandises : les mettre en vente (corruption, trafic d'êtres humains) détruit le lien social au lieu de le fonder.
\item \textbf{L'échange marchand ne suffit pas.} Une société a aussi besoin de dons gratuits (solidarité, entraide) et de services publics (école, santé, sécurité) qui échappent à la logique du marché.
\end{itemize}

\begin{attention}[Erreur fréquente]
Ne jamais affirmer que les échanges sont \emph{toujours} harmonieux et pacifiques. L'échange peut être inégal, conflictuel ou exploiteur. Dans une dissertation, la bonne démarche consiste à montrer d'abord que l'échange fonde le lien social (thèse), puis à en discuter les conditions : il ne fonde la société que s'il est libre, équitable et limité à ce qui peut légitimement s'échanger.
\end{attention}

\courssub{Méthode : montrer qu'une notion « fonde » une autre}

\begin{methode}[Prouver un rapport de fondement]
Pour montrer qu'une notion A (l'échange) « fonde » une notion B (la société), il faut établir que B est \emph{impossible} sans A. Procédez en quatre étapes :
\begin{enumerate}
\item \textbf{Définir} clairement les deux notions (qu'est-ce qu'échanger ? qu'est-ce qu'une société ?).
\item \textbf{Formuler l'hypothèse contraire} (raisonnement par l'absurde) : une société sans aucun échange est-elle concevable ? Imaginons des individus qui ne se donnent, ne se vendent et ne se disent jamais rien : ce ne serait pas une société, mais une juxtaposition de solitudes.
\item \textbf{Dégager le mécanisme} : montrer par quelles étapes A produit B (besoins $\rightarrow$ échange $\rightarrow$ interdépendance $\rightarrow$ lien social $\rightarrow$ société).
\item \textbf{Nuancer la conclusion} : A fonde B, mais sous certaines conditions (échange équitable, non-marchandisation de l'humain), ce qui montre l'esprit critique.
\end{enumerate}
\end{methode}

\begin{exoresolu}[Peut-on concevoir une société sans échanges ?]
\textbf{Énoncé.} Un élève affirme : « On pourrait imaginer une société où chacun produirait tout ce dont il a besoin, sans rien échanger avec personne. » En vous appuyant sur Platon, montrez que cette affirmation est intenable, puis rédigez un paragraphe argumenté (10 à 15 lignes).
\tcblower
\textbf{Solution détaillée.} D'après Platon, l'homme est un être de besoins qui ne se suffit pas à lui-même : produire seul sa nourriture, ses vêtements, sa maison, ses outils et ses soins dépasse les forces et le temps de tout individu. L'hypothèse de l'élève est donc matériellement impossible ; et même si elle était réalisable, les individus ainsi isolés ne formeraient pas une \emph{société}, puisqu'aucun lien ne les unirait.
\par\textbf{Paragraphe rédigé.} L'affirmation selon laquelle chacun pourrait vivre sans rien échanger se heurte à la nature même de l'homme. Comme le montre Platon dans \emph{La République}, l'homme est un être de besoins : il lui faut manger, se vêtir, se loger, se soigner. Or aucun individu ne possède à lui seul le temps, la force et les savoir-faire nécessaires pour produire tout cela : le temps passé à cultiver manquerait pour tisser, et le temps passé à construire manquerait pour cultiver. Il est donc nécessaire que chacun se spécialise et échange le fruit de son travail contre celui des autres : c'est la division sociale du travail. Mais l'échange fait plus que combler les besoins : en me rendant dépendant des autres et les autres dépendants de moi, il crée entre nous un lien de reconnaissance réciproque. Une prétendue société sans échanges serait donc doublement impossible : matériellement, parce que nul ne se suffit à soi ; socialement, parce que sans échange il n'y aurait aucun lien, donc pas de société du tout. L'échange apparaît ainsi comme le véritable fondement de la vie commune.
\end{exoresolu}

\begin{exercice}[Application au Cameroun]
Dans votre village ou votre quartier, identifiez trois pratiques d'échange (marché, tontine, travaux collectifs, contributions lors des cérémonies). Pour chacune, montrez en quelques lignes : (1) quel besoin elle satisfait ; (2) quel lien social elle crée ou renforce ; (3) si elle peut comporter une limite (inégalité, obligation excessive).
\end{exercice}

\begin{corrige}[Exercice 1]
Exemple de réponse attendue. \textbf{La tontine} : (1) elle satisfait le besoin d'argent pour un projet (commerce, frais scolaires) ; (2) elle crée entre les membres un lien de confiance et de solidarité durable, fondé sur la parole donnée ; (3) limite possible : la pression sur celui qui ne peut pas cotiser. \textbf{Le marché hebdomadaire} : (1) il permet à chacun d'acheter ce qu'il ne produit pas ; (2) il fait se rencontrer des communautés différentes, qui apprennent à se connaître et à négocier pacifiquement ; (3) limite : les prix peuvent être imposés par les plus forts. \textbf{Les travaux collectifs} (construction d'une case, entretien d'une piste) : (1) ils réalisent ce qu'une famille seule ne pourrait faire ; (2) ils expriment et renforcent la solidarité villageoise ; (3) limite : l'obligation de participer peut peser sur certains. L'essentiel est de montrer que l'échange y apparaît à la fois comme réponse aux besoins et comme tissu du lien social.
\end{corrige}

\begin{aretenir}[Ce qu'il faut retenir de la section]
\begin{itemize}
\item L'homme est un \cle{être de besoins} : nul ne peut produire seul tout ce dont il a besoin (Platon : la cité naît de l'échange des services).
\item La \cle{division sociale du travail} spécialise chacun et rend tous les membres de la société mutuellement dépendants.
\item L'échange crée le \cle{lien social} : échanger, c'est reconnaître l'autre comme partenaire et comme semblable.
\item Pour \cle{Lévi-Strauss}, la prohibition de l'inceste est le premier échange : elle fonde les alliances entre groupes, donc la société.
\item Le « \cle{doux commerce} » de Montesquieu : les échanges pacifient les relations entre les peuples.
\item Mais l'échange a des \cle{limites} : inégalités, exploitation, marchandisation de l'humain. Il ne fonde la société que s'il reste libre et équitable.
\end{itemize}
\end{aretenir}

\coursec{Méthode : appliquer les notions et rédiger une argumentation}

Après avoir analysé les échanges comme fondement de la société, nous devons maintenant maîtriser les outils intellectuels qui permettent de mobiliser ces connaissances face à un sujet d'examen ou à une situation concrète. La philosophie en classe de Première technique ne se limite pas à l'acquisition de savoirs : elle exige une \cle{pratique argumentée}, c'est-à-dire la capacité à problématiser, conceptualiser et argumenter avec rigueur. Cette section vous donne la méthode pas à pas, illustrée par des cas camerounais, pour réussir la dissertation et l'explication de texte au Probatoire.

\courssub{Appliquer une notion à une situation concrète : la démarche en quatre temps}

Face à un cas pratique (un fait de société, un article de presse, une scène de la vie quotidienne), l'élève ne doit pas « donner son avis » ni « raconter sa vie ». Il doit \cle{conceptualiser} : faire passer le singulier sous l'universel. La démarche s'opère en quatre temps rigoureux.

\begin{methode}[Appliquer une notion à une situation concrète]
\begin{enumerate}
    \item \textbf{Identifier les faits bruts} : décrire la situation sans la juger (qui ? quoi ? où ? quand ? comment ?).
    \item \textbf{Dégager la notion philosophique pertinente} : quelle notion du programme (travail, échange, aliénation, don, valeur d'usage/valeur d'échange, lien social, etc.) éclaire ce fait ?
    \item \textbf{Mobiliser un auteur ou une thèse} : citer explicitement un philosophe étudié (Hegel, Marx, Mauss, Locke, Rousseau, Arendt, etc.) dont la pensée permet d'analyser la situation.
    \item \textbf{Formuler un jugement justifié} : conclure en répondant au problème implicite ou explicite de la situation, en nuancant si nécessaire.
\end{enumerate}
\end{methode}

\begin{popfigure}
\begin{tikzpicture}[
    node distance=1cm and 1.2cm,
    box/.style={draw=popblue, thick, rounded corners, fill=popblueL, text width=2.6cm, align=center, minimum height=1.4cm, font=\small},
    arrow/.style={-Stealth, thick, popdark}
]
\node[box, align=center] (sit){Situation concrète\\ (faits observés)};
\node[box, right=of sit, align=center] (prob){Problème\\ philosophique};
\node[box, right=of prob, align=center] (th){Thèse / Notion\\ mobilisée};
\node[box, right=of th, align=center] (aut){Auteur / Référence\\ (Hegel, Marx, Mauss...)};
\node[box, right=of aut, align=center] (syn){Jugement\\ justifié / Synthèse};

\draw[arrow] (sit) -- (prob) node[midway, above, font=\small] {Questionner};
\draw[arrow] (prob) -- (th) node[midway, above, font=\small] {Conceptualiser};
\draw[arrow] (th) -- (aut) node[midway, above, font=\small] {Fonder};
\draw[arrow] (aut) -- (syn) node[midway, above, font=\small] {Argumenter};
\end{tikzpicture}
\legende{Organigramme de la démarche argumentative : de la situation concrète au jugement philosophique justifié.}
\end{popfigure}

\courssub{Étude de cas guidée 1 : le jeune diplômé mototaximan — travail dévalorisé ou dignité par l'effort ?}

Appliquons la méthode à une situation emblématique du Cameroun contemporain.

\begin{exemplebox}[Situation]
Paul, 26 ans, titulaire d'un BTS en Gestion des Entreprises, n'a pas trouvé d'emploi de bureau après deux ans de recherche. Pour subvenir à ses besoins et aider sa famille, il devient mototaximan (\emph{benskin}) à Douala. Il travaille 12h par jour, gagne environ 15\,000 FCFA nets par jour, mais subit le mépris de certains clients, l'insécurité, l'absence de protection sociale. Ses anciens camarades de promotion qui ont eu des concours le traitent d'« échec ».
\end{exemplebox}

\begin{exoresolu}[Analyse philosophique guidée]
Application de la méthode en 4 temps au cas de Paul.
\tcblower
\textbf{1. Faits bruts :} Diplôme technique + chômage structurel + reconversion dans l'informel (moto-taxi) + revenu décent mais précarité + stigmatisation sociale.

\textbf{2. Notions mobilisées :} \cle{Travail} (activité humaine transformant la nature), \cle{Aliénation} (perte de soi dans l'activité), \cle{Dignité} (valeur intrinsèque de la personne), \cle{Valeur d'usage / Valeur d'échange} (Marx), \cle{Reconnaissance} (Hegel/Honneth).

\textbf{3. Auteurs / Thèses :}
\begin{itemize}
    \item \textbf{Hegel (\emph{Phénoménologie de l'Esprit}) :} Le travail est l'acte par lequel l'homme se \cle{réalise} en objectivant sa conscience dans la matière. Même le travail manuel forme l'autonomie. \emph{Thèse :} Paul se réalise par l'effort ; il domine la machine (la moto) et la route.
    \item \textbf{Marx (\emph{Manuscrits de 1844}) :} Sous le capitalisme, le travail est \cle{aliéné} : l'ouvrier ne s'appartient pas dans son travail, il y perd sa vie. Paul vend sa force de travail à la journée, ne contrôle pas les conditions (horaires, sécurité), produit une valeur d'échange pour survivre, pas pour s'épanouir. \emph{Antithèse :} Paul est aliéné ; son diplôme est dévalorisé, il subit la contrainte économique.
    \item \textbf{Hannah Arendt (\emph{Condition de l'homme moderne}) :} Distinction \cle{Travail} (animal laborans, reproduction de la vie biologique, cyclique, nécessité) vs \cle{Œuvre} (homo faber, fabrication d'un monde durable). Le benskin est du pur « travail » (survie), pas de l'« œuvre » (construction d'un monde).
\end{itemize}

\textbf{4. Jugement justifié (Synthèse) :}
La situation de Paul révèle la \cle{tension entre la valeur anthropologique du travail (Hegel : dignité par l'effort, autonomie) et sa valeur sociale marchande (Marx : aliénation, dévalorisation du diplôme)}. Paul \emph{garde sa dignité humaine} par l'effort consenti et la responsabilité familiale (thèse hégélienne), \emph{mais} il subit une \emph{injustice structurelle} : la société ne lui offre pas les conditions d'un travail libre et reconnu (critique marxiste). Le mépris des camarades confond \emph{statut social} (emploi salarié protégé) et \emph{valeur morale} (courage, responsabilité). \textbf{Conclusion :} Le travail de Paul n'est pas « dévalorisé » en soi (il produit un service utile, valeur d'usage), mais il est \emph{socialement dévalorisé} par l'absence de droits et de reconnaissance. La philosophie invite à distinguer la \cle{dignité du travailleur} (inaliénable) de la \cle{condition du travail} (historiquement perfectible).
\end{exoresolu}

\begin{attention}[Piège à l'examen : Raconter sa vie au lieu de conceptualiser]
\emph{Erreur fréquente :} « Moi aussi mon oncle est mototaximan et il est fier, donc le travail libère. » \rightarrow \textbf{Zéro point.} L'exemple personnel n'est valable que s'il \textbf{illustre une idée générale} (thèse d'un auteur, notion du cours). \emph{Bonne rédaction :} « Comme le montre l'exemple de Paul, mototaximan diplômé, le travail manuel conserve une valeur anthropologique (Hegel : appropriation du monde), mais il devient aliénant lorsque les conditions sociales (précarité, absence de reconnaissance) empêchent le sujet de se reconnaître dans son œuvre (Marx). »
\end{attention}

\courssub{Étude de cas guidée 2 : la tontine comme échange économique et lien social}

\begin{exemplebox}[Situation]
Dans un quartier de Yaoundé, dix femmes commerçantes organisent une tontine mensuelle (\emph{njangi}). Chaque mois, chacune verse 20\,000 FCFA. Le pot (200\,000 FCFA) est attribué à tour de rôle à l'une d'elles, souvent pour un investissement (stock, frais scolaires, terrain). Pas de contrat écrit, pas d'intérêt, pas de banque. La confiance, la parole donnée et la solidarité du groupe garantissent le système. Si l'une défaut, le groupe la soutient ou la sanctionne par l'exclusion.
\end{exemplebox}

\begin{exoresolu}[Analyse philosophique guidée]
Application de la méthode en 4 temps au cas de la tontine.
\tcblower
\textbf{1. Faits :} Groupe fermé, cotisation fixe, redistribution tournante, absence d'intérêt financier, régulation par la confiance/parole, finalité d'investissement/entraide.

\textbf{2. Notions :} \cle{Échange} (don/contre-don), \cle{Don} (Mauss), \cle{Lien social}, \cle{Réciprocité}, \cle{Valeur d'usage} vs \cle{Valeur d'échange}, \cle{Contrat} vs \cle{Pacte}.

\textbf{3. Auteurs / Thèses :}
\begin{itemize}
    \item \textbf{Marcel Mauss (\emph{Essai sur le don}) :} La tontine est un \cle{fait social total} : économique (circulation d'argent), juridique (obligation de rendre), moral (honneur, confiance), religieux/rituel (réunions, repas). Le don n'est jamais « gratuit » : il crée une \cle{dette symbolique} qui lie les personnes. \emph{Thèse :} La tontine fonde la société par l'obligation triplement : donner, recevoir, rendre.
    \item \textbf{John Locke / Contractualistes :} Contraste avec le \cle{contrat marchand} (échange instantané, équivalence valeur/valeur, anonymat, droit). La tontine repose sur un \cle{pacte de confiance} (durée, identité, asymétrie temporelle : je donne maintenant, je recevrai plus tard).
    \item \textbf{Karl Polanyi (\emph{La Grande Transformation}) :} L'économie \cle{enchâssée} dans le social. La tontine n'est pas un marché autorégulé ; elle est encastrée dans les relations de parenté, de voisinage, de genre.
\end{itemize}

\textbf{4. Jugement justifié :}
La tontine n'est pas un « échange primitif » ou « informel » au sens péjoratif. C'est une \cle{forme sophistiquée d'échange non-marchand} qui produit du \cle{lien social} là où le marché (banque, microfinance) échoue (exclusion, taux d'intérêt, garanties foncières). Elle réalise l'idéal maussien : \textbf{l'échange comme ciment de la société}. Cependant, elle a des limites : montant plafonné, risque de défaillance, exclusion des plus pauvres (incapables de cotiser). \textbf{Synthèse :} La tontine démontre que \emph{l'échange précède et fonde le contrat} : sans confiance première (don), pas de marché possible. Elle est un \emph{rempart contre la désaffiliation} (Castel) dans l'économie informelle camerounaise.
\end{exoresolu}

\begin{popfigure}
\begin{tikzpicture}
\begin{axis}[
    width=\linewidth,
    height=7cm,
    axis lines=middle,
    xlabel={Temps (mois)},
    ylabel={Capital disponible (FCFA)},
    xtick={1,2,3,4,5,6,7,8,9,10},
    ytick={0,50000,100000,150000,200000},
    yticklabels={0, 50k, 100k, 150k, 200k},
    ymin=0, ymax=220000,
    xmin=0.5, xmax=10.5,
    grid=major, grid style={popline},
    title style={font=\small, text=popdark},
    title={Modélisation simplifiée d'une tontine à 10 membres (cotisation 20k/mois)}
]
% Capital cumulé du groupe (ligne montante)
\addplot[popblue, thick, mark=*] coordinates {
    (1,20000) (2,40000) (3,60000) (4,80000) (5,100000) (6,120000) (7,140000) (8,160000) (9,180000) (10,200000)
};
\addlegendentry{Capital total cotisé (flux entrant)}

% Flèche annotée pour le "pot"
\node[poporange, font=\small, align=center] at (axis cs:5,210000) {Le "Pot" = 200\,000 FCFA\\ attribué chaque mois\\ à une membre différente};
\draw[poporange, -Stealth, thick] (axis cs:5,205000) -- (axis cs:5,180000);
\end{axis}
\end{tikzpicture}
\legende{Visualisation du flux financier dans une tontine : l'échange n'est pas instantané (comme au marché), il est \textbf{différé} et \textbf{collectif}, créant un lien de dette/confiance durable (Mauss).}
\end{popfigure}

\courssub{Rédiger une démarche argumentative : la dissertation philosophique}

Au Probatoire, le sujet de dissertation (ex. : « Le travail libère-t-il l'homme ? », « L'échange est-il la base de la société ? ») exige une \cle{structure dialectique} : Thèse — Antithèse — Synthèse. Ce n'est pas un plan « pour / contre » binaire, mais une \cle{montée en profondeur} du problème.

\begin{methode}[Structure type de la dissertation (Probatoire / Bac Tech)]
\textbf{INTRODUCTION (environ 15--20 lignes)}
\begin{enumerate}
    \item \textbf{Accroche (Amorce) :} Une citation d'auteur, un fait d'actualité camerounaise, une définition étymologique, une opposition courante. \emph{Exemple :} « Au Cameroun, on dit : "Le travail nourrit son homme." Pourtant, Paul le mototaximan diplômé se demande s'il est vraiment libre. »
    \item \textbf{Définition des termes / Analyse du sujet :} Définir les concepts clés (Travail, Libérer, Homme). Montrer la tension : libération \emph{par} le travail (maîtrise) vs libération \emph{du} travail (loisirs) vs aliénation \emph{dans} le travail.
    \item \textbf{Problématique :} La question ouverte qui guide tout le devoir. \emph{Exemple :} « \textbf{Dans quelle mesure le travail, tout en étant condition de l'autonomie humaine, peut-il devenir source d'aliénation selon les conditions historiques de sa réalisation ?} »
    \item \textbf{Annonce du plan :} « Nous verrons d'abord que le travail est essence de l'humanité (Thèse : Hegel), puis qu'il peut être aliénation sous le capitalisme (Antithèse : Marx), pour conclure que la libération dépend de la transformation sociale des conditions de travail (Synthèse : Marx / Arendt / Droit du travail). »
\end{enumerate}

\textbf{DÉVELOPPEMENT (3 grandes parties, 2 à 3 sous-parties chacune)}
\begin{itemize}
    \item \textbf{Partie I (Thèse) :} Le travail libère / L'échange fonde la société. Arguments + Auteurs + Exemples concrets (Cameroun).
    \item \textbf{Partie II (Antithèse) :} Le travail aliène / L'échange marchand détruit le lien. Arguments + Auteurs + Exemples.
    \item \textbf{Partie III (Synthèse / Dépassement) :} Condition de possibilité d'un travail libre / d'un échange juste. Rôle de la technique, du droit, de la politique, de la reconnaissance. \textbf{Pas de simple résumé !} Nouvelle thèse qui intègre les deux précédentes.
\end{itemize}

\textbf{CONCLUSION (environ 10 lignes)}
\begin{enumerate}
    \item Reprise synthétique du parcours (réponse à la problématique).
    \item Ouverture : prolongement vers une autre notion du programme (ex. : Technique, État, Justice, Bonheur) ou une question actuelle (ex. : « L'intelligence artificielle libérera-t-elle enfin l'homme du travail pénible ? »).
\end{enumerate}
\end{methode}

\begin{methode}[Une bonne problématique : question ouverte, pas fermée]
\emph{Mauvaise problématique (fermée) :} « Le travail libère-t-il l'homme ? » (Réponse : Oui / Non). \rightarrow \textbf{Plan binaire, note faible.}
\emph{Bonne problématique (ouverte) :} « \textbf{Comment le travail, qui est la condition de l'autonomie humaine, peut-il devenir, sous certaines conditions historiques, le lieu de l'aliénation ?} » \rightarrow \textbf{Plan dialectique, note forte.}
\textbf{Astuce :} Commencez la problématique par « Comment », « Dans quelle mesure », « Pourquoi », « Quelles sont les conditions pour que... ».
\end{methode}

\courssub{Justifier une réponse : le schéma AEI (Affirmation -- Explication -- Illustration)}

En philosophie, \textbf{aucune affirmation ne vaut sans justification}. Au Probatoire, on attend : \cle{Affirmer} $\rightarrow$ \cle{Expliquer} $\rightarrow$ \cle{Illustrer}.

\begin{methode}[Schéma AEI pour chaque argument du développement]
\begin{enumerate}
    \item \textbf{A -- Affirmation (L'argument) :} Une phrase claire qui répond à la problématique. \emph{Ex :} « Le travail permet à l'homme de se reconnaître comme sujet libre. »
    \item \textbf{E -- Explication (Le concept / L'auteur) :} Déployer le sens philosophique. Citer l'auteur. \emph{Ex :} « Selon \textbf{Hegel}, dans la dialectique du Maître et de l'Esclave, c'est par le \textbf{travail} (négation de la nature) que la conscience servile accède à l'indépendance et à la vérité de soi. L'esclave transforme le monde extérieur et s'y reconnaît ; il devient "pour-soi". »
    \item \textbf{I -- Illustration (L'exemple concret / Contre-exemple) :} Ancrer dans le réel (Cameroun si possible). \emph{Ex :} « C'est le cas de l'artisan menuisier à Foumban qui, en sculptant le bois selon son projet, voit son intelligence s'objectiver dans l'œuvre ; il se sent "maître" de sa matière, contrairement à l'ouvrier à la chaîne qui répète un geste imposé. »
\end{enumerate}
\textbf{Règle d'or :} Un paragraphe = Un argument = Un schéma AEI complet.
\end{methode}

\begin{exemplebox}[Sujet type : « Le travail libère-t-il l'homme ? » — Plan détaillé AEI]
\textbf{Thèse (Hegel) :} Le travail libère par l'objectivation.
\begin{itemize}
    \item[A] Le travail est l'acte fondateur de la conscience de soi.
    \item[E] Hegel, \emph{Phénoménologie de l'Esprit} : l'esclave travaille, transforme la nature, se reconnaît dans l'objet produit $\rightarrow$ autonomie.
    \item[I] L'agriculteur du Nord-Cameroun qui maîtrise son champ, ses semences, son calendrier : il est libre \emph{dans} son travail.
\end{itemize}
\textbf{Antithèse (Marx) :} Le travail salarié capitaliste aliène.
\begin{itemize}
    \item[A] Le travail devient aliénation quand l'ouvrier ne possède ni le produit, ni le processus, ni lui-même.
    \item[E] Marx, \emph{Manuscrits de 1844} : 4 formes d'aliénation (produit, acte, espèce, homme). Travail forcé, non libre, mortification.
    \item[I] Paul le mototaximan ou l'ouvrier de la SOCAPALM / CDC : cadences, précarité, absence de propriété sur l'outil, travail \emph{pour} survivre, pas \emph{par} projet.
\end{itemize}
\textbf{Synthèse :} La libération n'est pas dans le travail \emph{en soi}, mais dans les \textbf{conditions sociales du travail}.
\begin{itemize}
    \item[A] Travail libre = travail associé, non aliéné, reconnu (Droit du travail, Coopératives, Autogestion).
    \item[E] Marx (\emph{Capital}, \emph{Critique du Programme de Gotha}) : « Dans une société communiste, le travail devient le premier besoin vital. » Arendt : distinction Travail (nécessité) / Œuvre (monde) / Action (politique). La technique doit libérer du \emph{travail} (peine) pour ouvrir à l'\emph{action} et à l'\emph{œuvre}.
    \item[I] Les coopératives paysannes (GIC) au Cameroun : mutualisation des moyens, décision collective, valorisation du produit $\rightarrow$ reprise de pouvoir sur le travail.
\end{itemize}
\end{exemplebox}

\courssub{Entraînement à l'explication de texte : Marx et Mauss}

L'explication de texte suit une méthode rigoureuse : \textbf{Comprendre $\rightarrow$ Analyser $\rightarrow$ Discuter}.

\begin{methode}[Méthode de l'explication de texte (3 étapes)]
\textbf{Étape 1 : Comprendre (Le "Quoi ?")}
\begin{itemize}
    \item Identifier l'auteur, l'œuvre, la date, le contexte.
    \item Repérer la \textbf{thèse principale} (une phrase).
    \item Découper le texte en \textbf{unités de sens} (paragraphes logiques, étapes du raisonnement).
    \item Définir les \textbf{concepts clés} du texte (aliénation, force de travail, don, réciprocité, fait social total...).
\end{itemize}

\textbf{Étape 2 : Analyser / Expliquer (Le "Comment ?")} -- \emph{Cœur du devoir}
\begin{itemize}
    \item Suivre le fil du texte \textbf{pas à pas} (linéaire) ou par \textbf{thèmes} (thématique).
    \item Pour chaque étape : \textbf{Reformuler} l'idée de l'auteur $\rightarrow$ \textbf{Expliquer} le mécanisme conceptuel $\rightarrow$ \textbf{Montrer} comment l'argument soutient la thèse.
    \item \textbf{Citer} le texte (guillemets, numéro de ligne) pour appuyer l'explication.
\end{itemize}

\textbf{Étape 3 : Discuter / Évaluer (Le "Combien ? / Et alors ?")}
\begin{itemize}
    \item \textbf{Portée :} Pourquoi ce texte est-il important ? (Rupture théorique, fondement d'une discipline).
    \item \textbf{Limites / Critiques :} Qu'est-ce que le texte ne dit pas ? Quelles objections ? (Ex. : Marx sous-estime le rôle de la technique comme libération ; Mauss essentialise le don ?).
    \item \textbf{Actualisation :} Lien avec l'actualité camerounaise (tontine, travail informel, économie collaborative).
\end{itemize}
\end{methode}

\begin{exemplebox}[Extrait type : Marx, \emph{Manuscrits de 1844} -- « L'objet que le travail produit... »]
\textbf{Thèse :} Le travail aliéné fait de l'ouvrier un esclave de son propre produit.
\textbf{Analyse guidée :}
\begin{enumerate}
    \item \textbf{Aliénation dans le produit :} « L'objet que le travail produit... lui est étranger, hostile, puissant. » $\rightarrow$ Explication : L'ouvrier met sa vie dans l'objet, mais l'objet ne lui appartient pas (capitaliste). Plus il produit, plus il s'appauvrit (valeur d'échange > valeur d'usage pour lui).
    \item \textbf{Aliénation dans l'acte :} « Le travail est extérieur à l'ouvrier... il ne s'affirme pas, il se nie. » $\rightarrow$ Explication : Le travail n'est pas volontaire, c'est une contrainte (salariat). Il ne déploie pas ses forces, il les use. « Mortification ».
    \item \textbf{Conséquence anthropologique :} « L'homme devient animal. » $\rightarrow$ Explication : Perte de l'humanité générique (créativité, conscience, liberté). Réduction à des fonctions biologiques (manger, boire, loger).
\end{enumerate}
\textbf{Discussion :} Pertinence pour analyser le travail précaire au Cameroun (CDD, sous-traitance, benskin). Limite : Marx néglige les formes de résistance ouvrière et la possible fierté du métier (sociologie du travail : Dubar, Dejours).
\end{exemplebox}

\begin{exemplebox}[Extrait type : Mauss, \emph{Essai sur le don} -- « Le don... oblige »]
\textbf{Thèse :} Le don n'est pas gratuit ; il crée un lien social obligatoire (donner, recevoir, rendre).
\textbf{Analyse guidée :}
\begin{enumerate}
    \item \textbf{Le don "libre" est une illusion :} « Il n'y a pas de don gratuit. » $\rightarrow$ Explication : Même sans contrat juridique, le don engage l'honneur, l'âme (\emph{hau}, esprit de la chose donnée chez les Maoris). Refuser de donner ou de rendre = guerre ou perte de statut.
    \item \textbf{Triple obligation :} Donner (générosité/puissance), Recevoir (soumission/confiance), Rendre (libération de la dette).
    \item \textbf{Fait social total :} Économique, juridique, moral, religieux, politique à la fois. La tontine camerounaise en est l'illustration parfaite.
\end{enumerate}
\textbf{Discussion :} Fondement de l'anthropologie économique. Critique : Peut-on tout réduire au don ? Le marché moderne (monnaie, contrat) a permis l'universalisation des échanges au-delà du groupe restreint. Polanyi : le "désenchâssement" de l'économie est historique, pas seulement moral.
\end{exemplebox}

\begin{savaistu}[Atout décisif pour le Bac Tech]
Au Baccalauréat technique, les sujets de philosophie portent \textbf{très fréquemment sur l'application des notions à la vie sociale, économique et technique} (ex. : « Technique et emploi », « Échange et développement », « Travail et dignité »). Les correcteurs valorisent les copies qui :
\begin{itemize}
    \item Mobilisent des \textbf{exemples camerounais précis} (tontine, GIC, benskin, agriculture, industries locales, startups, droit OHADA, Code du travail).
    \item Distinguent clairement \textbf{notion philosophique} et \textbf{réalité empirique}.
    \item Utilisent le \textbf{vocabulaire technique du cours} (aliénation, valeur d'usage, réciprocité, fétichisme, reconnaissance, etc.).
\end{itemize}
S'entraîner sur ces cas concrets \emph{tout au long de l'année} (et pas seulement en révision) est la meilleure préparation à l'épreuve.
\end{savaistu}

\begin{exercice}[Entraînement : Appliquer la méthode AEI]
\textbf{Sujet :} « L'échange marchand détruit-il le lien social ? »
\begin{enumerate}
    \item Rédigez une \textbf{problématique ouverte} (style « Dans quelle mesure... » / « Comment... »).
    \item Proposez un \textbf{plan dialectique} (Thèse / Antithèse / Synthèse) avec les auteurs principaux.
    \item Rédigez \textbf{un paragraphe complet (AEI)} pour la Thèse (L'échange marchand crée du lien / pacifie les relations) en citant \textbf{Montesquieu} (\emph{Esprit des lois}, « doux commerce ») et en illustrant par un exemple camerounais (ex. : marché transfrontalier, plateforme e-commerce).
    \item Rédigez \textbf{un paragraphe complet (AEI)} pour l'Antithèse (L'échange marchand dégrade le lien / tout devient marchandise) en citant \textbf{Marx} (fétichisme de la marchandise) ou \textbf{Polanyi} et en illustrant (ex. : privatisation de l'eau, travail des enfants dans les mines, rupture du lien communautaire).
\end{enumerate}
\end{exercice}

\begin{corrige}[Exercice n°1 -- Éléments de correction]
\textbf{1. Problématique :} « \textbf{Dans quelle mesure l'échange marchand, en universalisant les relations humaines, risque-t-il de les vider de leur substance morale et communautaire ?} »
\textbf{2. Plan :}
\begin{itemize}
    \item \textbf{I. Thèse (Montesquieu, Smith) :} Le "doux commerce" pacifie les mœurs, crée l'interdépendance, fondement de la société civile moderne (contrat, confiance abstraite). Ex : Marché central de Douala, Zone de libre-échange continentale (ZLECAf).
    \item \textbf{II. Antithèse (Marx, Polanyi, Sandel) :} Le fétichisme de la marchandise masque les rapports sociaux ; la "grande transformation" désenchâsse l'économie, tout devient marchandise (corps, nature, éducation), érosion de la solidarité. Ex : Accaparement des terres, uberisation, dettes étudiantes.
    \item \textbf{III. Synthèse (Sen, Droit, Économie sociale) :} Le marché a besoin de \textbf{non-marchand} pour fonctionner (confiance, institutions, don, État de droit). L'échange juste = marché \textbf{encadré} + sphères protégées (service public, communs, tontine, coopérative).
\end{itemize}
\textbf{3. Paragraphe Thèse (AEI) :}
[A] L'échange marchand, loin de détruire le lien, est le ciment de la paix civile. [E] Pour \textbf{Montesquieu} (\emph{De l'esprit des lois}, Livre XX), le commerce « poli les mœurs » : là où il y a commerce, les hommes apprennent à se connaître, à négocier plutôt qu'à se faire la guerre ; l'intérêt propre pousse à respecter l'autrui comme partenaire. [I] Au Cameroun, les marchés transfrontaliers (ex. : Kousseri/N'Djaména, Garoua Boulai/Bertoua) obligent des populations de cultures différentes à coopérer quotidiennement ; la monnaie (FCFA) et le contrat verbal remplacent la violence ; les plateformes comme Jumia ou Glovo créent de nouveaux liens de confiance notée entre vendeurs et acheteurs anonymes.
\textbf{4. Paragraphe Antithèse (AEI) :}
[A] Mais l'extension illimitée de la forme marchande corrompt les liens les plus intimes. [E] \textbf{Marx} montre que le \textbf{fétichisme de la marchandise} fait apparaître les rapports sociaux entre hommes comme des rapports entre choses (prix) ; \textbf{Polanyi} (\emph{La Grande Transformation}) dénonce les "fictitious commodities" (travail, terre, monnaie) : traiter le travail comme une marchandise détruit la communauté villageoise. [I] Dans les zones minières de l'Est (Bétaré-Oya), l'orpaillage artisanal transforme la terre ancestrale en concession vendue, brise l'autorité des chefs, pousse des enfants à la carrière au lieu de l'école : le lien clanique se brise sous la pression de la valeur d'échange immédiate.
\end{corrige}

\begin{aretenir}[Check-list finale pour l'épreuve de Philosophie (Probatoire)]
\begin{itemize}
    \item \textbf{Sujet disserté :} Problématique ouverte $\rightarrow$ Plan dialectique (3 parties) $\rightarrow$ AEI dans chaque sous-partie $\rightarrow$ Auteurs cités $\rightarrow$ Exemples camerounais $\rightarrow$ Conclusion + Ouverture.
    \item \textbf{Explication de texte :} Thèse + Découpage $\rightarrow$ Explication linéaire rigoureuse (citations) $\rightarrow$ Discussion critique + Actualisation.
    \item \textbf{Cas pratique / Question de cours :} Méthode 4 temps (Faits $\rightarrow$ Notion $\rightarrow$ Auteur $\rightarrow$ Jugement) + Schéma AEI.
    \item \textbf{Vocabulaire :} Maîtriser les définitions précises (Travail, Aliénation, Échange, Don, Valeur d'usage/échange, Fétichisme, Lien social, Reconnaissance).
    \item \textbf{Attitude :} Pas d'opinion non argumentée, pas de récit de vie brut, pas de hors-sujet. \textbf{Conceptualisez !}
\end{itemize}
\end{aretenir}

\newpage

\coursec{Activité d'intégration}

\coursec{Le travail et les échanges}

\courssub{Objectifs de l'activité}

Cette activité d'intégration a pour but de faire travailler ensemble les deux savoirs de la leçon --- le \cle{travail} (définition et valeur) et les \cle{échanges} (fondements de la société) --- dans une situation concrète de la vie courante camerounaise. À l'issue de l'activité, l'élève doit être capable de :

\begin{itemize}
\item repérer, dans des documents réels (photographie, tableau statistique, témoignage), les différentes \cle{formes de travail} et les différentes \cle{formes d'échanges} ;
\item mobiliser les thèses des auteurs étudiés (\cle{Platon}, \cle{Aristote}, \cle{Marcel Mauss}, \cle{Karl Marx}, \cle{Hegel}) pour analyser ces situations ;
\item conduire un débat argumenté sur la valeur du travail (libération ou aliénation) ;
\item rédiger collectivement une conclusion argumentée d'une quinzaine de lignes selon la démarche \cle{thèse -- antithèse -- synthèse}, puis la présenter oralement à la classe.
\end{itemize}

\courssub{Situation}

Dans le cadre de la semaine philosophique du lycée technique de la Briqueterie (Yaoundé), le professeur de philosophie organise une \textbf{enquête de quartier} intitulée : \emph{« Le travail et les échanges dans mon quartier »}. La classe de Première technique est divisée en groupes de cinq élèves. Chaque groupe reçoit une chemise contenant trois documents.

\textbf{Document 1 --- Photographie.} Une photographie du marché Mokolo (Yaoundé) un samedi matin : étals de légumes, poissons fumés, tissus ; vendeuses assises derrière leurs marchandises ; mototaximen déposant des clients ; porteurs transportant des sacs ; acheteurs marchandant les prix.

\textbf{Document 2 --- Tableau statistique.} Résultats d'un mini-recensement réalisé par les élèves auprès de 200 chefs de ménage du quartier (enquête fictive mais réaliste) :

\begin{center}
\begin{tabularx}{\linewidth}{|l|X|c|c|}\hline
\rowcolor{popblueL}
\textbf{Métier principal} & \textbf{Nature de l'activité} & \textbf{Effectif} & \textbf{Part} \\ \hline
Cultivateurs (maraîchage) & Production de biens & 40 & 20 \% \\ \hline
Commerçants et vendeuses & Échange marchand & 60 & 30 \% \\ \hline
Mototaximen & Service de transport & 30 & 15 \% \\ \hline
Couturiers, coiffeurs, menuisiers & Artisanat (transformation) & 40 & 20 \% \\ \hline
Enseignants, infirmiers, fonctionnaires & Services publics & 30 & 15 \% \\ \hline
\textbf{Total} & & \textbf{200} & \textbf{100 \%} \\ \hline
\end{tabularx}
\end{center}

\textbf{Document 3 --- Témoignage.} Maman Élise, 45 ans, vendeuse de légumes au marché Mokolo, déclare : \emph{« Je vends ici depuis vingt ans. Mon travail est dur : je me lève à 4 heures du matin. Mais grâce à lui, j'ai scolarisé mes quatre enfants. Je fais partie d'une tontine de douze vendeuses : chaque mois, chacune verse \num{10000} francs CFA, et l'une de nous reçoit les \num{120000} francs à tour de rôle. C'est comme ça que j'ai acheté ma congélateur. Ici au marché, nous sommes comme une famille : quand l'une est malade, les autres gardent sa place. »}

\textbf{Problème philosophique posé à la classe :} le travail et les échanges observés dans ce quartier ne sont-ils que des moyens de survivre, ou sont-ils aussi ce qui fait de ce quartier une véritable \cle{société} ?

\courssub{Consignes}

\begin{enumerate}
\item \textbf{Analyser les documents.} Pour chacun des trois documents, identifiez : (a) les formes de travail observées (travail productif, travail de service, travail artisanal) ; (b) les formes d'échanges observées (échange marchand, échange de services, don et contre-don dans la tontine, entraide entre vendeuses).
\item \textbf{Appliquer les notions.} À l'aide du tableau du document 2, montrez que le quartier illustre la thèse de \cle{Platon} dans \emph{La République} : la cité naît de la division du travail, parce que nul ne peut se suffire à lui-même. Calculez la part des activités d'échange et de services (commerçants, mototaximen, fonctionnaires) et commentez le résultat.
\item \textbf{Interpréter la tontine.} En vous appuyant sur la théorie du \cle{don} de Marcel Mauss (\emph{Essai sur le don}, 1925) --- triple obligation de donner, recevoir et rendre ---, montrez que la tontine de Maman Élise n'est pas un simple échange d'argent, mais un lien social.
\item \textbf{Débattre.} Le travail de la vendeuse au marché est-il \cle{libérateur} ou \cle{aliénant} ? Chaque groupe formule d'abord deux arguments pour chaque position (on pourra penser à Hegel : le travail forme et libère l'homme ; et à Marx : le travail peut aliéner quand le travailleur est exploité et épuisé).
\item \textbf{Rédiger.} Rédigez collectivement une conclusion argumentée d'environ quinze lignes répondant au problème posé, selon la démarche \cle{thèse -- antithèse -- synthèse} : thèse (le travail et les échanges ne sont que des moyens de survie), antithèse (ils fondent le lien social), synthèse (ils sont les deux à la fois : en échangeant, les hommes satisfont leurs besoins et créent la société).
\item \textbf{Présenter.} Un rapporteur par groupe présente la conclusion à la classe en cinq minutes ; les autres groupes posent des questions.
\end{enumerate}

\courssub{Ressources à mobiliser}

\begin{aretenir}[Notions et repères utiles]
\begin{itemize}
\item \textbf{Le travail} : activité humaine, manuelle ou intellectuelle, par laquelle l'homme transforme la nature pour satisfaire ses besoins. Il a une valeur \cle{vitale} (survie), une valeur \cle{économique} (revenu) et une valeur \cle{sociale et morale} (dignité, reconnaissance, lien social).
\item \textbf{Platon} : la cité naît de la \cle{division du travail} ; chacun fait ce pour quoi il est doué, et l'échange des produits unit les citoyens.
\item \textbf{Aristote} : distingue l'\emph{économie} (échange naturel pour satisfaire des besoins) de la \emph{chrématistique} (recherche illimitée du gain) ; l'échange juste repose sur la réciprocité et la proportionnalité, fondement de la justice sociale.
\item \textbf{Hegel} : par le travail, l'homme transforme le monde et se transforme lui-même ; le travail est \cle{libérateur}.
\item \textbf{Marx} : dans certaines conditions, le travail devient \cle{aliénation} : le travailleur est épuisé, ne se reconnaît plus dans ce qu'il produit.
\item \textbf{Mauss} : le \cle{don} obéit à trois obligations --- donner, recevoir, rendre --- et crée des liens durables ; l'échange n'est pas seulement économique, il est un « fait social total ».
\item \textbf{Méthode de rédaction} : thèse -- antithèse -- synthèse ; chaque affirmation doit être justifiée par un argument et illustrée par un exemple tiré des documents.
\end{itemize}
\end{aretenir}

\courssub{Démarche et corrigé commenté}

\begin{exoresolu}[Corrigé commenté de l'enquête philosophique]
\textbf{Énoncé.} Analyser les trois documents, mobiliser Platon et Mauss, débattre de la valeur du travail de la vendeuse, puis rédiger une conclusion argumentée (thèse -- antithèse -- synthèse).
\tcblower
\textbf{Étape 1 --- Analyse des documents.}

Document 1 (photographie) : on observe des \textbf{travaux} variés (vendeuses, porteurs, mototaximen) et des \textbf{échanges marchands} (marchandage, achat contre argent). Document 2 (tableau) : on distingue le travail \emph{productif} (cultivateurs, artisans : 40 \%), le travail d'\emph{échange} (commerçants : 30 \%) et le travail de \emph{service} (mototaximen, fonctionnaires : 30 \%). Document 3 (témoignage) : le travail de Maman Élise a une valeur vitale (nourrir la famille), économique (scolarisation des enfants) et sociale (la tontine, l'entraide).

\textbf{Étape 2 --- Application de Platon.}

Part des activités d'échange et de services : $30\,\% + 15\,\% + 15\,\% = 60\,\%$. Plus de la moitié des actifs du quartier ne produisent pas directement de biens : ils vivent de l'échange et des services. Cela confirme la thèse de Platon : nul ne se suffit à lui-même ; le cultivateur a besoin du mototaximen pour transporter sa récolte, le commerçant a besoin du cultivateur, tous ont besoin de l'enseignant. La \cle{division du travail} rend chacun dépendant de tous, et cette dépendance réciproque \emph{fonde la société}.

\textbf{Étape 3 --- Application de Mauss.}

La tontine illustre la triple obligation du don : chaque vendeuse \emph{donne} (\num{10000} francs par mois), \emph{reçoit} (\num{120000} francs à son tour) et s'engage à \emph{rendre} en continuant de verser après avoir reçu. L'argent circule, mais ce qui circule surtout, c'est la \cle{confiance} et la \cle{solidarité} : « quand l'une est malade, les autres gardent sa place ». L'échange est donc un lien social, pas seulement un calcul économique.

\textbf{Étape 4 --- Débat : libération ou aliénation ?}

\emph{Arguments pour la libération} : le travail de Maman Élise lui donne une autonomie économique (elle scolarise ses enfants), une dignité et une reconnaissance sociale (elle est intégrée dans un réseau de solidarité). Selon Hegel, en travaillant, elle se réalise. \emph{Arguments pour l'aliénation} : conditions dures (lever à 4 heures), fatigue, faibles revenus, absence de protection sociale ; selon Marx, un tel travail peut épuiser l'homme sans lui laisser le temps de s'épanouir.

\textbf{Étape 5 --- Exemple de conclusion rédigée (quinze lignes).}

\emph{Thèse} : on pourrait penser que le travail et les échanges du quartier ne sont que des moyens de survie : Maman Élise vend pour nourrir sa famille, le mototaximen transporte pour gagner sa journée. \emph{Antithèse} : mais l'enquête montre davantage. Avec 60 \,\% d'actifs vivant de l'échange et des services, le quartier vérifie la thèse de Platon : la division du travail rend les habitants nécessaires les uns aux autres. Et la tontine, selon l'analyse de Mauss, crée par le don et le contre-don des liens de confiance et de solidarité qui dépassent le simple calcul d'argent. \emph{Synthèse} : le travail et les échanges sont donc les deux à la fois : en cherchant à satisfaire leurs besoins, les hommes tissent sans toujours le savoir le lien social. Certes, le travail peut devenir aliénant quand il épuise ; mais il reste, quand il est reconnu, un chemin de dignité et de liberté. Le quartier n'est pas une simple addition d'individus : c'est une société, parce qu'on y travaille et qu'on y échange.
\end{exoresolu}

\courssub{Bilan de l'activité}

Cette enquête a permis de mettre en évidence trois résultats essentiels. D'abord, les notions philosophiques ne sont pas abstraites : le \cle{travail} et l'\cle{échange} s'observent concrètement au marché Mokolo comme dans tout quartier camerounais. Ensuite, les thèses des auteurs éclairent la réalité : Platon explique pourquoi la diversité des métiers unit le quartier, Aristote précise la nature de l'échange juste, Mauss révèle la dimension sociale de la tontine, Hegel et Marx permettent de nuancer la valeur du travail. Enfin, l'activité a entraîné les deux savoir-faire de l'unité : \emph{appliquer les notions à des situations concrètes} et \emph{rédiger une démarche argumentée et justifiée}, compétence directement mobilisable dans l'épreuve de dissertation au Probatoire et au Baccalauréat technique.

\begin{attention}[Pièges à éviter dans la rédaction]
Ne pas confondre échange marchand (achat-vente) et don (tontine, entraide) ; ne pas citer un auteur sans expliquer sa thèse ; ne pas opposer libération et aliénation de façon absolue : un même travail peut être l'un et l'autre selon les conditions ; ne pas oublier de justifier chaque affirmation par un exemple tiré des documents.
\end{attention}

\newpage

\coursec{Méthodes et exercices résolus}

\courssub{Savoir-faire 1 : Appliquer les notions de travail et d'échange à une situation concrète}

\begin{methode}[Analyser une situation concrète avec les notions de la leçon]
\begin{enumerate}
\item \textbf{Lire attentivement la situation} et repérer les éléments pertinents : qui agit ? sur quoi ? avec quels outils ? dans quel but ? contre quelle contrepartie ?
\item \textbf{Identifier la notion mobilisée} : s'agit-il de travail (activité humaine transformant la nature, avec outils et finalité) ou d'échange (transfert réciproque de biens, de services ou de paroles) ?
\item \textbf{Appliquer la définition du cours} : vérifier que la situation remplit les critères de la définition (pour le travail : intentionnalité, transformation, outillage ; pour l'échange : réciprocité, utilité, règles).
\item \textbf{Qualifier la situation} avec le vocabulaire précis du cours : travail manuel ou intellectuel, travail libérateur ou aliéné (Marx), échange marchand (acheter/vendre), échange symbolique (don/contre-don, Mauss), échange de paroles.
\item \textbf{Conclure en justifiant} : dire ce que la situation montre (valeur du travail, fonction sociale de l'échange) en une phrase claire.
\end{enumerate}
\textbf{Réflexes} : toujours citer la définition avant de l'appliquer ; toujours nommer l'auteur de la thèse mobilisée (Marx, Mauss, Locke) ; toujours rattacher l'exemple à la problématique.
\end{methode}

\begin{exoresolu}[Le tailleur de pierre de Foumban]
\textbf{Énoncé.} À Foumban, un artisan tailleur de pierre façonne des sculptures qu'il vend aux touristes sur le marché artisanal. Il dit : « Mon travail me fatigue, mais il me fait vivre et il me rend fier. » Appliquez les notions de la leçon pour analyser cette situation : en quoi cette activité est-elle un travail, et quelles valeurs du travail y apparaissent ?

\tcblower
\textbf{Solution détaillée.}

\textbf{Étape 1 : repérage des éléments.} Un homme (l'artisan) agit sur une matière (la pierre) avec des outils (ciseaux, marteau), dans un but précis (produire des sculptures vendables), contre une contrepartie (l'argent des touristes).

\textbf{Étape 2 : identification de la notion.} Il s'agit d'un \cle{travail} au sens propre : une activité humaine, consciente et finalisée, qui transforme la nature (la pierre brute) à l'aide d'outils pour produire un objet utile.

\textbf{Étape 3 : application de la définition.} Les trois critères sont remplis : (a) \emph{intentionnalité} — l'artisan a un projet, la sculpture existe d'abord dans sa pensée ; (b) \emph{transformation} — la pierre brute devient œuvre ; (c) \emph{outillage} — les instruments prolongent la main. Comme le souligne le cours, le travail distingue l'homme de l'animal : l'abeille construit par instinct, l'artisan construit d'après un plan.

\textbf{Étape 4 : qualification.} C'est un \cle{travail manuel} artisanal. Les paroles de l'artisan révèlent les trois \cle{valeurs du travail} : la \emph{valeur économique} (« il me fait vivre » : le travail produit une richesse échangeable) ; la \emph{valeur morale et personnelle} (« il me rend fier » : le travail réalise et dignifie l'homme) ; et la \emph{valeur pénible} (« il me fatigue » : le travail est aussi effort, ce qui rappelle l'étymologie latine \emph{tripalium}, instrument de torture).

\textbf{Conclusion.} Cette situation montre que le travail n'est pas seulement une contrainte : il est à la fois source de revenu, moyen d'épanouissement et effort pénible. L'artisan de Foumban illustre la thèse selon laquelle le travail humanise l'homme en le rendant créateur de son monde.
\end{exoresolu}

\courssub{Savoir-faire 2 : Distinguer travail libérateur et travail aliéné}

\begin{methode}[Mobiliser la dialectique libération/aliénation]
\begin{enumerate}
\item \textbf{Rappeler les deux thèses} : le travail libère (il arrache l'homme à la nature, développe ses capacités, lui donne une dignité — thèse hégélienne) ; le travail aliène (l'ouvrier est dépossédé du produit, du geste et de lui-même — thèse de Marx dans les \emph{Manuscrits de 1844}).
\item \textbf{Repérer dans l'énoncé les indices d'aliénation} : travail répétitif à la chaîne, produit qui n'appartient pas au travailleur, salaire de simple survie, absence de créativité.
\item \textbf{Repérer les indices de libération} : autonomie, création, reconnaissance sociale, maîtrise du produit.
\item \textbf{Nuancez toujours} : une même activité peut être libératrice sous un aspect et aliénante sous un autre ; la conclusion doit être problématisée, pas manichéenne.
\end{enumerate}
\textbf{Formule à retenir} : « Le travail est à la fois malédiction (peine), libération (humanisation) et risque d'aliénation (dépossession) : tout dépend des conditions dans lesquelles il s'exerce. »
\end{methode}

\begin{exoresolu}[L'ouvrière de la zone industrielle de Douala]
\textbf{Énoncé.} Dans une usine de la zone industrielle de Douala, une ouvrière répète huit heures par jour le même geste : coudre des étiquettes sur des chemises destinées à l'exportation. Elle ne verra jamais les chemises finies et son salaire couvre à peine ses besoins. Montrez, à l'aide de Marx, que ce travail est aliéné, puis indiquez ce qui le distinguerait d'un travail libérateur.

\tcblower
\textbf{Solution détaillée.}

\textbf{Étape 1 : rappel de la thèse.} Pour Marx, dans le \emph{travail salarié capitaliste}, l'ouvrier subit une quadruple \cle{aliénation} : il est séparé du \emph{produit} de son travail (qui appartient au patron), de l'\emph{acte} de production (geste imposé, répétitif), de sa \emph{nature humaine} (le travail ne réalise plus ses capacités créatrices) et des \emph{autres hommes} (rapports de concurrence et d'exploitation).

\textbf{Étape 2 : application point par point.}
\begin{itemize}
\item \emph{Aliénation du produit} : « elle ne verra jamais les chemises finies » — le produit lui est étranger, il part à l'exportation.
\item \emph{Aliénation de l'acte} : « répète huit heures par jour le même geste » — le travail est morcelé, sans initiative ni créativité.
\item \emph{Aliénation de soi} : le salaire « couvre à peine ses besoins » — elle travaille pour survivre, non pour se réaliser ; sa force de travail est devenue une marchandise.
\end{itemize}

\textbf{Étape 3 : contraste avec le travail libérateur.} Ce travail serait libérateur si l'ouvrière maîtrisait l'ensemble du processus (voir la chemise aboutir), si elle pouvait exercer son jugement et sa créativité, et si le produit et la juste valeur de son effort lui revenaient. Le travail libère lorsqu'il est \emph{conscient, créatif et reconnu} : c'est le cas de l'artisan qui conçoit et achève son œuvre.

\textbf{Conclusion.} L'exemple de Douala vérifie l'analyse marxienne : le travail n'est pas libérateur par nature, il l'est seulement sous certaines conditions économiques et sociales. La même activité peut donc être humanisante ou aliénante selon le rapport du travailleur à son produit et à son geste.
\end{exoresolu}

\courssub{Savoir-faire 3 : Analyser un échange et ses fondements}

\begin{methode}[Distinguer les formes d'échange et montrer leur rôle social]
\begin{enumerate}
\item \textbf{Identifier le type d'échange} : échange \cle{marchand} (achat-vente, équivalence monétaire), échange \cle{symbolique} (don et contre-don, selon Marcel Mauss dans \emph{Essai sur le don}), échange de \cle{paroles} (communication, langage).
\item \textbf{Dégager la réciprocité} : tout échange engage celui qui donne et celui qui reçoit ; Mauss montre que le don crée une triple obligation : donner, recevoir, rendre.
\item \textbf{Montrer la fonction sociale} : l'échange crée du lien, institue la confiance, fonde la société (on ne vit pas seul : la division du travail rend les hommes mutuellement dépendants).
\item \textbf{Conclure} : l'échange n'est pas seulement économique, il est le \emph{fondement du lien social}.
\end{enumerate}
\textbf{Réflexe} : dans un échange marchand, demander « quelle est l'équivalence ? » ; dans un don, demander « quelle est l'obligation créée ? ».
\end{methode}

\begin{exoresolu}[Le don lors d'un mariage traditionnel à Bafoussam]
\textbf{Énoncé.} Lors d'un mariage traditionnel à Bafoussam, la famille du marié remet à la famille de la mariée des cadeaux (dot) ; en retour, la famille de la mariée offre un banquet et des présents. Un observateur dit : « Ce n'est pas du commerce, c'est de l'amitié. » À l'aide de Mauss, analysez cet échange et montrez en quoi il fonde la société.

\tcblower
\textbf{Solution détaillée.}

\textbf{Étape 1 : identification du type d'échange.} Il ne s'agit pas d'un échange marchand (aucun prix n'est fixé, aucune monnaie n'équivalant exactement aux biens n'est exigée), mais d'un échange \cle{symbolique} de type \emph{don/contre-don}, étudié par Marcel Mauss dans l'\emph{Essai sur le don} (1925).

\textbf{Étape 2 : analyse de la réciprocité.} Mauss dégage trois obligations : \emph{donner} (la famille du marié offre la dot), \emph{recevoir} (refuser serait rejeter l'alliance), \emph{rendre} (la famille de la mariée répond par le banquet et les présents). Le don n'est donc jamais gratuit au sens d'unilatéral : il engage celui qui reçoit et crée un lien durable entre les deux familles.

\textbf{Étape 3 : fonction sociale.} L'observateur a raison de dire que ce n'est pas du commerce, mais il a tort d'y voir seulement de l'amitié spontanée : c'est une \emph{institution sociale réglée}. Par l'échange des dons, deux lignages s'allient, la paix et la confiance s'établissent, la communauté se renforce. L'échange précède et fonde la société : comme le montre le cours, les hommes, dépendants les uns des autres par la division du travail, ne peuvent vivre ensemble qu'en échangeant biens, services et paroles.

\textbf{Conclusion.} Le mariage de Bafoussam illustre la thèse de Mauss : l'échange, même non marchand, est le ciment du lien social. Échanger, c'est reconnaître l'autre et instituer la vie en communauté.
\end{exoresolu}

\courssub{Savoir-faire 4 : Rédiger une argumentation de type dissertation}

\begin{methode}[Construire une argumentation philosophique en cinq étapes]
\begin{enumerate}
\item \textbf{Analyser le sujet} : souligner les mots-clés (« travail », « libère », « aliène »), définir chacun, dégager la \cle{problématique} sous forme de question (ex. : « Le travail est-il seulement une contrainte ou aussi un moyen de libération ? »).
\item \textbf{Formuler une thèse} claire et personnelle, annoncée dès l'introduction.
\item \textbf{Construire le plan} en deux ou trois parties : thèse adverse (le travail comme malédiction), thèse défendue (le travail comme libération), dépassement (conditions d'un travail libérateur, risque d'aliénation).
\item \textbf{Argumenter chaque partie} selon le schéma : idée $\rightarrow$ définition $\rightarrow$ argument $\rightarrow$ exemple concret (Cameroun de préférence) $\rightarrow$ référence d'auteur $\rightarrow$ mini-conclusion.
\item \textbf{Conclure} en répondant nettement à la problématique et en ouvrant éventuellement sur une question voisine.
\end{enumerate}
\textbf{Formules à retenir} : « On pourrait objecter que... mais... » ; « Cet exemple montre que... » ; « Ainsi, ... ». Une idée = un paragraphe ; jamais de catalogue d'auteurs sans argument.
\end{methode}

\begin{exoresolu}[Sujet de dissertation : « Le travail libère-t-il l'homme ? »]
\textbf{Énoncé.} Rédigez l'introduction et la première partie (thèse adverse) d'une dissertation sur le sujet : \emph{Le travail libère-t-il l'homme ?}

\tcblower
\textbf{Solution détaillée.}

\textbf{Introduction (accroche, définitions, problématique, annonce du plan).}

\emph{Accroche} : Chaque matin, des millions de Camerounais — planteurs de cacao, moto-taxi, enseignants — se mettent au travail, souvent avec peine. \emph{Définition} : le travail est l'activité humaine consciente qui transforme la nature à l'aide d'outils pour produire des biens utiles ; « libérer » signifie arracher à une servitude. \emph{Problème} : le travail apparaît d'abord comme une fatigue et une obligation ; pourtant on dit aussi qu'il dignifie et humanise. \emph{Problématique} : le travail est-il une malédiction subie, ou le moyen par lequel l'homme se libère de la nature et de lui-même ? \emph{Annonce du plan} : nous verrons d'abord le travail comme malédiction, puis comme libération, enfin les conditions qui font basculer l'un dans l'autre.

\textbf{Première partie : le travail comme malédiction.}

\emph{Idée} : le travail est d'abord peine et contrainte. \emph{Argument 1 (étymologie et tradition)} : le mot vient du latin \emph{tripalium}, instrument de torture ; dans la tradition biblique, le travail est une punition (« tu gagneras ton pain à la sueur de ton front »). \emph{Argument 2 (Marx)} : dans le salariat, le travailleur est aliéné — dépossédé du produit et de son geste. \emph{Exemple} : l'ouvrière de la zone industrielle de Douala qui répète le même geste huit heures par jour pour un salaire de survie. \emph{Mini-conclusion} : dans ces conditions, le travail ne libère pas : il asservit.

\textbf{Conclusion de la rédaction.} Cette première partie prépare le renversement dialectique : c'est précisément parce que le travail peut asservir que la question de sa valeur libératrice doit être posée — ce qui sera l'objet de la seconde partie (Hegel : par le travail, l'homme se forme et se reconnaît dans son œuvre).
\end{exoresolu}

\courssub{Savoir-faire 5 : Rédiger et justifier une conclusion argumentée}

\begin{methode}[Justifier une réponse ou une conclusion]
\begin{enumerate}
\item \textbf{Reprendre la question posée} et y répondre par oui, non, ou (le plus souvent) une réponse nuancée et conditionnelle.
\item \textbf{Rappeler brièvement les deux arguments décisifs} qui fondent la réponse (un par thèse examinée).
\item \textbf{Formuler la nuance} sous la forme : « X, mais à condition que Y » (ex. : le travail libère, à condition d'être créatif, reconnu et justement rétribué).
\item \textbf{Éviter les conclusions vides} : interdiction de terminer par « ainsi le travail est un vaste sujet » ; la conclusion doit trancher.
\item \textbf{Ouvrir} éventuellement sur une question liée (ex. : le loisir a-t-il aussi une valeur ?).
\end{enumerate}
\textbf{Réflexe de rédaction} : une conclusion = réponse + justification + nuance + (ouverture). Quatre phrases suffisent.
\end{methode}

\begin{exoresolu}[Conclure un devoir sur les échanges]
\textbf{Énoncé.} Un élève a montré, dans son devoir, que les échanges marchands organisent la vie économique (marché de Mokolo à Yaoundé) et que les échanges symboliques (dons, salutations, paroles) créent du lien. Rédigez la conclusion justifiée de ce devoir dont la problématique était : « Les échanges sont-ils le fondement de la société ? »

\tcblower
\textbf{Solution détaillée.}

\textbf{Étape 1 : réponse nette à la problématique.} « Nous pouvons maintenant répondre : oui, les échanges sont bien le fondement de la société. »

\textbf{Étape 2 : rappel des arguments décisifs.} « D'une part, l'échange marchand rend possible la division du travail : nul ne produisant tout ce dont il a besoin, le vendeur de vivres de Mokolo et l'enseignant dépendent l'un de l'autre ; sans échange, pas de survie collective. D'autre part, l'échange symbolique — don et contre-don selon Mauss, parole échangée selon le langage — crée la confiance et l'alliance sans lesquelles aucune communauté ne tient. »

\textbf{Étape 3 : nuance conditionnelle.} « Toutefois, l'échange ne fonde la société que s'il reste juste et réciproque : un échange inégal ou trompeur détruit le lien qu'il prétend créer. »

\textbf{Étape 4 : ouverture.} « Reste alors à se demander si la justice de l'échange suffit à fonder la justice de la société tout entière. »

\textbf{Conclusion de l'exercice.} Cette conclusion remplit les quatre critères : elle répond, justifie par deux arguments du cours, nuance et ouvre. Elle serait sanctionnée par la totalité des points prévus au barème pour la conclusion.
\end{exoresolu}

\courssub{Savoir-faire 6 : Commenter un texte philosophique court}

\begin{methode}[Commenter une citation d'auteur]
\begin{enumerate}
\item \textbf{Situer} : nommer l'auteur, l'œuvre et le thème (ex. : Marx, \emph{Manuscrits de 1844}, thème de l'aliénation).
\item \textbf{Dégager la thèse du texte} en une phrase : que soutient l'auteur ?
\item \textbf{Expliquer les termes clés} du texte (travail, produit, étranger, échange...).
\item \textbf{Reconstituer l'argumentation} : quelles raisons l'auteur donne-t-il ? Quels exemples ?
\item \textbf{Discuter} brièvement : accord, objection ou limite, appuyée sur un exemple concret.
\end{enumerate}
\textbf{Réflexe} : ne jamais paraphraser mot à mot ; toujours reformuler la thèse avec ses propres mots avant de la discuter.
\end{methode}

\begin{exoresolu}[Commentaire d'une citation de Marx]
\textbf{Énoncé.} « Le travailleur se sent d'autant plus pauvre qu'il produit plus de richesse ; son travail est pour lui quelque chose d'étranger. » (Marx, \emph{Manuscrits de 1844}.) Expliquez la thèse de ce texte et discutez-la à l'aide d'un exemple.

\tcblower
\textbf{Solution détaillée.}

\textbf{Étape 1 : situation.} Ce texte est extrait des \emph{Manuscrits de 1844} de Karl Marx, où le philosophe analyse l'\cle{aliénation} du travailleur dans la production capitaliste.

\textbf{Étape 2 : thèse.} Marx soutient que, dans le salariat, plus l'ouvrier produit de richesse, plus il s'appauvrit lui-même, parce que le produit de son travail ne lui appartient pas : il lui devient « étranger ».

\textbf{Étape 3 : explication des termes.} « Produit plus de richesse » : l'ouvrier fabrique des biens de grande valeur. « Se sent d'autant plus pauvre » : son salaire ne croît pas avec la valeur produite (la plus-value revient au propriétaire). « Étranger » : le produit se dresse contre lui comme une puissance indépendante — c'est le cœur de l'aliénation.

\textbf{Étape 4 : argumentation reconstituée.} Le raisonnement de Marx : (1) le travailleur vend sa force de travail ; (2) le produit appartient au capitaliste ; (3) donc l'accroissement de la production enrichit le capital, non l'ouvrier ; (4) donc le travail, au lieu de réaliser l'homme, le dépossède.

\textbf{Étape 5 : discussion.} L'exemple de l'ouvrière de Douala confirme la thèse : plus elle coud de chemises exportées, plus l'entreprise s'enrichit, sans que son salaire change. \emph{Limite possible} : dans le travail artisanal ou indépendant (le sculpteur de Foumban), le produit appartient au travailleur, qui se reconnaît dans son œuvre — le travail n'y est plus « étranger ». La thèse de Marx vaut donc pour le travail salarié dans certaines conditions, non pour tout travail.

\textbf{Conclusion.} Le texte de Marx révèle que le problème n'est pas le travail lui-même, mais le rapport de propriété qui sépare le travailleur de son produit : c'est là une analyse toujours pertinente pour comprendre le monde du travail camerounais contemporain.
\end{exoresolu}

\begin{attention}[Les 5 erreurs qui coûtent des points au Bac]
\begin{enumerate}
\item \textbf{Confondre travail et simple activité.} Jouer, se reposer ou agir par instinct ne sont pas du travail : sans intentionnalité, transformation et outillage, la définition ne s'applique pas. Toujours vérifier les trois critères avant de qualifier une situation.
\item \textbf{Traiter le don comme un geste gratuit sans suite.} Erreur classique : oublier que, selon Mauss, le don crée une obligation de recevoir et de rendre. Un échange symbolique sans réciprocité analysée est un échange mal compris.
\item \textbf{Présenter Marx comme un ennemi du travail.} Marx ne condamne pas le travail, il condamne les \emph{conditions} capitalistes qui l'aliènent ; pour lui, le travail est par nature l'activité qui humanise l'homme. Inverser cette thèse est une faute de contenu.
\item \textbf{Rédiger une conclusion qui ne répond pas à la problématique.} « Le travail est un vaste débat » ne vaut aucun point : la conclusion doit trancher (réponse + justification + nuance). Relire la problématique avant de conclure.
\item \textbf{Empiler les noms d'auteurs sans argument ni exemple.} Citer Marx, Mauss et Hegel sans expliquer leurs thèses ni les illustrer (artisan de Foumban, ouvrière de Douala, mariage de Bafoussam) donne un catalogue, pas une dissertation : une idée = un argument + un exemple + une référence.
\end{enumerate}
\end{attention}

\newpage

\coursec{Exercices}

\courssub{Je vérifie mes connaissances}

\begin{exercice}[QCM : définitions et distinctions]
Parmi les propositions suivantes, une seule est exacte. Laquelle ? Justifiez brièvement votre choix.
\begin{enumerate}
    \item Le travail se définit uniquement par l'effort physique fourni.
    \item L'échange marchand suppose toujours l'usage de la monnaie.
    \item Le travail transforme la nature pour satisfaire des besoins humains.
    \item Selon Marx, l'aliénation du travailleur disparaît avec l'augmentation du salaire.
\end{enumerate}
\end{exercice}

\begin{exercice}[Vrai ou Faux justifié]
Indiquez si les affirmations suivantes sont vraies ou fausses, en justifiant votre réponse par une référence précise au cours (auteur, concept, distinction).
\begin{enumerate}
    \item « Travailler, c'est toujours s'aliéner. »
    \item « Le don, chez Marcel Mauss, est un acte purement gratuit et désintéressé. »
    \item « La valeur d'usage d'un bien est identique à sa valeur d'échange. »
    \item « Pour Hannah Arendt, l'\emph{animal laborans} est l'homme qui agit et parle dans la cité. »
\end{enumerate}
\end{exercice}

\begin{exercice}[Définitions philosophiques]
Donnez une définition philosophique précise (genre + différence spécifique) pour chacun des concepts suivants :
\begin{enumerate}
    \item Le travail (au sens anthropologique).
    \item L'aliénation (chez Marx).
    \item L'échange symbolique.
    \item La plus-value.
\end{enumerate}
\end{exercice}

\begin{exercice}[Mise en relation de concepts]
Reliez chaque auteur à la thèse qui lui correspond en justifiant le lien.
\begin{center}
\begin{tabularx}{\linewidth}{|l|X|}\hline
\textbf{Auteur} & \textbf{Thèse} \\ \hline
Karl Marx & A. Le travail libère l'homme en objectivant sa liberté. \\ \hline
Hegel & B. L'échange de dons crée le lien social par l'obligation de rendre. \\ \hline
Marcel Mauss & C. Le travail salarié capitaliste sépare l'ouvrier du produit de son travail. \\ \hline
Hannah Arendt & D. Le travail (labor) assure la survie biologique, l'œuvre (work) bâtit un monde durable. \\ \hline
\end{tabularx}
\end{center}
\end{exercice}

\courssub{Je m'entraîne}

\begin{exercice}[Analyse d'une situation concrète : le secteur informel à Douala]
Un jeune soudeur travaille 12 heures par jour dans un atelier informel à New Bell (Douala). Il fabrique des grilles de fenêtres sur commande. Il est payé à la pièce, n'a pas de contrat, pas de couverture sociale, et utilise des outils qu'il a lui-même achetés.
\begin{enumerate}
    \item Qualifiez cette activité au regard de la distinction \emph{travail / emploi / ouvrage}.
    \item En quoi cette situation illustre-t-elle la thèse de Marx sur l'aliénation du travailleur (rapport au produit, à l'acte de production, à soi, aux autres) ?
    \item L'échange qui a lieu ici (vente de la grille) est-il un échange marchand pur ? Quel autre type d'échange peut-on identifier (relation patron/ouvrier, insertion dans le quartier) ?
\end{enumerate}
\end{exercice}

\begin{exercice}[La tontine : échange économique et lien social]
Dans un quartier de Yaoundé, un groupe de dix femmes commerçantes organise une tontine mensuelle. Chacune verse 10 000 FCFA ; le pot (100 000 FCFA) revient à tour de rôle à l'une d'elles. Elles se retrouvent le soir de la remise pour un repas partagé.
\begin{enumerate}
    \item Identifiez les trois obligations du don chez Mauss (donner, recevoir, rendre) dans ce fonctionnement.
    \item Montrez que l'échange économique (l'argent) est indissociable d'un échange symbolique (confiance, honneur, parole donnée).
    \item En quoi la tontine constitue-t-elle un « fondement de la société » au niveau local (solidarité, mutualisation du risque) ?
\end{enumerate}
\end{exercice}

\begin{exercice}[Transformation d'une question fermée en problématique]
Sujet : « Le travail est-il pénible ? »
\begin{enumerate}
    \item Expliquez pourquoi cette question, posée ainsi, ne constitue pas encore un sujet de dissertation philosophique.
    \item Proposez \textbf{deux} reformulations problématisées (ouvertes, tension entre deux thèses).
    \item Pour chaque reformulation, citez un auteur qui pourrait soutenir la thèse affirmative et un auteur pour la thèse négative.
\end{enumerate}
\end{exercice}

\begin{exercice}[Comparaison de deux conceptions de la valeur du travail]
\begin{enumerate}
    \item Résumez en trois lignes la conception \emph{libérale} de la valeur du travail (source de propriété, d'échange, de richesse).
    \item Résumez en trois lignes la conception \emph{marxiste} de la valeur (travail abstrait, substance de la valeur, exploitation).
    \item Appliquez ces deux grilles de lecture à la situation suivante : « Une start-up camerounaise développe une application agricole et la vend à une multinationale. » Comment chaque courant analyse-t-il la création de valeur ici ?
\end{enumerate}
\end{exercice}

\courssub{Je me prépare au Bac}

\begin{exercice}[Dissertation type Baccalauréat Technique]
\textbf{Sujet :} « Le travail n'est-il qu'une contrainte ? »
\begin{enumerate}
    \item Rédigez une \textbf{introduction complète} (accroche, définition des termes, analyse du sujet, problématique, annonce du plan).
    \item Proposez un \textbf{plan détaillé en trois parties} (avec au moins deux sous-parties chacune), en indiquant pour chaque sous-partie : l'idée directrice, l'argument principal, l'auteur/référence mobilisé, et un exemple concret (idéalement camerounais).
    \item Rédigez la \textbf{transition} entre la première et la deuxième partie.
\end{enumerate}
\end{exercice}

\begin{exercice}[Explication de texte : Marx et l'aliénation]
\textbf{Texte :} « L'ouvrier ne se sent pas chez lui dans son travail ; il se sent chez lui hors de son travail. [...] Son travail n'est donc pas volontaire, mais forcé, travail forcé. Il n'est donc pas la satisfaction d'un besoin, mais seulement un moyen pour satisfaire des besoins en dehors de lui. [...] Le travail extérieur à l'ouvrier, c'est-à-dire le travail dans lequel il ne s'affirme pas, mais se nie, ne se développe pas, mais s'atrophie physiquement et se ruine moralement. » \hfill \emph{Karl Marx, Manuscrits de 1844}
\begin{enumerate}
    \item Dégagez la \textbf{thèse principale} du texte en une phrase.
    \item Identifiez et expliquez \textbf{trois arguments} que Marx utilise pour fonder cette thèse (citez les passages correspondants).
    \item \textbf{Discutez l'actualité} de cette analyse au regard du secteur informel camerounais (moto-taxi, vendeurs à la sauvette, artisans non déclarés). La précarité renforce-t-elle l'aliénation ou peut-elle être vécue comme une forme d'autonomie ?
\end{enumerate}
\end{exercice}

\begin{exercice}[Question de réflexion structurée]
\textbf{Sujet :} « Peut-on vivre en société sans échanger ? »
Construisez une réponse en \textbf{trois étapes dialectiques} (Thèse / Antithèse / Synthèse). Pour chaque étape :
\begin{itemize}
    \item Formulez l'idée force en une phrase.
    \item Développez l'argumentation en 4-5 lignes.
    \item Mobilisez \textbf{un auteur précis} et \textbf{un exemple camerounais concret} (ex. : tontine, marché de quartier, relations inter-ethniques, diaspora, administration publique).
\end{itemize}
La synthèse doit dépasser l'opposition en précisant la nature de l'échange indispensable (économique, symbolique, communicationnel).
\end{exercice}

\newpage

\coursec{Corrigés des exercices}

\coursec{Leçon 13 : Le travail et les échanges — Corrigés des activités et exercices}

\begin{corrige}[Exercice 1 — QCM : définitions et distinctions]
\textbf{Réponse exacte : la proposition 3.}

\textbf{Justification :} Le travail, au sens anthropologique, est l'activité humaine \emph{consciente et finalisée} par laquelle l'homme transforme la nature (matière première, milieu) pour produire des biens et des services destinés à satisfaire ses besoins. C'est la définition retenue notamment par Marx (\emph{Le Capital}, livre I) : le travail est un « procès » entre l'homme et la nature.

\textbf{Pourquoi les autres propositions sont fausses :}
\begin{itemize}
    \item \textbf{Proposition 1 :} fausse. Le travail ne se réduit pas à l'effort physique : il comprend l'intelligence, la technique, la créativité (travail intellectuel, artistique, de service). Un enseignant ou un programmeur travaille sans effort musculaire dominant.
    \item \textbf{Proposition 2 :} fausse. Le \emph{troc} est un échange marchand sans monnaie (échange direct de biens). La monnaie facilite l'échange mais ne le définit pas.
    \item \textbf{Proposition 4 :} fausse. Chez Marx, l'aliénation tient au \emph{mode de production capitaliste} (séparation du travailleur d'avec le produit, les moyens de production, son activité et les autres hommes), non au niveau du salaire. Augmenter le salaire ne supprime pas l'exploitation ni la dépossession.
\end{itemize}

\textbf{Points de vigilance :} au QCM, toujours justifier par une définition ou une référence du cours, jamais par la seule intuition. Attention à l'adverbe « uniquement » ou « toujours », souvent indice d'une proposition fausse.
\end{corrige}

\begin{corrige}[Exercice 2 — Vrai ou Faux justifié]
\begin{enumerate}
    \item \textbf{Faux.} L'aliénation n'est pas une conséquence nécessaire de tout travail : elle caractérise, chez Marx, le travail \emph{salarié dans le mode de production capitaliste}, où l'ouvrier est séparé du produit, de l'acte de production, de son essence et des autres. Le travail peut aussi être \emph{libérateur} : chez Hegel (\emph{Phénoménologie de l'esprit}, dialectique du maître et de l'esclave), le travail forme et émancipe l'esclave en objectivisant sa liberté dans l'objet façonné.
    \item \textbf{Faux.} Dans \emph{Essai sur le don} (1925), Marcel Mauss montre que le don obéit à \textbf{trois obligations} : donner, recevoir, rendre. Le don n'est jamais purement gratuit : il engage celui qui reçoit et crée un lien social durable (l'esprit de la chose donnée, le \emph{hau} maori, oblige à restituer).
    \item \textbf{Faux.} La \emph{valeur d'usage} est l'utilité concrète d'un bien (un pagne habille) ; la \emph{valeur d'échange} est le rapport quantitatif dans lequel il s'échange contre d'autres biens. Un bien peut avoir une grande valeur d'usage et une faible valeur d'échange (l'eau), et inversement (le diamant). C'est la distinction classique reprise par Marx au début du \emph{Capital}.
    \item \textbf{Faux.} Chez Hannah Arendt (\emph{Condition de l'homme moderne}, 1958), l'\emph{animal laborans} est l'homme réduit au \emph{travail} (labor) : la production répétitive nécessaire à la survie biologique. L'homme qui agit et parle dans la cité relève de l'\emph{action} (praxis), troisième activité de la \emph{vita activa}, distincte du travail et de l'œuvre.
\end{enumerate}
\textbf{Points de vigilance :} exiger toujours la référence précise (auteur + concept) ; une réponse « vrai/faux » sans justification ne vaut aucun point.
\end{corrige}

\begin{corrige}[Exercice 3 — Définitions philosophiques]
\begin{enumerate}
    \item \textbf{Le travail (sens anthropologique) :} activité proprement humaine (\emph{genre} : activité consciente et finalisée) par laquelle l'homme transforme la nature et lui-même pour produire des biens répondant à ses besoins (\emph{différence spécifique} : médiation technique entre l'homme et la nature, impliquant un projet préalablement pensé).
    \item \textbf{L'aliénation (chez Marx) :} processus (\emph{genre}) par lequel le travailleur est dépossédé du produit de son travail, de son activité, de son essence humaine et de sa relation aux autres, qui lui deviennent étrangers et hostiles (\emph{différence spécifique} : la chose produite par l'homme se retourne contre lui comme puissance extérieure).
    \item \textbf{L'échange symbolique :} forme d'échange (\emph{genre}) dans laquelle ce qui circule (dons, paroles, signes, honneur) vise moins l'utilité économique que la création et l'entretien du lien social (\emph{différence spécifique} : la valeur des objets échangés est dans la relation qu'ils instituent, non dans leur prix).
    \item \textbf{La plus-value :} différence (\emph{genre}) entre la valeur produite par le travailleur pendant sa journée de travail et la valeur de sa force de travail (le salaire) (\emph{différence spécifique} : travail non payé, approprié par le capitaliste, source du profit et de l'exploitation selon Marx).
\end{enumerate}
\textbf{Points de vigilance :} une bonne définition philosophique suit le schéma \emph{genre prochain + différence spécifique} ; éviter les définitions par l'exemple (« le travail, c'est quand on... »).
\end{corrige}

\begin{corrige}[Exercice 4 — Mise en relation de concepts]
\begin{center}
\begin{tabularx}{\linewidth}{|l|l|X|}\hline
\textbf{Auteur} & \textbf{Thèse} & \textbf{Justification} \\ \hline
Karl Marx & C & Dans les \emph{Manuscrits de 1844} et \emph{Le Capital}, Marx montre que le salariat capitaliste sépare l'ouvrier du produit (qui appartient au patron), des moyens de production et de sa propre activité : c'est l'aliénation. \\ \hline
Hegel & A & Dans la dialectique du maître et de l'esclave (\emph{Phénoménologie de l'esprit}), le travail de l'esclave transforme les choses et le transforme lui-même : il objectivise sa conscience et accède à la liberté. \\ \hline
Marcel Mauss & B & \emph{Essai sur le don} : le cycle donner--recevoir--rendre fonde le lien social ; refuser de rendre, c'est rompre l'alliance et risquer la guerre. \\ \hline
Hannah Arendt & D & \emph{Condition de l'homme moderne} : distinction entre \emph{labor} (travail, survie, consommation), \emph{work} (œuvre, fabrication d'un monde durable) et \emph{action}. \\ \hline
\end{tabularx}
\end{center}
\textbf{Points de vigilance :} ne pas se contenter de relier : la justification par l'œuvre ou le concept est exigée. Piège classique : attribuer à Arendt la thèse hégélienne de la libération par le travail.
\end{corrige}

\begin{corrige}[Exercice 5 — Analyse d'une situation concrète : le secteur informel à Douala]
\begin{enumerate}
    \item \textbf{Qualification :} c'est un \emph{travail} au sens anthropologique (transformation de la matière — le fer — en produit utile — la grille) et un \emph{ouvrage} au sens d'Arendt (fabrication d'un objet durable). Ce n'est pas un \emph{emploi} au sens juridique strict : pas de contrat, pas de salariat formel, pas de protection sociale. On parle de \emph{travail informel} ou d'auto-emploi précaire.
    \item \textbf{Aliénation au sens de Marx :}
    \begin{itemize}
        \item \emph{Rapport au produit :} la grille, dès sa fabrication, appartient au client qui l'a commandée ; le soudeur ne dispose pas du fruit de son travail.
        \item \emph{Rapport à l'acte de production :} le rythme (12 h/jour) est imposé par la nécessité de survivre, non choisi ; le travail est « forcé » (Marx).
        \item \emph{Rapport à soi :} l'absence de couverture sociale et l'usure physique montrent que le travail « s'atrophie physiquement » au lieu d'épanouir.
        \item \emph{Rapport aux autres :} la concurrence entre artisans informels et la dépendance envers les clients dégradent les relations humaines en rapports marchands.
    \end{itemize}
    \textbf{Nuance :} possédant ses propres outils, il n'est pas aliéné au sens strict du prolétaire dépossédé de ses moyens de production ; l'aliénation est ici partielle (précarité plus que salariat).
    \item \textbf{Type d'échange :} la vente de la grille est un échange \emph{marchand} (bien contre monnaie, payé à la pièce). Mais on identifie aussi un échange \emph{symbolique} : la confiance de voisinage, la réputation dans le quartier de New Bell, la parole donnée sur les délais, les relations de fidélité patron/client. L'économie informelle camerounaise repose largement sur ce mélange d'échange marchand et de lien social.
\end{enumerate}
\textbf{Points de vigilance :} appliquer les quatre dimensions de l'aliénation \emph{une par une} ; savoir nuancer (l'artisan possède ses outils) montre la maîtrise du concept plutôt que son application mécanique.
\end{corrige}

\begin{corrige}[Exercice 6 — La tontine : échange économique et lien social]
\begin{enumerate}
    \item \textbf{Les trois obligations de Mauss :}
    \begin{itemize}
        \item \emph{Donner :} chaque femme verse ses 10 000 FCFA mensuels, y compris les mois où elle ne reçoit pas le pot.
        \item \emph{Recevoir :} accepter le pot quand vient son tour, c'est accepter d'être liée au groupe ; refuser serait une rupture.
        \item \emph{Rendre :} celle qui a reçu le pot continue de verser ses cotisations jusqu'à ce que toutes les autres aient bénéficié à leur tour : elle « rend » ce qu'elle a reçu.
    \end{itemize}
    \item \textbf{Échange économique et symbolique indissociables :} économiquement, la tontine est un système d'épargne et de crédit rotatif (chaque membre reçoit 100 000 FCFA, soit la somme des cotisations). Mais sa valeur dépasse le calcul : la confiance mutuelle remplace les garanties bancaires, l'honneur de chacune est engagé (défaillir, c'est perdre la face devant le groupe), et le repas partagé scelle symboliquement la solidarité. La parole donnée vaut contrat.
    \item \textbf{Fondement de la société locale :} la tontine mutualise le risque (financement d'un commerce, d'une urgence, de frais scolaires), crée de la solidarité entre femmes, renforce les liens de voisinage et permet l'inclusion financière de personnes exclues des banques. Elle illustre la thèse de Mauss : l'échange, loin d'être seulement économique, \emph{fabrique du lien social} — c'est en ce sens que les échanges sont le fondement de la société.
\end{enumerate}
\textbf{Points de vigilance :} ne pas réduire la tontine à un « prêt informel » ; l'exigence est de montrer la dimension symbolique (honneur, confiance, réciprocité) conformément à la lecture de Mauss.
\end{corrige}

\begin{corrige}[Exercice 7 — Transformation d'une question fermée en problématique]
\begin{enumerate}
    \item \textbf{Pourquoi ce n'est pas un sujet de dissertation :} « Le travail est-il pénible ? » appelle une réponse fermée (oui/non) et relève du constat empirique ou psychologique (chacun répond selon son expérience). Elle ne contient aucune \emph{tension conceptuelle} : elle n'interroge ni l'essence du travail, ni sa valeur, ni sa nécessité. Un sujet philosophique doit opposer des thèses et exiger une argumentation rationnelle universelle, non un témoignage.
    \item \textbf{Deux reformulations problématisées :}
    \begin{itemize}
        \item R1 : « La peine est-elle l'essence du travail, ou le travail peut-il être source d'épanouissement ? »
        \item R2 : « Le travail aliène-t-il nécessairement l'homme, ou est-il la voie de sa libération ? »
    \end{itemize}
    \item \textbf{Auteurs mobilisables :}
    \begin{itemize}
        \item Pour R1 : thèse affirmative (le travail est peine) — la tradition biblique (« tu gagneras ton pain à la sueur de ton front », Genèse) et Schopenhauer ; thèse négative — Hegel (le travail forme et libère), Paul Lafargue (\emph{Le Droit à la paresse}) ou les éloges de l'activité créatrice (Arendt pour l'œuvre).
        \item Pour R2 : thèse affirmative — Marx (l'aliénation du travail salarié) ; thèse négative — Hegel (dialectique du maître et de l'esclave : par le travail, l'esclave s'émancipe).
    \end{itemize}
\end{enumerate}
\textbf{Points de vigilance :} une problématique doit faire apparaître deux pôles en tension ; vérifier que chaque reformulation reste fidèle au sujet initial (la peine du travail) sans le trahir.
\end{corrige}

\begin{corrige}[Exercice 8 — Comparaison de deux conceptions de la valeur du travail]
\begin{enumerate}
    \item \textbf{Conception libérale (Locke, Smith) :} le travail est la source légitime de la \emph{propriété} (Locke : mélanger son travail à la nature rend la chose sienne) et la mesure de la \emph{valeur d'échange} (Smith : la richesse des nations vient du travail et de sa division). Le travail crée la richesse, et l'échange libre des produits du travail organise harmonieusement la société.
    \item \textbf{Conception marxiste :} la valeur d'une marchandise est déterminée par la quantité de \emph{travail abstrait} socialement nécessaire à sa production. Mais le capitaliste ne paie au travailleur que la valeur de sa force de travail (salaire) : la différence, la \emph{plus-value}, est accaparée comme profit. Le travail est donc à la fois source de toute valeur et siège de l'\emph{exploitation}.
    \item \textbf{Application à la start-up camerounaise :}
    \begin{itemize}
        \item \emph{Lecture libérale :} les développeurs et fondateurs ont créé par leur travail (code, connaissances agronomiques) une valeur nouvelle ; la vente à la multinationale est un échange libre et légitime qui récompense l'innovation, enrichit les créateurs et injecte des capitaux dans l'économie camerounaise.
        \item \emph{Lecture marxiste :} la valeur de l'application vient du travail des salariés (développeurs) ; si le prix de vente dépasse largement les salaires versés, la plus-value a été appropriée par les actionnaires. La cession à une multinationale peut aussi s'analyser comme une extraction de la valeur produite localement vers des capitaux étrangers.
    \end{itemize}
\end{enumerate}
\textbf{Points de vigilance :} respecter la contrainte des trois lignes (concision) ; dans l'application, présenter les deux lectures \emph{symétriquement} sans prendre parti — la discussion critique vient ensuite.
\end{corrige}

\begin{corrige}[Exercice 9 — Dissertation : « Le travail n'est-il qu'une contrainte ? »]
\textbf{Barème indicatif (sur 20) :} introduction 4 pts ; plan détaillé 10 pts (cohérence, auteurs, exemples) ; transition 2 pts ; rigueur philosophique et expression 4 pts.

\begin{enumerate}
    \item \textbf{Introduction complète :}

    \emph{Accroche :} Chaque matin, des millions de Camerounais — du moto-taximan de Douala à la fonctionnaire de Yaoundé — se lèvent pour travailler, souvent sans enthousiasme, par nécessité. \emph{Définition :} le travail est l'activité humaine finalisée de transformation de la nature et de production de biens ; la contrainte est ce qui s'impose à la volonté contre son gré. \emph{Analyse du sujet :} le « n'est-il \emph{que} » invite à tester une réduction : le travail se réduirait-il à une obligation subie (survie, fatigue, discipline) ? \emph{Problématique :} le travail est-il essentiellement une servitude imposée par les besoins, ou peut-il être aussi libération et accomplissement de l'homme ? \emph{Annonce du plan :} nous verrons d'abord que le travail apparaît comme contrainte, puis qu'il peut être libérateur, enfin que sa valeur dépend des conditions sociales dans lesquelles il s'exerce.

    \item \textbf{Plan détaillé :}

    \textbf{I. Le travail comme contrainte.}
    \begin{itemize}
        \item I.1. \emph{Idée :} le travail est d'abord nécessité naturelle : il faut produire pour survivre. \emph{Argument :} l'homme est l'être des besoins ; ne pas travailler, c'est périr. \emph{Référence :} la tradition biblique (Genèse : « tu gagneras ton pain à la sueur de ton front ») ; Arendt sur le \emph{labor} asservi au cycle biologique. \emph{Exemple :} le cultivateur de l'Extrême-Nord soumis aux saisons.
        \item I.2. \emph{Idée :} le travail salarié ajoute une contrainte sociale : l'aliénation. \emph{Argument :} le travailleur ne choisit ni son produit ni son rythme ; le travail est « forcé ». \emph{Référence :} Marx, \emph{Manuscrits de 1844}. \emph{Exemple :} l'ouvrier payé à la pièce dans une usine de Bonabéri.
    \end{itemize}

    \textbf{II. Le travail comme libération.}
    \begin{itemize}
        \item II.1. \emph{Idée :} le travail arrache l'homme à la nature et le forme. \emph{Argument :} en transformant les choses, l'homme se transforme et prend conscience de sa liberté. \emph{Référence :} Hegel, dialectique du maître et de l'esclave. \emph{Exemple :} l'apprenti menuisier de Bafoussam qui devient maître-artisan.
        \item II.2. \emph{Idée :} le travail est source de reconnaissance sociale et de dignité. \emph{Argument :} par son œuvre, l'homme s'insère dans un monde commun et durable. \emph{Référence :} Arendt, la distinction travail/œuvre ; Hegel, la reconnaissance. \emph{Exemple :} l'artiste sculpteur de Foumban reconnu par sa communauté.
    \end{itemize}

    \textbf{III. La valeur du travail dépend de ses conditions.}
    \begin{itemize}
        \item III.1. \emph{Idée :} ce n'est pas le travail en soi qui aliène, mais le mode de production. \emph{Argument :} la même activité peut être servitude ou épanouissement selon l'organisation sociale. \emph{Référence :} Marx, critique du capitalisme, non du travail lui-même. \emph{Exemple :} la coopérative agricole où le producteur maîtrise son produit.
        \item III.2. \emph{Idée :} il faut donc transformer les conditions du travail (droit, protection, sens) pour qu'il devienne libération. \emph{Référence :} le droit du travail et les conventions sociales ; réflexion personnelle argumentée. \emph{Exemple :} l'enjeu de la formalisation du secteur informel au Cameroun.
    \end{itemize}

    \item \textbf{Transition I $\rightarrow$ II :} « Ainsi, le travail se présente d'abord comme une nécessité subie, naturelle puis sociale, au point que l'étymologie latine (\emph{tripalium}, instrument de torture) semble le condamner. Mais cette évidence résiste-t-elle à l'analyse ? Ne confond-on pas la peine du travail avec son essence ? C'est ce qu'il faut examiner en montrant que le travail est aussi ce par quoi l'homme se libère. »
\end{enumerate}
\textbf{Points de vigilance :} l'introduction doit traiter le « \emph{que} » du sujet ; chaque sous-partie exige idée + argument + référence + exemple ; la transition doit faire le bilan de la partie I et annoncer le problème de la partie II, jamais une simple phrase de liaison.
\end{corrige}

\begin{corrige}[Exercice 10 — Explication de texte : Marx et l'aliénation]
\textbf{Barème indicatif (sur 20) :} thèse 4 pts ; trois arguments 9 pts (3 pts chacun) ; discussion 5 pts ; expression 2 pts.

\begin{enumerate}
    \item \textbf{Thèse principale :} dans le travail salarié capitaliste, le travail est extérieur et imposé à l'ouvrier : au lieu de s'y réaliser, il s'y nie et s'y détruit — le travail est aliénation, non accomplissement de l'homme.

    \item \textbf{Trois arguments :}
    \begin{itemize}
        \item \emph{Argument 1 — le travail comme dépossession de soi :} « L'ouvrier ne se sent pas chez lui dans son travail ; il se sent chez lui hors de son travail. » Le travail n'est pas le lieu de l'épanouissement mais un temps perdu pour la vie : l'ouvrier n'existe vraiment qu'en dehors de l'atelier. Le travail lui est \emph{étranger}.
        \item \emph{Argument 2 — le caractère forcé du travail :} « Son travail n'est donc pas volontaire, mais forcé, travail forcé. » L'ouvrier ne travaille pas par désir ou par choix, mais sous la contrainte du besoin (il faut un salaire pour vivre). Ce qui devrait être libre activité devient servitude.
        \item \emph{Argument 3 — le travail réduit à un simple moyen :} « Il n'est donc pas la satisfaction d'un besoin, mais seulement un moyen pour satisfaire des besoins en dehors de lui. » Le travail n'a pas de fin en lui-même : il n'est qu'un instrument pour gagner de quoi manger. D'où la conséquence finale : l'ouvrier « ne s'affirme pas, mais se nie [...] s'atrophie physiquement et se ruine moralement » — le travail détruit au lieu de développer.
    \end{itemize}

    \item \textbf{Discussion — actualité au Cameroun :} le secteur informel semble d'abord confirmer Marx : le moto-taximan de Douala ou le vendeur à la sauvette de Yaoundé travaillent de longues heures, sans protection, usant leur corps, pour un revenu de pure survie — le travail y est bien « moyen pour satisfaire des besoins en dehors de lui » et la précarité aggrave l'aliénation (insécurité, absence de reconnaissance). \emph{Cependant}, une nuance s'impose : l'artisan ou le commerçant informel possède souvent ses outils, choisit ses horaires et garde le produit de son travail ; il peut vivre cette activité comme une \emph{autonomie} face au salariat subordonné, voire comme une fierté (le « débrouillard »). L'aliénation ne tient donc pas seulement à la précarité, mais au rapport au produit et à l'activité : l'informel échappe partiellement à l'aliénation salariale tout en subissant une aliénation par la nécessité économique. La formalisation et la protection sociale apparaissent ainsi comme des conditions de désaliénation.
\end{enumerate}
\textbf{Points de vigilance :} citer littéralement le texte pour chaque argument ; dans la discussion, éviter le manichéisme (l'informel n'est ni pure aliénation ni pure liberté) ; mobiliser la distinction marxiste travail/moyen de subsistance.
\end{corrige}

\begin{corrige}[Exercice 11 — Question de réflexion structurée : « Peut-on vivre en société sans échanger ? »]
\textbf{Barème indicatif (sur 20) :} thèse 5 pts ; antithèse 5 pts ; synthèse 7 pts (dépassement effectif) ; auteurs et exemples 3 pts.

\textbf{Thèse — Non, la société est impossible sans échange.}
\emph{Idée force :} l'échange est le ciment de tout lien social. \emph{Développement :} aucun homme ne se suffit à lui-même : la division du travail rend chacun dépendant des autres pour se nourrir, se vêtir, se soigner. Plus profondément, Marcel Mauss (\emph{Essai sur le don}) montre que l'échange des dons (donner--recevoir--rendre) précède et fonde le contrat social : c'est en échangeant que les groupes passent de l'hostilité à l'alliance. \emph{Exemple :} la tontine des commerçantes de Yaoundé, où la circulation de l'argent crée confiance, solidarité et réciprocité durables.

\textbf{Antithèse — Si, des formes de vie sociale semblent exister sans échange.}
\emph{Idée force :} certains liens sociaux reposent sur le partage ou la gratuité, non sur l'échange. \emph{Développement :} dans la famille, la mère nourrit l'enfant sans contrepartie ; l'entraide communautaire villageoise (travail collectif des champs, « njangi » d'urgence) semble obéir au don pur. On pourrait donc penser qu'une société fondée sur la solidarité gratuite est concevable, et que l'échange marchand corrompt plutôt le lien social (critique de la marchandisation). \emph{Exemple :} les travaux champêtres collectifs dans les villages du Nord, où chacun aide sans paiement.

\textbf{Synthèse — L'échange indispensable n'est pas seulement économique : il est symbolique et communicationnel.}
\emph{Idée force :} on ne peut vivre en société sans échanger, mais l'échange fondamental est celui des signes, des paroles et des dons, non seulement des marchandises. \emph{Développement :} l'antithèse se dissout si l'on élargit le concept d'échange : le don maternel appelle la reconnaissance filiale ; le travail collectif villageois fonctionne par rotation (j'aide aujourd'hui, on m'aidera demain) — c'est bien le cycle de Mauss. Même le langage est un échange : parler, c'est donner et recevoir des signes. Ce qui fonde la société, c'est donc la \emph{réciprocité} sous toutes ses formes : économique (marché de quartier de Mokolo), symbolique (don et contre-don), communicationnelle (dialogue, salutations, palabre). Vivre en société sans aucun échange serait vivre sans relation — donc hors de l'humanité même. \emph{Auteur :} Mauss, complété par Arendt (l'action et la parole comme espace commun). \emph{Exemple :} la palabre villageoise, échange de paroles qui règle les conflits et refonde l'unité du groupe.

\textbf{Points de vigilance :} la synthèse ne doit pas être un compromis mou (« un peu des deux ») mais un \emph{dépassement} : ici, l'élargissement du concept d'échange (économique $\rightarrow$ symbolique et communicationnel) résout l'opposition. Chaque étape doit mobiliser un auteur nommé et un exemple camerounais précis.
\end{corrige}

\newpage

\coursec{Fiche bilan}

\begin{popfigure}
\begin{tikzpicture}[mindmap, grow cyclic, every node/.style={concept, font=\small}, concept color=popblue!60,
level 1/.append style={sibling angle=60, level distance=3.4cm},
level 2/.append style={sibling angle=45, level distance=2.2cm, font=\footnotesize}]
\node[concept color=popink!70, font=\bfseries\small]{Travail et échanges}
  child[concept color=popblue!50]{ node {Définir le travail}
    child { node {Activité humaine transformant la nature} }
    child { node {Travail $\neq$ labeur animal} } }
  child[concept color=popgreen!50]{ node {Valeurs du travail}
    child { node {Valeur morale : dignité} }
    child { node {Valeur sociale : intégration} }
    child { node {Valeur économique : richesse} } }
  child[concept color=poporange!60]{ node {Malédiction ?}
    child { node {Travail-punition (mythes)} }
    child { node {Fatigue, peine} } }
  child[concept color=popgold!60]{ node {Libération ?}
    child { node {Hegel : se humaniser par le travail} }
    child { node {Maîtrise de la nature} } }
  child[concept color=poppink!60]{ node {Aliénation ?}
    child { node {Marx : ouvrier dépossédé} }
    child { node {Travail à la chaîne} } }
  child[concept color=poppurple!50]{ node {L'échange}
    child { node {Donner-recevoir-rendre} }
    child { node {Biens, paroles, services} } }
  child[concept color=popgreen!60]{ node {Fondement de la société}
    child { node {Lien social et réciprocité} }
    child { node {Division du travail} } };
\end{tikzpicture}
\legende{Carte mentale de la leçon : le travail (définition, valeurs, trois regards : malédiction, libération, aliénation) et les échanges comme fondement du lien social.}
\end{popfigure}

\begin{bilan}
\textbf{Définitions clés}
\begin{itemize}
\item \cle{Travail} : activité humaine, consciente et finalisée, par laquelle l'homme transforme la nature et produit des biens utiles ; il se distingue de l'activité animale par la conscience et le projet.
\item \cle{Valeur du travail} : ce qui fait du travail un bien : valeur morale (dignité), sociale (reconnaissance, intégration) et économique (production de richesses).
\item \cle{Aliénation} (Marx) : situation où le travailleur est dépossédé du produit de son travail, de son activité et de son humanité ; le travail devient étranger et subi.
\item \cle{Échange} : transfert réciproque de biens, de services ou de paroles entre des personnes ; il repose sur la \cle{réciprocité} (donner, recevoir, rendre).
\item \cle{Division du travail} : répartition des tâches entre les membres de la société, qui rend chacun dépendant des autres et renforce la solidarité.
\end{itemize}

\textbf{Auteurs et thèses à connaître par cœur}
\begin{center}
\begin{tabularx}{\linewidth}{|l|X|}\hline
\rowcolor{popblueL} \textbf{Auteur / source} & \textbf{Thèse} \\ \hline
Récits bibliques et mythes & Le travail comme punition : « tu gagneras ton pain à la sueur de ton front » (malédiction). \\ \hline
Hegel & Le travail humanise : en transformant la nature, l'homme se transforme et se reconnaît dans son œuvre (libération). \\ \hline
Marx & Dans le capitalisme, le travail est aliénation : l'ouvrier ne possède ni le produit ni les moyens de production. \\ \hline
Durkheim & La division du travail crée la solidarité sociale : nous dépendons les uns des autres. \\ \hline
Mauss & L'échange obéit à la triple obligation : donner, recevoir, rendre ; il fonde le lien social. \\ \hline
\end{tabularx}
\end{center}

\textbf{Méthodes express}
\begin{itemize}
\item Dissertation : définir les termes dès l'introduction, problématiser (le travail est-il une contrainte ou une libération ?), puis construire thèse / antithèse / synthèse.
\item Toujours illustrer par un exemple concret camerounais : le planteur de cacao, l'artisan de Foumban, le commerçant du marché, le tontine (échange de solidarité).
\item Conclure en répondant nettement à la question posée et en ouvrant sur une perspective.
\end{itemize}
\end{bilan}

\begin{aretenir}[Checklist avant le Bac]
\begin{itemize}
\item[$\square$] Je sais définir le travail et le distinguer de l'activité animale.
\item[$\square$] Je connais les trois valeurs du travail (morale, sociale, économique).
\item[$\square$] Je sais expliquer la thèse du travail-malédiction avec un exemple.
\item[$\square$] Je sais expliquer la thèse de Hegel : le travail comme libération.
\item[$\square$] Je sais définir l'aliénation selon Marx et citer un exemple (travail à la chaîne).
\item[$\square$] Je sais définir l'échange et la réciprocité (donner, recevoir, rendre).
\item[$\square$] Je sais montrer que les échanges fondent la société (Durkheim, Mauss).
\item[$\square$] Je sais citer au moins deux exemples camerounais (tontine, marché, coopérative).
\item[$\square$] Je sais construire une problématique à partir d'un sujet sur le travail.
\item[$\square$] Je sais rédiger une introduction : définition, problématique, annonce du plan.
\item[$\square$] Je sais justifier ma conclusion par des arguments et des références précises.
\end{itemize}
\end{aretenir}

\findecours

\end{document}
